\documentclass{article}
\usepackage{graphicx} % Required for inserting images
\usepackage{amsfonts, amsmath, amssymb, amsthm, float}
\usepackage{tikz-cd}

\title{Complementi di topologia, geometria e algebra}
\author{mathcal notes}
\date{September 2025}

\begin{document}

\maketitle
\tableofcontents
\newpage
\section{Topologia}
\textbf{Teorema di De Morgan:} siano $X$ un insieme e $\{ A_i \}_{i \in I} \subset X$ una famiglia di sottoinsiemi. Allora:
\begin{gather}
    X \setminus \bigcup_{i \in I} A_i = \bigcap_{i \in I} (X \setminus A_i) \notag \\
    X \setminus \bigcap_{i \in I} A_i = \bigcup_{i \in I} (X \setminus A_i) \notag
\end{gather} \\

\textbf{Dimostrazione:} se $x \in X \setminus \bigcup_{i \in I} A_i$ si ha $x \not\in \bigcup_{i \in I} A_i$, pertanto $x \not\in A_i$ per ogni $i \in I$, ciò equivale a $x \in X \setminus A_i$ per ogni $i \in I$, dunque $x \in \bigcap_{i \in I} (X \setminus A_i)$. Viceversa, se $y \in \bigcap_{i \in I} (X \setminus A_i)$ segue che $y \in X \setminus A_i$ per ogni $i \in I$, pertanto $y \not\in A_i$ per ogni $i \in I$, dunque $y \in X \setminus \bigcup_{i \in I} A_i$. \\ \\
Se $x \in X \setminus \bigcap_{i \in I} A_i$ si ha che $x \not\in \bigcap_{i \in I} A_i$, pertanto esiste $j \in I$ tale che $x \not\in A_j$, ciò comporta che $x \in X \setminus A_j$, quindi $x \in \bigcup_{i \in I} (X \setminus A_i)$. Viceversa, se $y \in \bigcap_{i \in I} (X \setminus A_i)$ segue che esiste $k \in I$ tale che $y \in X \setminus A_k$, ovvero $y \not\in A_k$, pertanto $y \not\in \bigcap_{i \in I} A_i$, dunque $y \in X \setminus \bigcap_{i \in I} A_i$. \qed \\ \\ \\
\textbf{W.O.P (assioma):} sia $S \subseteq \mathbf{N}$, con $\mathbf{N} := \{0, 1, 2, \dots\}$, allora esiste $m \in S$ tale che $m \leq x$ per ogni $x \in S$.
\subsection{Spazi topologici e applicazioni continue}
\textbf{Definizione 1.1.1:} sia $X$ un insieme. Una \textit{topologia} su $X$ è un sottoinsieme $\tau \subseteq P(X)$ tale che: \\
\textbf{1.} $\emptyset, X \in \tau$. \\
\textbf{2.} se $A, B \in \tau$ allora $A \cap B \in \tau$. \\
\textbf{3.} se $\{ A_i \}_{i \in I} \subset \tau$ allora $\bigcup_{i \in I} A_i \in \tau$. \\ \\
Si dice che la coppia $(X, \tau)$ è uno \textit{spazio topologico} e gli elementi di $\tau$ si dicono \textit{aperti}. \\ \\
\textbf{Esempi:} - sia $X$ un insieme, $\tau = P(X)$ è una topologia su $X$, detta \textit{topologia discreta}. \\
- sia $X$ un insieme, $\tau = \{ \emptyset, X \}$ è una topologia su $X$, detta \textit{topologia banale}. \\
- sia $X = \mathbf{R}$, sia $A \in \tau$ se e solo se $A$ è della forma $A = \bigcup_{k \in \mathbf{N}} (a_k, b_k)$. Allora $\tau$ è una topologia, detta \textit{topologia euclidea}. Sia $B = \bigcup_{j \in \mathbf{N}} (c_j, d_j)$, si ha che $A \cap B = \bigcup_{k \in \mathbf{N}} (a_k, b_k) \cap \bigcup_{j \in \mathbf{N}} (c_j, d_j) = \bigcup_k \bigcup_j ((a_k, b_k) \cap (c_j, d_j))$, pertanto anche $A \cap B \in \tau$. \\ \\ \\
\textbf{Definizione 1.1.2:} sia $(X, \tau)$ uno spazio topologico, sia $C \subseteq X$, $C$ è detto \textit{chiuso} se $X \setminus C \in \tau$. \\ \\
Notiamo che è possibile definire una topologia su un insieme partendo da una definizione di insieme chiuso. Denotiamo con $\mathcal{F}$ la famiglia dei sottoinsiemi chiusi di $X$, dalla definizione di insieme aperto si ottiene che $\mathcal{F}$ gode delle seguenti proprietà: \\
\textbf{1.}$\emptyset, X \in \mathcal{F}$. \\
\textbf{2.} L'intersezione di una qualsiasi famiglia di elementi di $\mathcal{F}$ è un elemento di $\mathcal{F}$. \\
\textbf{3.} L'unione di un numeri finito di elementi di $\mathcal{F}$ è una elemento di $\mathcal{F}$. \\ \\
Pertanto dato un insieme non vuoto $X$ ed una famiglia $\mathcal{F}$ che gode delle precedenti proprietà, si ha che la famiglia $\tau \subset \mathcal{P}(X)$ definita da
\begin{gather}
    A \in \tau \iff X \setminus A \in \mathcal{F} \notag
\end{gather}
è una topologia su $X$. \\ \\
\textbf{Osservazioni:} se $C, D$ sono chiusi allora anche $C \cap D$ è chiuso, infatti
\begin{gather}
    X \setminus (C \cap D) = X \setminus C \cup X \setminus D \notag
\end{gather}
unione di aperti, pertanto $X \setminus (C \cap D) \in \tau$. \\ \\
Se $\{ C_i \}_{i \in I}$ è una famiglia di chiusi, allora $\bigcap_{i \in I} C_i$ è chiuso, infatti
\begin{gather}
    X \setminus \bigcap_{i \in I} C_i = \bigcup_{i \in I} (X \setminus C_i) \notag
\end{gather}
unione numerabile di aperti, pertanto $\bigcup_{i \in I} (X \setminus C_i) \in \tau$. \\ \\ \\
\textbf{Esempi:} - sia $X = \mathbf{R}$, allora $\bigcap_{n \in \mathbf{N}} \left(-\frac{1}{n}, \frac{1}{n}\right) = \emptyset$, ovvero in generale un'intersezione numerabile di aperti non è un aperto. \\
- sia $X = \mathbf{R}$, allora $\bigcup_{-1 < a < b < 1} [a, b] = (-1, 1)$, pertanto in generale l'unione (in questo caso non numerabile) di chiusi non è un chiuso. \\ \\ \\
\textbf{Definizione 1.1.3:} sia $(X, \tau)$ uno spazio topologico, sia $Y \subseteq X$. Allora $\tau_Y := \{ A \cap Y : A \in \tau \}$ è una topologia su $Y$, detta \textit{topologia indotta} su $Y$ da $X$, in tal caso si dice che $Y$ è un \textit{sottospazio} di $X$. Notiamo che gli aperti di $(Y, \tau_Y)$ NON sono uguali agli aperti di $(X, \tau)$. \\ \\ \\
\textbf{Definizione 1.1.4:} siano $(X, \tau)$ uno spazio topologico, $x_0 \in X$ e $A \subseteq X$. $A$ è detto \textit{intorno} di $x_0$ se esiste $W \in \tau$ tale che $W \subseteq A$ e $x_0 \in W$. Poniamo $\mathcal{N}(x_0) := \{ V \in \tau : V \text{ intorno di } x_0 \}$. \\ \\
Un \textit{sistema fondamentale di intorni di un punto} $x$ dello spazio topologico $X$ è una famiglia $\mathcal{B}(x)$ di intorni di $x$ tale che per ogni intorno $M$ di $x$ esiste $N \in \mathcal{B}(x)$ tale che $N \subseteq M$. \\
L'insieme $\mathcal{N}(x)$ di tutti gli intorni di $x$ è un particolare esempio di sistema fondamentale di intorni di $x$. \\ \\ \\
\textbf{Proposizione 1.1.1:} sia $(X, \tau)$ uno spazio topologico, sia $A \subseteq X$. $A$ è aperto se e solo se per ogni $x \in A$ esiste $V \in \mathcal{N}(x)$ tale che $V \subseteq A$. \\

\textbf{Dimostrazione:} ($\rightarrow$) sia $A$ aperto, poniamo $V := A$, si ha che $V \in \tau$ e $V \subseteq A$, infine $x \in V$ per ogni $x \in A$, dunque $A$ è un intorno di $x$ per ogni $x \in A$. \\
($\leftarrow$) supponiamo ora che per ogni $x \in A$ esiste $V_x \in \mathcal{N}(x)$ tale che $V_x \subseteq A$, allora per definizione di intorno esiste $W_x \subseteq V_x$ tale che $W_x \in \tau$ per ogni $x \in A$, dunque $A = \bigcup_{x \in A} W_x$, unione di aperti, pertanto $A \in \tau$. \qed \\ \\ \\
\textbf{Definizione 1.1.5:} siano $(X, \tau)$ uno spazio topologico e $\mathcal{B} \subseteq \tau$ una famiglia di aperti. $\mathcal{B}$ è una \textit{base} per $\tau$ se ogni aperto $\tau$ di è unione di aperti di $\mathcal{B}$. \\ \\
Dalla definizione segue che se $\mathcal{B}$ è una base per $\tau$ e $\mathcal{B} \subset \mathcal{C} \subset \tau$ si ha che anche $\mathcal{C}$ è una base per $\tau$. \\ \\
Una \textit{sottobase} è una famiglia $\mathcal{S}$ di insiemi aperti tali che la famiglia
\begin{gather}
    \mathcal{B}_\mathcal{S} := \{ U_1 \cap \dots \cap U_n : U_1, \dots , U_n \in \mathcal{S} \} \notag
\end{gather}
costituita dalle intersezioni finite di elementi di $\mathcal{S}$ sia una base della topologia $\tau$. \\ \\ \\
\textbf{Proposizione 1.1.2:} sia $X$ un insieme non vuoto e $\mathcal{B}$ una famiglia di sottoinsiemi di $X$ tali che: \\
\textbf{1.} $\bigcup_{B \in \mathcal{B}} B = X$. \\
\textbf{2.} dati comunque $A, B \in \mathcal{B}$, si ha che $A \cap B$ è unione di elementi di $\mathcal{B}$. \\ \\
Allora esiste una topologia $\tau_\mathcal{B}$ su $X$ di cui $\mathcal{B}$ è una base; $\tau_\mathcal{B}$ è l'unica topologia con tale proprietà. \\

\textbf{Dimostrazione:} se $\tau_\mathcal{B}$ esiste è unica poiché ogni suo aperto deve essere unione di elementi di $\mathcal{B}$ e, viceversa, ogni unione di elementi di $\mathcal{B}$ deve appartenere a $\tau_\mathcal{B}$, perché gli elementi di $\mathcal{B}$ sono aperti in $\tau_\mathcal{B}$. Dunque $\tau_\mathcal{B}$ deve coincidere con la famiglia di sottoinsiemi di $X$ che sono unione di elementi di $\mathcal{B}$. \\ \\
Dimostriamo ora che $\tau_\mathcal{B}$ è una topologia: innanzitutto $\emptyset, X \in \tau_\mathcal{B}$, il primo come unione della famiglia vuota di elementi di $\mathcal{B}$, il secondo per l'ipotesi \textbf{1}. Consideriamo ora una famiglia $\{ U_j \}_{j \in J}$ di elementi di $\tau_\mathcal{B}$, allora si ha
\begin{gather}
    U_j = \bigcup_{k \in K(j)} B_k, \quad B_k \in \mathcal{B} \quad \forall j \in J \notag
\end{gather}
dove $K(j)$ sono opportuni insiemi di indici. Quindi si ottiene
\begin{gather}
    \bigcup_{j \in J} U_j = \bigcup_{j \in J} \left( \bigcup_{k \in K(j)} B_k \right) \in \tau_\mathcal{B} \notag
\end{gather}
Infine se
\begin{gather}
    U_1 = \bigcup_{h \in H} B_h \in \tau_\mathcal{B}, \quad U_2 = \bigcup_{k \in K} B_k \in \tau_\mathcal{B} \notag
\end{gather}
si ottiene
\begin{align}
    U_1 \cap U_2 &= \bigcup_{h \in H} B_h \cap \bigcup_{k \in K} B_k = \notag \\
    &= \bigcup_{\substack{h \in H \\ k \in K}} (B_h \cap B_k) \notag
\end{align}
poiché per ipotesi $B_h \cap B_k \in \tau_\mathcal{B}$ per ogni $(h, k) \in H \times K$, si ottiene la tesi. \qed \\ \\ \\
Uno spazio topologico $X$ soddisfa il \textit{primo assioma di numerabilità} se ogni suo punto possiede un sistema fondamentale di intorni numerabile, cioè se per ogni $x \in X$ esiste una famiglia
\begin{gather}
    \{ N_i : i = 1, 2, \dots \} \notag
\end{gather}
che è una base di intorni di $x$. Si può sempre supporre che un sistema fondamentale di intorni numerabile soddisfi l'ulteriore condizione $N_{i + 1} \subset N_i$ per ogni $i$, basta sostituire $\{ N_i \}$ con
\begin{gather}
    \{ M_i := \bigcap_{j \leq i} N_i : i = 1, 2, \dots \} \notag
\end{gather}
Uno spazio topologico $X$ soddisfa il \textit{seconda assioma di numerabilità} se $X$ possiede una base numerabile. \\
Ogni spazio metrizzabile soddisfa il primo assioma di numerabilità, infatti per ogni $x \in X$ la famiglia
\begin{gather}
    \mathcal{B}(x) := \{ D_{1/n}(x) : n \in \mathbf{N} \} \notag
\end{gather}
è un sistema fondamentale di intorni numerabile. \\
In particolare $\mathbf{R}^n$ soddisfa il secondo assioma di numerabilità perché la famiglia
\begin{gather}
    \mathcal{B} := \{ D_{1/n}(\mathbf{x}) := n \in \mathbf{N}, \mathbf{x} \in \mathbf{Q}^n \} \notag
\end{gather}
è una base numerabile della topologia euclidea. \\ \\ \\
\textbf{Definizione 1.1.6:} siano $X, Y$ spazi topologici, $x_0 \in X$, $f : X \longrightarrow Y$, allora $f$ è detta \textit{continua} in $x_0$ se per ogni $V \in \mathcal{N}(f(x_0))$ esiste $W \in \mathcal{N}(x_0)$ tale che $f(W) \subset V$. \\
$f$ si dice \textit{continua} se è continua per ogni $x \in X$. \\ \\ \\
\textbf{Proposizione 1.1.3:} siano $X, Y$ spazi topologici, $f : X \longrightarrow Y$, allora $f$ è continua se e solo se $f^{-1}(A) \subseteq X$ è un aperto per ogni aperto $A \subseteq Y$. \\

\textbf{Dimostrazione:} ($\rightarrow$) siano $f$ continua per ipotesi, $A \subseteq Y$ un aperto, \\ $x \in f^{-1}(A)$ e infine $y := f(x) \in A$. Poiché $A \in \mathcal{N}(f(x))$, per ipotesi di continuità, si ha che esiste $W \in \mathcal{N}(x)$ tale che $f(W) \subset A$, ciò comporta $W \subset f^{-1}(A)$, pertanto, dato che vale per ogni $x \in f^{-1}(A)$, per la \textit{proposizione 1.1.1} si ottiene che $f^{-1}(A)$ è un aperto di $X$. \\ \\
$(\leftarrow)$ Supponiamo per ipotesi che per ogni aperto $A \subseteq Y$ si ha che $f^{-1}(A)$ è un aperto di $X$. Sia $x \in X$, consideriamo un intorno $V \in \mathcal{N}(f(x))$, per definizione esiste un aperto $W$ di $Y$ tale che $f(x) \in W \subseteq V$, inoltre per ipotesi si ha che $f^{-1}(W)$ è un aperto di $X$. Pertanto otteniamo $f(f^{-1}(W)) \subseteq W \subseteq V$, quindi $f$ è continua. \qed \\ \\ \\
\textbf{Osservazione:} $f$ è continua se e solo se $f^{-1}(C)$ è un chiuso per ogni chiuso $C \subseteq Y$. \\ \\
$f$ continua NON implica in generale che $f(A)$ è un aperto per ogni aperto $A \subseteq X$; infatti sia $X := \mathbf{R}$, topologia euclidea, e sia $f(x) = x^2$ per ogni $x \in \mathbf{R}$, allora $f$ è continua, consideriamo l'intervallo aperto $(-1, 1)$, si ha che $f((-1, 1)) = [0, 1)$ che non è aperto né chiuso. \\ \\ \\
\textbf{Definizione 1.1.7:} sia $f : X \longrightarrow Y$ un'applicazione tra spazi topologici, $f$ è detta \textit{omeomorfismo} se: \\
\textbf{1.} $f$ è continua. \\
\textbf{2.} $f$ è biunivoca (biettiva e suriettiva). \\
\textbf{3.} $f^{-1}$ è continua. \\ \\
Se tra due spazi topologici $X, Y$ esiste un omeomorfismo si dice che $X$ è omeomorfo a $Y$ (o viceversa) e si scrive $X \cong Y$. \\ \\ \\
\textbf{Definizione 1.1.8:} sia $\mathcal{P}$ una proprietà, $\mathcal{P}$ è detta \textit{proprietà topologica} se per $X \cong Y$ si ha che $X$ ha $\mathcal{P}$ se e solo se $Y$ ha $\mathcal{P}$, ovvero se la proprietà è invariante per omeomorfismi. \\ \\ \\
\textbf{Definizione 1.1.9:} si dice che $f : X \longrightarrow X$ è un \textit{automorfismo} se $f$ è un omeomorfismo. \\ \\ \\
\textbf{Definizione 1.1.10:} sia $X$ uno spazio topologico, sia $A \subseteq X$, è detta \textit{parte interna di $A$} l'insieme
\begin{gather}
    \mathring{A} := \bigcup_{\substack{B \subseteq A \\ B \text{ aperto}}} B \notag
\end{gather}
ovvero il più "grande" aperto contenuto in $A$. \\ \\
Mentre è detta \textit{chiusura di $A$} l'insieme
\begin{gather}
    \overline{A} := \bigcap_{\substack{C \supseteq A \\ C \text{ chiuso}}} C \notag
\end{gather}
ovvero il più piccolo chiuso che contiene $A$. \\ \\ \\
\textbf{Osservazione:} ovviamente vale che $\mathring{A} \subseteq A \subseteq \overline{A}$, inoltre se $A$ è aperto si ha $\mathring{A} = A$, mentre se $A$ è chiuso si ha $\overline{A} = A$. \\ \\ \\
\textbf{Definizione 1.1.11:} siano $X$ uno spazio topologico, $A \subseteq X$ e $x_0 \in A$. $x_0$ è detto \textit{punto di accumulazione} per $A$ se per ogni $V \in \mathcal{N}(x_0)$ si ha che
\begin{gather}
    A \cap (V \setminus \{ x_0 \}) \neq \emptyset \notag
\end{gather}
Il \textit{derivato di $A$}, $D(A)$, è l'insieme di tutti i punti di accumulazione di $A$. \\ \\ \\
\textbf{Proposizione 1.1.4:} siano $X$ uno spazio topologico, $A \subseteq X$, allora $D(A) \subseteq \overline{A}$. \\

\textbf{Dimostrazione:} sia $x \in D(A)$, supponiamo per assurdo che $x \in X \setminus \overline{A}$, poiché $\overline{A}$ è chiuso si ha che $X \setminus \overline{A}$ è aperto, pertanto esiste $V \in \mathcal{N}(x)$ tale che $V \subset X \setminus \overline{A}$, ciò comporta che $\overline{A} \cap (V \setminus \{x_0 \} = \emptyset$ e poiché $\overline{A} \supseteq A$ si ottiene $A \cap (V \setminus \{x_0 \} = \emptyset$, ma è un assurdo. \qed \\ \\ \\
\textbf{Definizione 1.1.12:} siano $X$ uno spazio topologico, $A \subseteq X$, la \textit{frontiera} (o \textit{bordo}) di $A$ è $\partial A := \overline{X \setminus A} \cap \overline{A}$. \\ \\ 
\textbf{Osservazione:} vale sempre $\overline{X \setminus A} = X \setminus \mathring{A}$, infatti si ha 
\begin{gather}
    X \setminus \mathring{A} = X \setminus \bigcup_{\substack{S \subseteq A \\ S \text{ aperto}}} S = \bigcap_{\substack{S \subseteq A \\ S \text{ aperto}}} (X \setminus S) = \bigcap_{\substack{R \supseteq X \setminus A \\ R \text{ chiuso}}} R = \overline{X \setminus A} \notag
\end{gather}
ponendo $R := X \setminus S$, infatti, poiché $S$ è un aperto contenuto in $A$, si ha che $X \setminus S$ è un chiuso che contiene il complementare di $A$ e dunque otteniamo per definizione la chiusura di $X \setminus A$. \\ \\
Analogamente vale sempre $\mathring{X \setminus A} = X \setminus \overline{A}$, infatti si ha
\begin{gather}
    X \setminus \overline{A} = X \setminus \bigcap_{\substack{R \supseteq A \\ R \text{ chiuso}}} R = \bigcup_{\substack{R \supseteq A \\ R \text{ chiuso}}} (X \setminus R) = \bigcap_{\substack{S \subseteq X \setminus A \\ S \text{ aperto}}} S = \mathring{X \setminus A} \notag
\end{gather}
avendo posto in questo caso $S := X \setminus R$, un aperto contenuto in $X \setminus A$. \\ \\ \\
\textbf{Definizione 1.1.13:} siano $X_1, X_2$ spazi topologici, un rettangolo aperto in $X := X_1 \times X_2$ è un insieme della forma $A_1 \times A_2$ dove $A_1$ è aperto in $X_1$ e $A_2$ è aperto in $A_2$. \\
La \textit{topologia prodotto} è la topologia dove gli aperti sono unione di rettangoli aperti. \\ \\
Consideriamo infatti la famiglia dei sottoinsiemi di $X$ che sono dati dal prodotto cartesiano di un aperto di $X_1$ con un aperto di $X_2$, cioè
\begin{gather}
    \mathcal{B} := \{ A_1 \times A_2 : A_1 \in \tau_1, A_2 \in \tau_2 \} \notag
\end{gather}
la famiglia $\mathcal{B}$ è un ricoprimento di $X$ perché $X = X_1 \times X_2$ è un elemento di $\mathcal{B}$. Inoltre se $A_1 \times A_2, B_1 \times B_2 \in \mathcal{B}$, allora
\begin{gather}
    (A_1 \times A_2) \cap (B_1 \times B_2) = (A_1 \cap B_1) \times (A_2 \times B_2) \in \mathcal{B} \notag
\end{gather}
pertanto, perché $A_1 \cap B_1 \in \tau_1$ e $A_2 \cap B_2 \in \tau_2$, per la \textit{proposizione 1.1.2}, si ottiene che $\mathcal{B}$ è base di un'unica topologia $\tau$ su $X$, detta \textit{topologia prodotto} (delle topologie $\tau_1$ e $\tau_2$). \\ \\
Per definizione, gli aperti di $X$ sono unione di elementi di $\mathcal{B}$; segue immediatamente che un sottoinsieme $N$ di $X$ è un intorno del punto $(x_1, x_2)$ se e solo se esistono $N_1$ e $N_2$ intorni di $x_1$ e $x_2$ rispettivamente tali che $N \supset N_1 \times N_2$. \\ \\
Ragionando analogamente si può definire una topologia prodotto per il prodotto cartesiano di un numero finito di spazi topologici. \\ \\ \\
\textbf{Proposizione 1.1.5:} siano $\mathcal{B}_1$ e $\mathcal{B}_2$ basi degli spazi topologici $X_1$ e $X_2$, rispettivamente. La famiglia di sottoinsiemi
\begin{gather}
    \mathcal{B} := \{ B_1 \times B_2 : B_1 \in \mathcal{B}_1, B_2 \in \mathcal{B}_2 \} \notag
\end{gather}
è una base per la topologia prodotto di $X = X_1 \times X_2$. \\

\textbf{Dimostrazione:} si verifica immediatamente che gli aperti di $X$ sono unione degli aperti di $\mathcal{B}$. \qed \\ \\ \\
\textbf{Definizione 1.1.14:} siano $X$ uno spazio topologico, $\{ x_n \}_{n \in \mathbf{N}} \subseteq X$ e $y \in X$. Si dice che $y$ è un limite per $\{ x_n \}_{n \in \mathbf{N}}$ per $n \rightarrow \infty$ se per ogni $V \in \mathcal{N}(y)$ esiste $M \in \mathbf{N}$ tale che $x_n \in V$ per ogni $n \geq M$, in questo caso scriviamo $x_n \rightarrow y$, per $n \rightarrow \infty$. \\ \\
\textbf{Esempio:} in topologia banale si ha che ogni successione converge ad ogni punto, poiché l'unico intorno è $X$ stesso. \\ \\ \\
\textbf{Definizione 1.1.15:} siano $X$ uno spazio topologico e $Y \subset X$, $Y$ è detto \textit{denso} se $\overline{Y} = X$. \\ \\ \\
\textbf{Definizione 1.1.16:} sia $X$ uno spazio topologico, $X$ è detto \textit{separabile} se esiste $Y \subset X$ denso e numerabile (numerabile se esiste biezione con $\mathbf{N}$).
\subsubsection{Esercizi}
\textbf{1.} siano $X := \mathbf{R}$, $\tau := \{ \emptyset, X \} \cup \{ [q, +\infty) : q \in \mathbf{R} \}$, verificare se $\tau$ è una topologia su $\mathbf{R}$. \\

\textbf{Soluzione:} siano $q, p \in \mathbf{R}$, allora
\begin{gather}
    [q, +\infty) \cap [p, +\infty) = [M, +\infty) \notag
\end{gather}
con $M := \max \{q, p\}$, pertanto è ancora un elemento della topologia, dunque $\tau$ è chiusa rispetto all'intersezione finita. \\
Tuttavia notiamo che
\begin{gather}
    \bigcup_{q > 1} [q, +\infty) = (1, +\infty) \not\in \tau \notag
\end{gather}
pertanto $\tau$ non è una topologia su $\mathbf{R}$. \\ \\ \\
\textbf{2.} siano $X$ un insieme non vuoto e $\tau := \{ \emptyset, X \} \cup \{ X \setminus A : A \subseteq X, |A| < +\infty \}$, verificare se $\tau$ è una topologia su $X$. \\

\textbf{Soluzione:} siano $A, B \subseteq X$ di cardinalità finita, si ha che
\begin{gather}
    X \setminus A \cap X \setminus B = X \setminus (A \cup B) \in \tau \text{ poiché } |A \cup B| < +\infty \notag
\end{gather}Consideriamo ora una famiglia infinita di insiemi di cardinalità infinita $\{ A_i \}_{i \in I}$, si ha che
\begin{gather}
    \bigcup_{i \in I} (X \setminus A_i) = X \setminus \bigcap_{i \in I} A_i \in \tau \notag
\end{gather}
infatti $|\bigcap_{i \in I} A_i| \leq |A_i| < +\infty$ per ogni $i \in I$, concludiamo che $\tau$ è una topologia su $X$, detta \textit{topologia cofinita}. \\ \\ \\
\textbf{3.} Siano $X := \mathbf{N}$ e $\tau := \{ \emptyset, X \} \cup \{q, q+1, \dots : q \in \mathbf{N} \}$. Determinare l'interno di $Y := \{ 5, 11, 12, 25 \}$. \\

\textbf{Soluzione:} per definizione si ha che $\mathring{Y}$ è dato dal più "grande" aperto contenuto in $Y$. Ovviamente vale $\emptyset \subset Y$, notiamo inoltre che per ogni $p \in \mathbf{N}$ si ha che $\{ p, p+1, \dots \} \not\subset Y$, infatti tutti gli aperti di $\tau$ (eccetto $\emptyset$) sono illimitati. Pertanto concludiamo che $\mathring{Y} = \emptyset$. \\ \\ \\
\textbf{4.} Siano $X := \mathbf{R}$ e $\tau := \{ \emptyset, X \} \cup \{ (q, +\infty) : q \in \mathbf{R} \}$. Determinare $\overline{\mathbf{Z}}$ come sottoinsieme di $\mathbf{R}$, successivamente determinare $D(\mathbb{\overline{Z}})$. \\

\textbf{Soluzione:} i chiusi per la topologia $\tau$ definita sono della forma $(-\infty, q]$ con $q \in \mathbf{R}$. Per ogni $q < +\infty$ si ha che $\mathbf{Z} \not\subset (-\infty, q]$, pertanto il più "piccolo" chiuso che contiene $\mathbf{Z}$ è $\mathbf{R}$, dunque $\overline{\mathbf{Z}} = \mathbf{R}$. \\ \\
Sia $x_0 \in \mathbf{R}$, consideriamo un intorno $V \in U(x_0)$, poiché $\overline{\mathbf{Z}} = \mathbf{R}$ si ha che $\overline{\mathbf{Z}} \cap V \setminus \{ x_0 \} \neq \emptyset$, ciò vale per ogni $x_0 \in \mathbf{R}$, dunque $D(\mathbf{Z}) = \mathbf{R}$. \\ \\ \\
\textbf{5.} Determinare $D([4, 9])$ come sottoinsieme di $\mathbf{R}$ per la topologia definita nell'esercizio precedente. \\

\textbf{Soluzione:} sia $p \leq 9$, per ogni $\varepsilon > 0$ l'aperto $(p - \varepsilon, +\infty) \in U(p)$ e verifica la condizione
\begin{gather}
    [4, 9] \cap (p - \varepsilon, +\infty) \setminus \{ p \} \neq \emptyset \notag
\end{gather}
pertanto $(-\infty, 9] \subset D([4, 9])$. Sia ora $p > 9$, si ha che esiste $\varepsilon > 0$ tale che l'intorno aperto $(p - \varepsilon, +\infty) \cap [4, 9] = \emptyset$, basta scegliere $\varepsilon = \frac{1}{2}(p - 9)$, pertanto $D([4, 9]) = (-\infty, 9]$. \\

Ovviamente questo risultato può essere generalizzato a qualsiasi intervallo chiuso della retta reale: siano $a, b \in \mathbf{R}$, allora $D([a, b]) = (-\infty, b]$. \\ \\ \\
\textbf{6.} Siano $X$ uno spazio topologico e $A \subseteq X$. Verificare se è sempre vero che $\overline{A} = \mathring{A} \cup \partial(A)$. \\

\textbf{Soluzione:} per definizione si ha $\partial(A) = \overline{X \setminus A} \cap \overline{A} = (X \setminus \mathring{A}) \cap \overline{A}$. Pertanto otteniamo:
\begin{gather}
    \mathring{A} \cup [(X \setminus \mathring{A}) \cap \overline{A}] = X \cap \overline{A} = \overline{A} \notag
\end{gather}
che è la tesi. Notiamo inoltre che ciò equivale a $\partial(A) = \overline{A} \setminus \mathring{A}$. \\ \\ \\
\textbf{7.} Sia $X$ uno spazio topologico, dimostrare che comunque dati due sottoinsiemi $A, B \in X$ si ha: \\
\textbf{i.} $\partial (A \cap B) \subset \partial (A) \cap \partial (B)$ (N.B. sbagliato sul Sernesi). \\
\textbf{ii.} $\partial (A \cup B) \subset \partial (A) \cup \partial (B)$. \\
\textbf{iii.} $\mathring{(A \cap B)} = \mathring{A} \cap \mathring{B}$. \\
\textbf{iv.} $\mathring{(A \cup B)} \supset \mathring{A} \cup \mathring{B}$. \\
\textbf{v.} $\overline{A \cup B} = \overline{A} \cup \overline{B}$. \\
\textbf{vi.} $\overline{A \cap B} \subset \overline{A} \cap \overline{B}$. \\

\textbf{Dimostrazione:} \textbf{iii.} se $x \in \mathring{A} \cap \mathring{B}$ si ha $x \in \mathring{A}$ e $x \in \mathring{B}$, quindi per definizione esistono $S, R \in \tau$ tali che $S \subseteq A$ e $R \subseteq B$ e $x \in S \cap R$. Pertanto, poiché $S \cap R \in \tau$ e $S \cap R \subseteq A \cap B$, si ha, per definizione, che $x \in \mathring{A \cap B}$, dunque $\mathring{A} \cap \mathring{B} \subseteq \mathring{A \cap B}$. Viceversa abbiamo che $\mathring{A \cap B} \subseteq A \cap B$, quindi $\mathring{A \cap B} \subseteq A$ e $\mathring{A \cap B} \subseteq B$. Essendo $\mathring{A}$ e $\mathring{B}$ i più "grandi" aperti contenuti rispettivamente in $A$ e $B$ si ha che $\mathring{A \cap B} \subseteq \mathring{A}$ e $\mathring{A \cap B} \subseteq \mathring{B}$, pertanto $\mathring{A \cap B} \subseteq \mathring{A} \cap \mathring{B}$. Combinando i risultati otteniamo $\mathring{A \cap B} = \mathring{A} \cap \mathring{B}$. \\ \\
\textbf{iv.} Se $x \in \mathring{A} \cup \mathring{B}$ si ha che esiste, per definizione, $S \subseteq A$ aperto e/o $R \subseteq B$ aperto tale che $x \in S$ e/o $x \in R$, pertanto $x \in S \cup R$. Poiché $S \cup R \subset A \cup B$ ed è aperto si ha, per definizione, che $x \in \mathring{A \cup B}$. \\ \\
\textbf{v.} Se $x \in \overline{A} \cup \overline{B}$ si ha che $x \in S$ e/o $x \in R$ per ogni $S \supseteq A$ chiuso e/o per ogni $R \supseteq B$ chiuso, pertanto otteniamo che $x \in S \cup R$ per ogni $S, R$. Poiché $S \cup R$ è un chiuso contenente $A \cup B$ per ogni $S, R$ otteniamo, per definizione, che $y \in \overline{A \cup B}$, dunque $\overline{A} \cup \overline{B} \subseteq \overline{A \cup B}$. Viceversa, si ha che $\overline{A \cup B} \supseteq A \cup B$, dunque $\overline{A \cup B} \supseteq A$ e $\overline{A \cup B} \supseteq B$, inoltre, poiché $\overline{A}$ e $\overline{B}$ sono i più "piccoli" chiusi contenenti rispettivamente $A$ e $B$ si ha che $\overline{A \cup B} \supseteq \overline{A}$ e $\overline{A \cup B} \supseteq \overline{B}$, pertanto si ottiene $\overline{A \cup B} \supseteq \overline{A} \cup \overline{B}$. Combinando i risultati si ha la tesi: $\overline{A \cup B} = \overline{A} \cup \overline{B}$. \\ \\
\textbf{vi.} Se $x \in \overline{A \cap B}$ si ha che $x \in Q$ per ogni $Q \supseteq A \cap B$ chiuso. Dato che $A \cap B \subseteq A$ e $A \cap B \subseteq B$ si ottiene che la classe dei chiusi $Q_A$ che contengono $A$ e la classe dei chiusi $Q_B$ che contengono $B$ sono sottoinsiemi della classe dei chiusi $Q$ che contengono $A \cap B$, pertanto $x \in Q_A$ e $x \in Q_B$ per ogni $Q_A \supseteq A$ e per ogni $Q_B \supseteq B$. Per definizione si ottiene che $x \in \overline{A} \cap \overline{B}$, dunque $\overline{A \cap B} \subset \overline{A} \cap \overline{B}$. \\ \\
\textbf{i.} Si ha che $\partial (A \cap B) = \overline{A \cap B} \setminus \mathring{A \cap B} \subset (\overline{A} \cap \overline{B}) \setminus (\mathring{A} \cap \mathring{B})$ e si ha che:
\begin{align}
    (\overline{A} \cap \overline{B}) \setminus (\mathring{A} \cap \mathring{B}) &= ((\overline{A} \cap \overline{B}) \setminus \mathring{A}) \cup ((\overline{A} \cap \overline{B}) \setminus \mathring{B}) \notag \\
    &= ((\overline{A} \setminus \mathring{A}) \cap (\overline{B} \setminus \mathring{A})) \cup ((\overline{B} \setminus \mathring{B}) \cap (\overline{A} \setminus \mathring{B})) \notag
\end{align}
pertanto si vede facilmente che $((\overline{A} \setminus \mathring{A}) \cap (\overline{B} \setminus \mathring{A})) \cup ((\overline{B} \setminus \mathring{B}) \cap (\overline{A} \setminus \mathring{B})) \subset\\\subset (\overline{A} \setminus \mathring{A}) \cup (\overline{B} \setminus \mathring{B}) = \partial (A) \cup \partial (B)$, che è la tesi. \\ \\
\textbf{ii.} Si ha che $\partial (A \cup B) = \overline{X \setminus (A \cup B)} \cap \overline{A \cup B}$, dunque:
\begin{gather}
    \overline{X \setminus (A \cup B)} \cap \overline{A \cup B} \subset \overline{(X \setminus A) \cap (X \setminus B)} \cap (\overline{A} \cup \overline{B}) \subset \notag \\
    \subset (\overline{X \setminus A} \cap \overline{X \setminus B}) \cap (\overline{A} \cup \overline{B}) \subset (\overline{X \setminus A} \cap \overline{A}) \cup (\overline{X \setminus B} \cap \overline{B}) \notag
\end{gather}
che è la tesi. \\ \\ \\
\textbf{8.} Siano $X$ uno spazio topologico, $M \subseteq X$. Dimostrare $M$ è chiuso se e solo se $\partial(M) \subset (M)$. \\

\textbf{Dimostrazione:} $(\rightarrow)$ sia $M$ chiuso, allora $M = \overline{M} = \mathring{M} \cup \partial(M)$, per l'\textit{esercizio 6}, dunque $\partial(M) \subseteq M$. \\ \\
$(\leftarrow)$ sia $\partial(M) \subseteq M$, per l'\textit{esercizio 6} si ha che $\partial(M) = \overline{M} \setminus \mathring{M}$, per definizione $\mathring{M} \subset M$, dunque anche $\mathring{M} \cup \partial(M) = \overline{M} \subset M$, cioè $M = \overline{M}$, dunque $M$ è chiuso.  \\ \\ \\
\textbf{9.} Sia $f : X \longrightarrow Y$ un'applicazione di spazi topologici, $x \in X$ e $\mathcal{B}(f(x))$ un sistema fondamentale di intorni di $f(x)$. Dimostrare che $f$ è continua in $x$ se e solo se per ogni $N \in \mathcal{B}(f(x))$ esiste $M \in \mathcal{N}(x)$ tale che $f(M) \subset N$. \\

\textbf{Dimostrazione:} $(\rightarrow)$ sia $f$ continua in $x$, per definizione si ha che per ogni $V \in \mathcal{N}(f(x))$ esiste $W \in \mathcal{N}(x)$ tale che $f(W) \subset V$, pertanto, dato che $\mathcal{B}(f(x)) \subseteq \mathcal{N}(f(x))$, ciò vale anche per ogni $N \in \mathcal{B}(f(x))$. \\ \\
$(\leftarrow)$ Supponiamo ora che per ogni $N \in \mathcal{B}(f(x))$ esiste $M \in \mathcal{N}(x)$ tale che $f(M) \subset N$. Per ogni $V$ intorno di $f(x)$, per definizione, esiste $W \in \mathcal{B}(f(x))$ tale che $W \subset V$, dunque, per ipotesi, si ha che esiste $M \in \mathcal{N}(x)$ tale che $f(M) \subset W \subset V$, quindi $f$ è continua in $x$.

\newpage
\subsection{Spazi metrici e proprietà di separazione}
\textbf{Definizione 1.2.1:} sia $X$ un insieme, sia $d_X : X \times X \longrightarrow [0, + \infty)$, $d_X$ è detta \textit{metrica} (o \textit{distanza}) su $X$ se: \\
\textbf{1.} $d(x_1, x_2) \geq 0$ per ogni $x_1, x_2 \in X$. \\
\textbf{2.} $d(x_1, x_2) = 0$ se e solo se $x_1 = x_2$. \\
\textbf{3.} $d(x_1, x_2) = d(x_2, x_1)$ per ogni $x_1, x_2 \in X$. \\
\textbf{4.} $d(x_1, x_2) \leq d(x_1, x_3) + d(x_3, x_2)$ per ogni $x_1, x_2, x_3 \in X$. \\ \\
La coppia $(X, d_X)$ è detta \textit{spazio metrico}. \\ \\ \\
\textbf{Definizione 1.2.2:} sia $(X, d_X)$ uno spazio metrico, è detta \textit{sfera aperta di raggio $r$ e centro $x_0$} l'insieme $D_r(x_0) := \{ x \in X : d_X(x, x_0) < r \}$. \\ \\
La metrica induce una topologia sullo spazio, sia infatti $(X, d_X)$ uno spazio metrico, definiamo $\tau$ come segue:
\begin{gather}
    A \in \tau \iff \forall x_0 \in A \text{ esiste } \varepsilon > 0 \text{ tale che } B_\varepsilon(x_0) \subseteq A \notag
\end{gather}
ovvero se e solo se $A$ è unione di sfere aperte. \\ \\
\textbf{Esempi:} - sia $X := \mathbf{R}^n$, la metrica $d(x_1, x_2) := \sqrt{\sum_{i = 1}^n (x_1^i - x_2^i)^2}$ per ogni $x_1, x_2 \in \mathbf{R}^n$ è detta \textit{metrica euclidea} e induce la topologia euclidea. \\
- sia $X := \mathbf{R}^n$, la metrica $d(x_1, x_2) := \sum_{i = 1}^n |x_1^i - x_2^i|$ per ogni $x_1, x_2 \in \mathbf{R}^n$ induce ancora una volta la metrica euclidea (in questo caso gli aperti sono rombi). \\
- siano $a, b \in \mathbf{R}$, $X := C^0([a, b])$, poniamo $d(f, g) := \max_{t \in [a, b]} |f(t) - g(t)|$ per ogni $f, g \in C^0([a, b])$, si verifica facilmente che $d$ è una metrica e dunque induce una topologia su $C^0([a, b])$. \\ \\ \\
\textbf{Definizione 1.2.3:} sia $V$ uno spazio vettoriale reale, una funzione $\| \cdot \| : V \times V \longrightarrow [0, +\infty)$ è detta \textit{norma} se: \\
\textbf{1.} $\| v \| \geq 0$ per ogni $v \in V$ e $\| v \| = 0$ se e solo se $v = \vec{0}$. \\
\textbf{2.} $\| v + w \| \leq \| v \| + \| w \|$ per ogni $v, w \in V$. \\
\textbf{3.} $\| \lambda v \| = |\lambda| \| v \|$ per ogni $v \in V$ e per ogni $\lambda \in \mathbf{R}$. \\ \\
La coppia $(V, \| \cdot \|)$ è detto \textit{spazio normato}. \\ \\
Notiamo che la norma induce una metrica sullo spazio $V$, infatti ponendo
\begin{gather}
    d(v, w) := \|v - w\| \notag
\end{gather}
si ottiene una metrica sullo spazio $V$. \\ \\ \\
\textbf{Definizione 1.2.4:} sia $X$ uno spazio topologico, $X$ è detto $T_2$ (o \textit{separato} o \textit{di Hausdorff}) se per ogni $x, y \in X$ con $x \neq y$ esistono $V \in \mathcal{N}(x)$ e $W \in \mathcal{N}(y)$ tali che $V \cap W = \emptyset$. \\ \\
Se $X$ è uno spazio metrico segue che $X$ è di Hausdorff, infatti basta porre $\varepsilon := \frac{1}{4}d(x, y)$ e prendere come intorni $V := D_\varepsilon(x)$ e $W := D_\varepsilon(y)$ per ogni $x, y \in X$. \\ \\ \\
\textbf{Proposizione 1.2.1:} se $Y$ è un sottoinsieme di uno spazio di Hausdorff è uno spazio di Hausdorff. \\

\textbf{Dimostrazione:} siano $p, q \in Y \subset X$, per ipotesi esistono $V \in \mathcal{N}(p)$ e $W \in \mathcal{N}(q)$ tali che $V \cap W = \emptyset$. Pertanto gli insiemi $Y \cap V$ e $Y \cap W$ sono rispettivamente intorni di $p$ e $q$ nella topologia indotta e si ha
\begin{gather}
    (Y \cap V) \cap (Y \cap W) = V \cap W \cap Y = \emptyset \cap Y = \emptyset \notag
\end{gather}
dunque $Y$ è uno spazio di Hausdorff. \qed \\ \\ \\
\textbf{Proposizione 1.2.2:} se $X$ è uno spazio di Hausdorff i punti sono chiusi. \\

\textbf{Dimostrazione:} sia $x \in X$, consideriamo l'insieme $X \setminus \{ x \}$ e mostriamo che è aperto. Per ogni $y \in X \setminus \{ x \}$ si ha, per ipotesi, che esistono $V \in \mathcal{N}(x)$ e $W \in \mathcal{N}(y)$ tali che $V \cap W = \emptyset$, allora si ha che $W \subseteq X \setminus V \subseteq X \setminus \{ x \}$. \qed \\ \\ \\
\textbf{Proposizione 1.2.3:} sia $X$ uno spazio di Hausdorff, allora se il limite di una successione esiste è unico. \\

\textbf{Dimostrazione:} sia $\{ x_n \}_{n \in \mathbf{N}} \subseteq X$ una successione, supponiamo per assurdo che essa converga in due punti $y_1$ e $y_2$ di $X$. Per ipotesi si ha che esistono $V \in \mathcal{N}(y_1)$ e $W \in \mathcal{N}(y_2)$ tali che $V \cap W = \emptyset$. Per definizione di limite esiste $M \in \mathbf{N}$ tale che $x_n \in V$ e $x_n \in W$ per ogni $n \geq M$, ma ciò è un assurdo. \qed \\ \\ \\
\textbf{Proposizione 1.2.4:} il prodotto topologico di una famiglia qualsiasi $\{ X_\mu \}_{\mu \in M}$ di spazi di Hausdorff è uno spazio di Hausdorff. \\

\textbf{Dimostrazione:} sia
\begin{gather}
    X := \prod_{\mu \in M} X_\mu \notag
\end{gather}
consideriamo due punti $p, q \in X$, $p \neq q$, allora esiste almeno un indice $\overline{\mu} \in M$ tale che $p(\overline{\mu}) \neq q(\overline{\mu})$. Per ipotesi $X_{\overline{\mu}}$ è uno spazio di Hausdorff, dunque esistono $V_{\overline{\mu}} \in \mathcal{N}(p(\overline{\mu}))$ e $W_{\overline{\mu}} \in \mathcal{N}(q(\overline{\mu}))$ tali che $V_{\overline{\mu}} \cap W_{\overline{\mu}} = \emptyset$. \\
Consideriamo per ogni $\mu \in M$, $\mu \neq \overline{\mu}$ gli intorni $V_\mu \in \mathcal{N}(p(\mu))$ e $W_\mu \in \mathcal{N}(q(\mu))$, si ha che
\begin{gather}
    V_{\overline{\mu}} \times \prod_{\mu \neq \overline{\mu}} V_\mu \in \mathcal{N}(p), \quad W_{\overline{\mu}} \times \prod_{\mu \neq \overline{\mu}} W_\mu \in \mathcal{N}(q) \notag
\end{gather}
inoltre
\begin{align}
    (V_{\overline{\mu}} \times \prod_{\mu \neq \overline{\mu}} V_\mu) \cap W_{\overline{\mu}} \times (\prod_{\mu \neq \overline{\mu}} W_\mu) &= (V_{\overline{\mu}} \cap W_{\overline{\mu}}) \times \prod_{\mu \neq \overline{\mu}} (V_\mu \cap W_\mu) = \notag \\
    &= \emptyset \times \prod_{\mu \neq \overline{\mu}} (V_\mu \cap W_\mu) = \emptyset \notag
\end{align}
Concludiamo che anche $X$ è uno spazio di Hausdorff. \qed \\ \\ \\
\textbf{Proposizione 1.2.5:} sia $X$ uno spazio topologico. Allora $X$ è di Hausdorff se e solo se $\Delta(X) := \{ (x_1, x_2) \in X \times X : x_1 = x_2 \}$ è chiuso in $X \times X$. \\

\textbf{Dimostrazione:} $(\rightarrow)$ Sia $X$ uno spazio di Hausdorff per ipotesi, siano $x, y \in X$, $x \neq y$. Per ipotesi esistono $V_x \in \mathcal{N}(x)$ e $V_y \in \mathcal{N}(y)$ tali che $V_x \cap V_y = \emptyset$. Consideriamo ora $(x, y) \in X \times X$, si ha che $(x, y) \in V_x \times V_y \in \mathcal{N}((x, y))$, inoltre $V_x \times V_y \cap \Delta (X) = \emptyset$ proprio perché $V_x \cap V_y = \emptyset$, pertanto $V_x \times V_y \subseteq X \times X \setminus \Delta (X)$, dunque, per la \textit{proposizione 1.1.1}, si ha che $X \times X \setminus \Delta (X)$ è un aperto. \\ \\
$(\leftarrow)$ Sia $\Delta (X)$ un chiuso di $X \times X$ per ipotesi. Allora $X \times X \setminus \Delta (X)$ è un aperto, pertanto per ogni $(x, y) \in X \times X \setminus \Delta (X)$ esistono $V_x \in \mathcal{N}(x)$ e $V_y \in \mathcal{N}(y)$ tali che $(x, y) \in V_x \times V_y \subseteq X \times X \setminus \Delta (X)$, ciò comporta che $V_x \cap V_y = \emptyset$ proprio per come è definito $\Delta (X)$. Poiché ciò vale per ogni $(x, y) \in X \times X \setminus \Delta (X)$ si ottiene la tesi. \qed \\ \\ \\
\textbf{Definizione 1.2.5:} uno spazio topologico $X$ è detto $T_1$ se i punti sono chiusi. \\ \\ \\
\textbf{Proposizione 1.2.6:} uno spazio topologico $X$ è $T_1$ se e solo se per ogni coppia di punti distinti $p, q \in X$ esistono $V, W$ aperti tali che $p \in V$ ma $q \not\in V$, $q \in W$ ma $p \not\in W$. \\

\textbf{Dimostrazione:} $(\rightarrow)$ sia $X$ uno spazio $T_1$, ovvero i punti sono chiusi, allora $X \setminus \{ p \}$ è un aperto e $q \in X \setminus \{ p \}$ ma $p \not\in X \setminus \{ p \}$. Analogamente, $X \setminus \{ q \}$ è un aperto e $p \in X \setminus \{ q \}$ ma $q \not\in X \setminus \{ q \}$. \\ \\
$(\leftarrow)$ supponiamo che per ogni coppia di punti distinti $p, q \in X$ esistono $V, W$ aperti tali che $p \in V$ ma $q \not\in V$ e $q \in W$ ma $p \not\in W$. Sia $u \in X$, per ogni $v \in X$ tale che $u \neq v$ esistono $V, W$ aperti tali che $u \in V \subset X \setminus W \subset X \setminus \{ v \}$, pertanto $X \setminus \{ v \}$ è un aperto, quindi i punti sono chiusi. \qed \\ \\ \\
\textbf{Proposizione 1.2.7:} sia $\{ x_n : n \in \mathbf{N} \}$ una successione di punti in uno spazio $X$ $T_1$ che soddisfa il primo assioma di numerabilità e sia $x \in X$ un punto di accumulazione per $\{ x_n \}$. Allora $\{ x_n \}$ ammette una sottosuccessione convergente in $x$. \\

\textbf{Dimostrazione:} ogni sottoinsieme finito di $\{ x_n \}$ è discreto (i punti sono chiusi), dunque non ammette punti di accumulazione, pertanto $x$ è punto di accumulazione per $\{ x_n : n \geq N \}$ per ogni $N \geq 1$. Sia $U := \{ U_k : k \in \mathbf{N} \}$ un sistema fondamentale di intorni numerabile di $x$ tale che $U_{k + 1} \subset U_k$ per ogni $k \geq 1$. Per ipotesi esiste $n_1 \in \mathbf{N}$ tale che $x_{n_1} \in U_1$. Analogamente esiste $n_2 \in \mathbf{N}$, $n_2 > n_1$, tale che $x_{n_2} \in U_2$. Ragionando per induzione possiamo costruire una sottosuccessione $\{ x_{n_k} : k \in \mathbf{N} \}$ tale che $x_{n_k} \in U_k$ per ogni $k \geq 1$, inoltre è chiaro che $x = \lim_{k \rightarrow \infty} x_{n_k}$. \qed \\ \\ \\
\textbf{Definizione 1.2.6:} uno spazio \textit{regolare} è uno spazio topologico $X$ che soddisfa le seguenti proprietà: \\
\textbf{1.} $X$ è $T_1$. \\
\textbf{2.} Per ogni sottoinsieme chiuso $F$ di $X$ e ogni $x \in X \setminus F$ esistono $V, W$ aperti tali che $x \in V$, $F \subset W$ e $V \cap W = \emptyset$. \\ \\ \\
\textbf{Proposizione 1.2.8:} ogni spazio regolare è uno spazio di Hausdorff. \\

\textbf{Dimostrazione:} segue direttamente dal fatto che uno spazio regolare è $T_1$, dunque i punti sono sottoinsiemi chiusi. \qed \\ \\ \\
\textbf{Definizione 1.2.7:} uno spazio \textit{normale} è uno spazio topologico $X$ che soddisfa le seguenti proprietà: \\
\textbf{1.} $X$ è $T_1$. \\
\textbf{2.} per ogni coppia di sottoinsiemi chiusi $F_1, F_2$ di $X$ esistono $V_1, V_2$ aperti tali che $F_1 \subset V_1$, $F_2 \subset V_2$ e $V_1 \cap V_2 = \emptyset$. \\ \\ \\
\textbf{Proposizione 1.2.9:} ogni spazio normale è uno spazio di Hausdorff. \\

\textbf{Dimostrazione:} segue direttamente dal fatto che uno spazio normale è $T_1$, dunque i punti sono sottoinsiemi chiusi. \qed \\ \\ \\
\textbf{Teorema 1.2.1:} ogni spazio metrizzabile $X$ è uno spazio normale. \\

\textbf{Dimostrazione:} sappiamo che ogni spazio metrizzabile è uno spazio di Hausdorff e dunque $T_1$. Per ogni sottoinsieme non vuoto $F$ di $X$ definiamo
\begin{gather}
    d_F(x) = \inf_{y \in F} d(x, y) \notag
\end{gather}
Dimostriamo ora un lemma necessario nel seguito. \\ \\
\textit{Lemma:} siano $X$ uno spazio metrizzabile, $M \subset X$ un insieme non vuoto, $C \subset X$ un chiuso non vuoto tale che $M \cap C = \emptyset$ e $p \in M$ (fissato). Allora $d_C(p) > 0$, ovvero esiste $\varepsilon > 0$ tale che $d(p, q) > \varepsilon$ per ogni $q \in C$. \\

\textit{Dimostrazione:} $C$ è chiuso, dunque $X \setminus C$ è aperto e $M \subset X \setminus C$. Per ipotesi esiste $\delta > 0$ tale che $D_\delta (p) \subset X \setminus C$, ovvero $d(p, q) > \delta$ per ogni $q \in C$, pertanto, scegliendo $\varepsilon = \delta$, si ottiene la tesi. \\ \\
Sia $x \in F_1$, allora $d_{F_1}(x) = 0$, mentre $d_{F_2}(x) > 0$, poiché $x \not\in F_2$, sia $r := d_{F_2}(x)$. Costruiamo due aperti disgiunti $W_1, W_2$ tali che $F_1 \subset W_1$ e $F_2 \subset W_2$, poniamo
\begin{gather}
    W_1 := \{ x \in X : d_{F_1}(x) - d_{F_2}(x) < 0 \} \notag \\
    W_2 := \{ x \in X : d_{F_1}(x) - d_{F_2}(x) > 0 \} \notag
\end{gather}
Si vede immediatamente che $W_1 \cap W_2 = \emptyset$. Inoltre si ha
\begin{gather}
    d_{F_1}(x) - d_{F_2}(x) = -d_{F_2}(x) < 0 \quad \forall x \in F_1 \notag
\end{gather}
e analogamente
\begin{gather}
    d_{F_1}(y) - d_{F_2}(y) = d_{F_1}(y) > 0 \quad \forall y \in F_2 \notag
\end{gather}
pertanto $F_1 \subset W_1$ e $F_2 \subset W_2$. Mostriamo ora che sono entrambi aperti, sia $x_0 \in W_1$, si ha che $r := d_{F_2}(x_0) - d_{F_1}(x_0) > 0$. Consideriamo la sfera aperta $D_{\frac{r}{3}}(x_0)$, per ogni $x \in D_{\frac{r}{3}}(x_0)$ si ha che
\begin{align}
    d_{F_1}(x) - d_{F_2}(x) &= [d_{F_1}(x) - d_{F_1}(x_0)] + [d_{F_1}(x_0) - d_{F_2}(x_0)] + [d_{F_2}(x_0) - d_{F_2}(x)] \notag \\
    &= [d_{F_1}(x) - d_{F_1}(x_0)] - r + [d_{F_2}(x_0) - d_{F_2}(x)] \notag \\
    &< \frac{r}{3} - r + \frac{r}{3} = -\frac{r}{3} < 0 \notag
\end{align}
per la disuguaglianza triangolare, dunque $D_{\frac{r}{3}}(x_0) \subset W_1$, poiché ciò vale comunque scelto $x_0 \in W_1$ si ha che $W_1$ è aperto. Analogamente si dimostra per $W_2$, dunque si ottiene la tesi. \qed \\ \\ \\
Le proprietà di separazione che uno spazio $X$ può avere sono collegate all'esistenza di applicazioni continue su $X$ a valori reali. \\
Ad esempio, se uno spazio $X$ ha la proprietà che per ogni coppia di punti distinti $p, q \in X$ esiste un'applicazione continua $f : X \longrightarrow \mathbf{R}$ tale che $f(p) \neq f(q)$, si ha che $X$ è uno spazio di Hausdorff, infatti ponendo $r := d(f(p), f(q))$ e
\begin{gather}
    V := f^{-1}(f(p) - \frac{r}{2}, f(p) + \frac{r}{2}) \notag \\
    W := f^{-1}(f(q) - \frac{r}{2}, f(q) + \frac{r}{2}) \notag
\end{gather}
$V$ e $W$ sono aperti disgiunti tali che $p \in V$ e $q \in W$.
\subsubsection{Esercizi}
\textbf{1.} Dimostrare che essere $T_1$, rispettivamente $T_2$, regolare e normale, è una proprietà topologica. \\

\textbf{Dimostrazione:} siano $X \cong Y$ e $f : X \longrightarrow Y$ un omeomorfismo. Sia $X$ uno spazio $T_1$, siano $p, q \in X$, $p \neq q$, per ipotesi esistono $V, W$ aperti di $X$ tali che $p \in V$, $q \not\in V$ e $q \in W$, $p \not\in W$. Si ha che $f(V), f(W)$ sono aperti in $Y$ (essere aperto è una proprietà topologica) tali che $f(p) \in f(V)$, $f(q) \not\in f(V)$ e $f(Q) \in f(W)$, $f(p) \not\in f(W)$, pertanto si ottiene la tesi. \\ \\
Ora sia $X$ uno spazio $T_2$, per ipotesi esistono $N \in \mathcal{N}(p)$ e $M \in \mathcal{N}(q)$ tali che $N \cap M = \emptyset$. Si ha che $f(N) \in \mathcal{N}(f(p))$ e $f(M) \in \mathcal{N}(f(q))$ (essere un intorno è una proprietà topologica) e $f(N) \cap f(M) = \emptyset$, per l'iniettività di $f$, pertanto segue la tesi. \\ \\
Sia ora $X$ uno spazio regolare, consideriamo un sottoinsieme chiuso $F$ di $X$ e un punto $x \in X \setminus F$, per ipotesi esistono $A, B$ aperti disgiunti tali che $F \subset A$ e $x \in B$. Si ha che $f(F) \subset f(A)$, $f(x) \in f(B)$ e $f(A) \cap f(B) = \emptyset$ per l'iniettività di $f$, pertanto segue la tesi. \\ \\
Sia ora $X$ normale, consideriamo due sottoinsiemi chiusi $F_1, F_2$ di $X$, per ipotesi esistono $U_1, U_2$ aperti disgiunti tali che $F_1 \subset U_1$ e $F_2 \subset U_2$. Si ha che $f(F_1) \subset f(U_1)$, $f(F_2) \subset f(U_2)$ e $f(U_1) \cap f(U_2) = \emptyset$ per l'iniettività di $f$, pertanto segue la tesi. \\ \\ \\
\textbf{2.} Dimostrare che un insieme $X$ infinito con topologia cofinita non è metrizzabile. \\

\textbf{Dimostrazione:} sia $X$ uno spazio topologico infinito e sia 
\begin{gather}
    \tau := \{ \emptyset, X \} \cup \{ X \setminus A : A \subset X, |A| < +\infty \} \notag
\end{gather}
Sappiamo che ogni spazio metrizzabile è anche normale e di Hausdorff. Siano $F_1, F_2 \subset X$ chiusi e disgiunti, gli aperti che contengono $F_1$ sono della forma $X \setminus M$ con $M \subset F_1$, analogamente gli aperti che contengono $F_2$ sono della forma $X \setminus N$ con $N \subset F_2$. Si ha che $X \setminus M \cap X \setminus N = X \setminus (M \cup N) = \emptyset$ se e solo se $M \cup N = X$, ma $M$ e $N$ sono insiemi finiti poiché sono contenuti rispettivamente in $F_1$ e $F_2$ che sono chiusi per la topologia cofinita, pertanto sono anch'essi finiti, dunque si ottiene un assurdo. \\ \\ \\
\textbf{3.} Dimostrare che $X$ è di Hausdorff se e solo se per ogni $x \in X$ si ha $\{ x \} = \bigcap_{M \in \mathcal{N}(x)} \overline{M}$. \\

\textbf{Dimostrazione:} $(\rightarrow)$ sia $X$ di Hausdorff, sia $p \in X$, per ogni $q \in X$, $q \neq p$, esistono $V_q \in \mathcal{N}(p)$ e $W_q \in \mathcal{N}(q)$ aperti tali che $V_q \cap W_q = \emptyset$. Si ottiene che
\begin{gather}
    \{ p \} = X \setminus \bigcup_{q \in X} W_q = \bigcap_{q \in X} (X \setminus W_q) \notag
\end{gather}
inoltre $V_q \subset X \setminus W_q$, pertanto $X \setminus W_q \in \mathcal{N}(p)$ ed è un chiuso per ogni $q \in X \setminus \{ p \}$. Dunque, dato che $X \setminus W_q$ sono una classe particolare di intorni chiusi di $p$, si ottiene
\begin{gather}
    \{ p \} \subset \bigcap_{M \in \mathcal{N}(p)} \overline{M} \subset \bigcap_{q \in X} (X \setminus W_q) \notag
\end{gather}
segue che
\begin{gather}
    \{ p \} = \bigcap_{M \in \mathcal{N}(p)} \overline{M} \notag
\end{gather} \\ \\
$(\leftarrow)$ sia $\{ p \} = \bigcap_{M \in \mathcal{N}(p)} \overline{M}$ per ogni $p \in X$, sia $q \in X \setminus \{ p \}$, si ha
\begin{gather}
    q \in X \setminus \bigcap_{M \in \mathcal{N}(p)} \overline{M} = \bigcup_{M \in \mathcal{N}(p)} (X \setminus \overline{M}) \notag
\end{gather}
ovvero esiste $M^* \in \mathcal{N}(p)$ tale che $q \in X \setminus \overline{M^*}$, che è un aperto, segue che $p \in \overline{M^*}$, dunque esiste un aperto $A^* \subset \overline{M^*}$ tale che $p \in A^*$ e
\begin{gather}
    A^* \cap (X \setminus \overline{M^*}) \subset \overline{M^*} \cap (X \setminus \overline{M^*}) = \emptyset \notag
\end{gather}
pertanto segue la tesi. \\ \\ \\
\textbf{4.} sia $f : X \longrightarrow Y$ continua e $Y$ di Hausdorff, dimostrare che il grafico di $f$, ovvero
\begin{gather}
    \Gamma_f := \{ (x, f(x)) \in X \times Y : x \in X \} \notag
\end{gather}
è chiuso in $X \times Y$. \\

\textbf{Dimostrazione:} sia $(x_0, y_0) \in X \times Y \setminus \Gamma_f$, allora $y_0 \neq f(x_0)$, pertanto esistono $V \in \mathcal{N}(y_0)$ e $W \in \mathcal{N}(f(x_0))$ tali che $V \cap W = \emptyset$, inoltre, per ipotesi di continuità, esiste $M \in \mathcal{N}(x_0)$ aperto tale che $f(M) \subset W$. Dunque si ha $(x_0, y_0) \in M \times V$ aperto in $X \times Y$, inoltre $M \times V \subset X \times Y \setminus \Gamma_f$, infatti se $x_0 \in M$ si ha che $f(x_0) \in f(M)$, ma $V \cap f(M) \subset V \cap W = \emptyset$. Ciò vale per ogni $(x, y) \in X \times Y \setminus \Gamma_f$, pertanto si ottiene la tesi. \\ \\ \\
\textbf{5.} Dimostrare che se $f, g : X \longrightarrow Y$ sono applicazioni continue, $Y$ di Hausdorff e $D \subset X$ è un denso tale che $f(x) = g(x)$ per ogni $x \in D$ si ha che $f = g$. \\

\textbf{Dimostrazione:} sia $x \in X$ tale che $f(x) \neq g(x)$, per ipotesi segue che $x \in X \setminus D$, inoltre esistono $V \in \mathcal{N}(f(x))$ e $W \in \mathcal{N}(g(x))$ aperti tali che $V \cap W = \emptyset$, ovvero $f^{-1}(V), f^{-1}(W) \subset X \setminus D$ e sono aperti (per la continuità di $f$). Tuttavia $D$ è denso, quindi interseca con ogni aperto non vuoto, pertanto $D \cap f^{-1}(V) \neq \emptyset$ e $D \cap f^{-1}(W) \neq \emptyset$, ma ciò è un assurdo, segue che $f = g$. \\ \\ \\
\textbf{6.} Siano $(X, d)$ uno spazio metrico, $x \in X$ e $B \subset X$. Dimostrare che $x$ è un punto di accumulazione di $B$ se e solo se esiste una successione $\{x_j : j \in \mathbf{N}\} \subset B\setminus X$ tale che $x_j \rightarrow x$. \\

\textbf{Dimostrazione:} $(\rightarrow)$ sia $x$ un punto di accumulazione per $B$, allora per definizione si ha che $V \cap B \setminus \{x\} \neq \emptyset$ per ogni $V \in \mathcal{N}(x)$. Consideriamo la seguente famiglia di intorni aperti di $x$
\begin{align}
    \mathcal{V} = \{D_{1/n}(x) : n \in \mathbf{N}\} \notag
\end{align}
Per ogni $n \in \mathbf{N}$ scegliamo un elemento $x_n \in D_{1/n}(x) \cap B$, abbiamo quindi costruito una successione $\{x_n : n \in \mathbf{N}\} \subset B$ tale che
\begin{align}
    d(x_n, x) < \frac{1}{n} \notag
\end{align}
per ogni $n \in \mathbf{N}$, ovvero $x_n \rightarrow x$. \\
$(\leftarrow)$ Sia $\{x_n : n \in \mathbf{N}\} \subset B$ una successione convergente a $x$. Per definizione, per ogni $\varepsilon > 0$ esiste $N_\varepsilon \in \mathbf{N}$ tale che
\begin{align}
    d(x_n, x) < \varepsilon \notag
\end{align}
ovvero per ogni $n > N_\varepsilon$ si ha che $x_n \in D_{\varepsilon}(x)$. Sia $V \in \mathcal{N}(x)$, poiché $V$ è aperto esiste $\delta > 0$ tale che $D_\delta(x) \subset V$, ovvero $x$ è interno a $V$, quindi $x_m \in V \cap B \setminus \{x\}$ per ogni $m > N_\delta$.

\newpage
\subsection{Connessione}
\textbf{Definizione 1.3.1:} sia $X$ uno spazio topologico, $X$ è detto \textit{sconnesso} se esistono $A, B \subset X$ aperti non vuoti tali che $A \cup B = X$ e $A \cap B = \emptyset$. \\ Altrimenti $X$ è detto \textit{connesso}. \\ \\
Notiamo che ciò equivale a dire che $X$ è \textit{connesso} se e solo se gli unici insiemi ad essere sia chiusi che aperti sono $\emptyset$ e $X$. Infatti se esistono $A, B$ aperti disgiunti tali che $X = A \cup B$, si ha che $X \setminus A = B$ è sia aperto, per ipotesi, che chiuso, poiché il complementare è aperto. \\ \\ \\
\textbf{Definizione 1.3.2:} $Y \subset X$ è \textit{sconnesso} se e solo se è sconnesso per la topologia indotta da $X$, ovvero $\tau_Y := \{ Y \cap A : A \text{ aperto in } X \}$. \\ \\
Ciò equivale a dire che $Y \subset X$ è \textit{sconnesso} se e solo se esistono $(Y \cap A), (Y \cap B) \subset\\\subset Y$ aperti non vuoti tali che $Y = Y \cap (A \cup B)$ e $(Y \cap A) \cap (Y \cap B) = \emptyset$. \\ \\
\textbf{Osservazione:} $Y \subset X$ è \textit{connesso} se e solo se per ogni $A, B \subset X$ aperti non vuoti e disgiunti tali che $Y \subseteq A \cup B$ si ha che $Y \subseteq A$ oppure $Y \subseteq B$. \\ \\ \\
\textbf{Proposizione 1.3.1:} siano $Y, Z \subset X$ tali che $Y \subseteq Z \subseteq \overline{Y}$. Allora se $Y$ è connesso anche $Z$ è connesso. \\

\textbf{Dimostrazione:} sia $M$ un sottoinsieme non vuoto, aperto e chiuso di $\overline{Y}$, poiché $\mathring{(\overline{Y} \setminus Y)} = \emptyset$ si ha che $M \cap Y \neq \emptyset$, allora
\begin{gather}
    M \cap Y = Y \notag
\end{gather}
perché $Y$ è connesso per ipotesi, dunque è l'unico insieme non vuoto, aperto e chiuso contenuto in $Y$ stesso. Ciò comporta $M \supseteq Y$, pertanto, dato che $M$ è chiuso, si ha che
\begin{gather}
    M = \overline{Y} \implies M \cap Z = Z \notag
\end{gather}
ovvero l'unico insieme non vuoto, aperto e chiuso contenuto in $Z$ è $Z$ stesso, la tesi segue. \\ \\
Nella dimostrazione si è usato che $\mathring{(\overline{Y} \setminus Y)} = \emptyset$, ciò è di immediata dimostrazione, infatti consideriamo un aperto $A$ contenuto in $\mathring{(\overline{Y} \setminus Y)}$, si ha che $Y \subset \overline{Y} \setminus A$, tuttavia $\overline{Y} \setminus A$ è un chiuso in topologia relativa (rispetto a $\overline{Y})$, pertanto deve valere $\overline{Y} \setminus A = \overline{Y}$, cioè $A = \emptyset$. \qed \\ \\ \\
\textbf{Proposizione 1.3.2:} siano $Y, Z \subset X$ connessi tali che $Y \cap Z \neq \emptyset$, allora $Y \cup Z$ è connesso. \\

\textbf{Dimostrazione:} supponiamo che $A \neq \emptyset$ sia un sottoinsieme aperto e chiuso di $Y \cup Z$, per fissare le idee poniamo $A \cap Y \neq \emptyset$. $A \cap Y$ è un insieme aperto e chiuso di $Y$, poiché $Y$ è connesso si ha che $Y \cap A = Y$, allora anche $A \cap Z \neq \emptyset$ ed essendo anche $Z$ connesso si ha che $Z \cap A = Z$, pertanto concludiamo che $A = Y \cup Z$, dunque gli unici insiemi aperti e chiusi di $Y \cup Z$ sono $\emptyset$ e $Y \cup Z$ stesso, dunque segue la tesi. \qed \\ \\ \\
\textbf{Definizione 1.3.3:} due punti di uno spazio topologico $X$ sono detti \textit{connessi} se esiste un sottoinsieme connesso di $X$ che li contiene entrambi. \\ \\ \\
\textbf{Proposizione 1.3.3:} se comunque scelti due punti $p, q$ di uno spazio topologico $X$ essi sono connessi, allora $X$ è connesso. \\

\textbf{Dimostrazione:} supponiamo che esistono $A, B \in X$ aperti non vuoti e disgiunti tali che $X = A \cup B$. Consideriamo due punti $p \in A$ e $q \in B$, per ipotesi si ha che esiste un sottoinsieme $K$ connesso tale $p, q \in K$, inoltre si ha che $K = K \cap (A \cup B)$ e $(K \cap A) \cap (K \cap B) = \emptyset$, ma ciò è un assurdo perché $K$ è connesso per ipotesi, pertanto si ottiene la tesi. \qed \\ \\ \\
\textbf{Proposizione 1.3.4:} siano $X$ uno spazio topologico, $Y \subseteq X$ e $\{ Y_\alpha \}_{\alpha \in I}$ una famiglia di sottospazi connessi di $Y$ tali che per ogni $\alpha, \beta \in I$ si ha $Y_\alpha \cap Y_\beta \neq \emptyset$, allora $\bigcup_{\alpha \in I} Y_\alpha$ è connesso. \\

\textbf{Dimostrazione:} siano $p_1, p_2 \in \bigcup_\alpha Y_\alpha$, allora esistono $\alpha_1, \alpha_2 \in I$ tali che $p_1 \in Y_{\alpha_1}$ e $p_2 \in Y_{\alpha_2}$, inoltre per ipotesi si ha $Y_{\alpha_1} \cap Y_{\alpha_2} \neq \emptyset$, pertanto $Y_{\alpha_1} \cup Y_{\alpha_2}$ è connesso per la \textit{proposizione 1.3.1}. Pertanto, poiché ciò vale per ogni coppia di punti $p_1, p_2 \in \bigcup_\alpha Y_\alpha$, segue dalla \textit{proposizione 1.3.2} che $\bigcup_\alpha Y_\alpha$ è connesso. \\ \\
Notiamo che se esiste $\gamma \in I$ tale che $p_1, p_2 \in Y_\gamma$ il risultato è ancora valido poiché $Y_\gamma$ è connesso per ipotesi. \qed \\ \\ \\
\textbf{Teorema 1.3.1:} gli intervalli di $\mathbf{R}$, in particolare $\mathbf{R}$ stesso, sono connessi. \\

\textbf{Dimostrazione:} supponiamo che $\mathbf{J}$ sia un intervallo sconnesso, allora esistono $A, B \subset \mathbf{J}$ non vuoti, aperti e chiusi tali che $\mathbf{J} = A \cup B$ e tali che $A \cap B = \emptyset$, siano $a \in A$ e $b \in B$, supponiamo $a < b$. Poiché $\mathbf{J}$ è un intervallo si ha che $[a, b] \subset \mathbf{J}$, inoltre $A \cap [a, b]$ è un chiuso, pertanto contiene il suo estremo superiore, che chiameremo $z$. Inoltre dato che $z = \sup (A \cap [a, b])$ si ha che $(z, b) \subset B$, ma $B$ anche è chiuso, pertanto $z \in B$, ciò è un assurdo. \qed \\ \\ \\
\textbf{Proposizione 1.3.5:} se $f : X \longrightarrow Y$ un'applicazione continua e $Z$ è un sottoinsieme connesso di $X$, $f(Z)$ è un sottoinsieme connesso di $Y$. \\

\textbf{Dimostrazione:} sia $M$ un sottoinsieme non vuoto, aperto e chiuso di $f(Z)$. Per ipotesi di continuità si ha che $f^{-1}(M)$ è un sottoinsieme non vuoto, aperto e chiuso di $X$ tale che $f^{-1}(M) \cap Z = Z$, perché $Z$ è connesso. Dunque $f^{-1}(M) \supseteq Z$ e ciò comporta che $M = f(f^{-1}(M)) = f(Z)$, la tesi segue. \qed \\ \\
Notiamo che la \textit{connessione} è una \textit{proprietà topologica}, infatti se $X \cong Y$ esiste un omeomorfismo $f : X \longrightarrow Y$, per definizione $f$ è continuo e anche $f^{-1}$ lo è, dunque se $X$ è connesso anche $Y = f(X)$ e viceversa se $Y$ è connesso anche $X = f^{-1}(X)$ lo è (ricordiamo $f$ omeomorfismo se è continuo, biettivo e con inversa continua). \\ \\ \\
In uno spazio topologico $X$ possiamo introdurre una relazione di equivalenza binaria $\rho$ ponendo
\begin{gather}
    p \rho q \text{ se e solo se } p,q \text{ sono connessi in } X \notag
\end{gather}
Verifichiamo che si tratta effettivamente di una relazione di equivalenza: \\
\textbf{1.} $p \rho q \iff q \rho p$, segue dalla \textit{prop. 1.3.3}. \\
\textbf{2.} $p \rho q$ e $q \rho r$ implica $p \rho r$, segue dalle \textit{prop. 1.3.3} e \textit{prop. 1.3.2}, infatti se $p, q$ sono connessi esiste $M$ connesso tale che $p, q \in M$, se $q, r$ sono connessi esiste $N$ connesso tale che $q, r \in N$, pertanto anche $p, r$ sono connessi poiché $M \cup N$ è un connesso (\textit{prop. 1.3.2}) tale che $p, r \in M \cup N$. \\ \\
Le classi di $\rho$\textit{-equivalenza} sono dette componenti connesse di X. Per ogni $p \in X$ denotiamo con $C(p)$ la \textit{componente connessa di} $p$. \\ \\
\textbf{Osservazione:} dalla \textit{proposizione 1.3.3} segue che $X$ è connesso se e solo se possiede un'unica componente connessa. \\ \\ \\
\textbf{Proposizione 1.3.6:} per ogni $p \in X$, $C(p)$ è il più grande sottoinsieme connesso di $X$ contenente $p$. \\

\textbf{Dimostrazione:} $C(p)$ è connesso poiché unione di insiemi connessi. Supponiamo ora che esiste $Z \supset C(p)$ connesso, allora esiste $q \in Z$ tale che $q \not\in C(p)$ e $p, q$ sono connessi, ciò comporterebbe, per definizione, $q \in C(p)$, ma ciò è un assurdo. \\ \\ Notiamo quindi che $C(p)$ è un chiuso per la \textit{proposizione 1.3.1}. \qed \\ \\ \\
\textbf{Teorema 1.3.2:} il prodotto topologico di una qualsiasi famiglia di spazi connessi è connesso. \\

\textbf{Dimostrazione:} iniziamo con il prodotto topologico di due insiemi connessi $X, Y$. Consideriamo i punti $(x_1, y_1), (x_2, y_2) \in X \times Y$, gli insiemi $X \times \{ y_2 \}$ e $\{ x_1 \} \times Y$ sono omeomorfi a $X$ e $Y$ rispettivamente, infatti $f_X$ e $f_Y$ definite ponendo
\begin{gather}
    f_X(x) = (x, y_2) \qquad f_Y(y) = (x_1, y) \notag
\end{gather}
sono continue, biettive e con inversa continua. Pertanto $X \times \{ y_2 \}$ e $\{ x_1 \} \times Y$ sono connessi (la connessione è una proprietà topologica), quindi in $X \times Y$ vale
\begin{gather}
    (x_1, y_1) \rho (x_1, y_2) \notag \\
    (x_1, y_1) \rho (x_2, y_1) \notag
\end{gather}
e per la transitività si ha
\begin{gather}
    (x_1, y_2) \rho (x_2, y_1) \notag
\end{gather}
Concludiamo che $X \times Y$ è connesso per la \textit{proposizione 1.3.3}. Ragionando per induzione si ottiene che il prodotto topologico di un numero finito di insiemi connessi è connesso. \\ \\
Dimostriamo ora il caso generale. Sia $\{ X_\mu : \mu \in M \}$ una famiglia di insiemi connessi e sia 
\begin{gather}
    X := \prod_{\mu \in M} X_\mu \notag
\end{gather}
vogliamo definire la topologia prodotto in modo tale che renda continue le \textit{proiezioni canoniche} $\pi_j : \prod_{\mu \in M} X_\mu \longrightarrow X_j$. Pertanto imponiamo che per ogni $j \in M$ si ha
\begin{gather}
    \pi_j^{-1}(A_j) = \{ x \in X : x(j) \in A_j \} = \prod_{\mu \neq j} X_\mu \times A_j \notag
\end{gather}
sia un aperto per ogni aperto $A_j \subset X_j$. Sia $\mu_1, \dots, \mu_k$ un numero finito di indici, segue da quanto detto prima che anche
\begin{gather}
    \bigcap_{n = 1}^k \pi_{\mu_n}^{-1}(A_{\mu_n}) = \{ x \in X : x(\mu) \in A_{\mu}, \mu = \mu_1, \dots, \mu_k \} = \prod_{\mu \in M} V_\mu \notag
\end{gather}
con $V_\mu = A_\mu$ per $\mu = \mu_1, \dots, \mu_k$ e $V_\mu = X_\mu$ altrimenti, è un aperto. Questo procedimento ci ha permesso di avere una topologia prodotto che ha le stesse proprietà della topologia prodotto per il prodotto cartesiano di un numero finito di spazi topologici. \\
Consideriamo un punto $p \in X$ e un aperto arbitrario $B$ della base della topologia prodotto, per quanto visto prima esiste un numero finito di indici $\mu_1, \dots, \mu_s \in M$ e di aperti non vuoti $A_1 \in X_{\mu_1}, \dots, A_s \in X_{\mu_s}$ tali che
\begin{gather}
    B = \{ x \in X : x(\mu_i) \in A_i, i = 1, \dots, s \} \notag
\end{gather}
Fissiamo $a_1 \in A_1, \dots, a_s \in A_s$; denotiamo il punto $q$ di $B$ ponendo
\begin{gather}
    q(\mu) = \left\{ \begin{array}{cc} 
         a_i& \text{ se } \mu = \mu_i \text{ per qualche } i = 1, \dots, s  \\
         p(\mu)& \text{ altrimenti} \notag
    \end{array} \right.
\end{gather}
Consideriamo poi il sottoinsieme di $X$
\begin{gather}
    G := \{ x \in X : x(\mu) = p(\mu), \mu \neq \mu_1, \dots, \mu_s \} \notag
\end{gather}
$G$ è connesso poiché omeomorfo a $X_{\mu_1} \times \dots \times X_{\mu_s}$, infatti
\begin{gather}
    G = X_{\mu_1} \times \dots \times X_{\mu_s} \times \prod_{\mu \neq \mu_i} p(\mu), \quad i = 1, \dots, s \notag
\end{gather}
inoltre si ha $q, p \in G$, pertanto $G \subset C(p)$ e $C(p) \cap B \neq \emptyset$, dunque, essendo $B$ un aperto qualsiasi, si ha che $C(p)$ è denso ed è chiuso per la \textit{proposizione 1.3.6}, otteniamo allora che $C(p) = X$, cioè $X$ è connesso. \qed
\subsubsection{Esercizi}
\textbf{1.} Sia $(X, \tau)$ uno spazio topologico, sia $\tau' \subseteq \tau$. Dimostrare che se $(X, \tau)$ è connesso anche $(X, \tau')$ è connesso. \\

\textbf{Dimostrazione:} sia $(X, \tau)$ connesso, per ipotesi non esistono due aperti $A, B \in \tau$ disgiunti e non vuoti tali che $X = A \cup B$, allora in particolare non esistono neanche in $\tau'$, proprio perché ogni aperto di $\tau'$ è anche un aperto di $\tau$. \\ \\ \\
\textbf{2.} Dimostrare se è sempre vero che uno spazio topologico $X$ è connesso se e solo se $\partial(Y) \neq 0$ per ogni sottoinsieme $Y$, $Y \neq \emptyset, X$. \\

\textbf{Dimostrazione:} $(\rightarrow)$ Sia $X$ connesso, ciò comporta che gli unici sottoinsiemi chiusi e aperti di $X$ siano $X$ e $\emptyset$, pertanto segue la tesi. \\ \\
$(\leftarrow)$ Supponiamo che esista un sottoinsieme $Y$ di $X$, diverso da $X$ stesso e da $\emptyset$, tale che $\partial(Y) = \emptyset$. Allora, poiché $\partial(Y) = \overline{Y} \setminus \mathring{Y}$ si ha che $\overline{Y} = \mathring{Y}$, ovvero $Y$ è sia chiuso che aperto, segue che $X$ non è connesso. \\ \\ \\
\textbf{3.} Sia $X$ uno spazio topologico $T_2$ e sia $E \subset X$ connesso. Dimostrare se è sempre vero che se $|E| \geq 2$ allora $|E| = +\infty$. \\

\textbf{Dimostrazione:} supponiamo $|E| = 2$, siano $p, q$ gli elementi di $E$. Per ipotesi esistono $V, W$ intorni aperti tali che
\begin{gather}
    p \in V, \quad q \in W, \quad V \cap W = \emptyset \notag
\end{gather}
Dunque $V \cap E$ e $W \cap E$ sono due aperti disgiunti non vuoti tali che
\begin{gather}
    E = (V \cap E) \cup (W \cap E) \notag
\end{gather}
pertanto $E$ non è connesso. \\
Ragionando per induzione si può generalizzare questo risultato ad un numero finito $n$ di elementi. Infatti, supponiamo $E = \{ x_i : i = 1, \dots, n \}$ e prendiamo $x_1 \in E$. Per ipotesi, per ogni $i = 1, \dots, n$ esistono $V_i, W_i$ intorni aperti tali che
\begin{gather}
    x_1 \in V_i, \quad x_i \in W_i, \quad V_i \cap W_i = \emptyset \notag
\end{gather}
Allora si ottiene che
\begin{gather}
    x_1 \in V := \bigcap_{i = 1}^n V_i, \quad E \setminus \{ x_1 \} \subset W := \bigcup_{i = 1}^n W_i \notag
\end{gather}
inoltre
\begin{gather}
    V \cap W = \emptyset \notag
\end{gather}
ovvero $V, W$ sono due aperti disgiunti tali che
\begin{gather}
    E = (V \cap E) \cup (W \cap E) \notag
\end{gather}
dunque $V \cap E, W \cap E$ sono sottoinsiemi non vuoti aperti e chiusi di $E$ distinti da $E$ stesso, segue che $E$ non è connesso e quindi la tesi.

\newpage
\subsection{Connessione per archi}
Un \textit{arco} (o \textit{cammino}) in uno spazio topologico $X$ è un'applicazione continua $\gamma : \mathbf{I} \longrightarrow X$, con $\mathbf{I} = [0, 1]$. I punti $\gamma(0)$ e $\gamma(1)$ si dicono \textit{estremi} di $\gamma$, e rispettivamente \textit{punto iniziale} e \textit{punto finale} di $\gamma$. \\ \\
\textbf{Definizione 1.4.1:} se per ogni $p, q \in X$ esiste un arco $\gamma$ tale che $\gamma(0) = p$ e $\gamma(1) = q$, $X$ si dice \textit{connesso per archi}. \\ \\
La connessione per archi è una proprietà più forte della connessione. \\ \\ \\
\textbf{Proposizione 1.4.1:} se $X$ è connesso per archi si ha che $X$ è connesso. \\

\textbf{Dimostrazione:} siano $p, q \in X$, per ipotesi esiste $\gamma : \mathbf{I} \longrightarrow X$ continua tale che $\gamma(0) = p$ e $\gamma(1) = q$. Consideriamo l'insieme $\gamma(\mathbf{I})$, per ipotesi di continuità di $\gamma$ si ha che $\gamma(\mathbf{I})$ è connesso (gli intervalli sono connessi), inoltre $p, q \in \gamma(\mathbf{I})$. Concludiamo che per ogni coppia di punti di $X$ esiste un sottoinsieme connesso che li contiene, segue dalla \textit{proposizione 1.3.3} che $X$ è connesso. \qed \\ \\
Il viceversa non è vero in generale, esistono insiemi connessi, ma non connessi per archi. Inoltre segue che un insieme sconnesso non è connesso per archi. \\ \\ \\
\textbf{Proposizione 1.4.2:} sia $f : X \longrightarrow Y$ un'applicazione continua, sia $Z$ un sottoinsieme di $X$ connesso per archi, allora $f(Z)$ è connesso per archi. \\ 

\textbf{Dimostrazione:} per ipotesi per ogni coppia di punti $p, q \in Z$ esiste $\gamma : \mathbf{I} \longrightarrow Z$ continua tale che $\gamma(0) = p$ e $\gamma(1) = q$. Consideriamo la restrizione di $f$ a $Z$, ovvero $f_Z : Z \longrightarrow f(Z)$ e la funzione composta $f_Z \cdot \gamma : \mathbf{I} \longrightarrow f(Z)$, per ipotesi di continuità di $\gamma$ e $f_Z$ anche $f_Z \cdot \gamma$ è continua, inoltre $f_Z(\gamma(0)) = f_Z(p) = f(p)$ e $f_Z(\gamma(1)) = f_Z(q) = f(q)$. Pertanto per ogni coppia di punti $f(p), f(q) \in f(Z)$ esiste una funzione continua $f_Z \cdot \gamma : \mathbf{I} \longrightarrow f(Z)$ tale che $f_Z \cdot \gamma (0) = f(p)$ e $f_Z \cdot \gamma (1) = f(q)$, dunque $f(Z)$ è connesso per archi. \qed \\ \\
Un'analogia tra connessione per archi e connessione si ha anche per il concetto di \textit{componente connessa per archi}. \\
Sia $X$ uno spazio topologico, introduciamo una relazione binaria $\varepsilon$ ponendo $x \varepsilon y$ se e solo se esiste un arco $\gamma$ in $X$ tale che $\gamma(0) = x$ e $\gamma (1) = y$. \\
Verifichiamo che si tratta di una relazione di equivalenza. $x \varepsilon x$ per ogni $x \in X$ segue dal fatto che l'applicazione costante $\gamma_x(t) = x$ per ogni $t \in \mathbf{I}$ è un arco. Se $x \varepsilon y$ si ha $y \varepsilon x$, basta considerare l'arco $\alpha = \beta (1 - t)$, dove $\beta$ è un arco da $x$ a $y$. Infine se $x \varepsilon y$ e $y \varepsilon z$ si ha $x \varepsilon z$, infatti per ipotesi esistono $\alpha, \beta$ archi tali che $\alpha(0) = x$, $\alpha(1) = y$, $\beta(0) = y$ e $\beta(1) = z$, pertanto possiamo definire $c : \mathbf{I} \longrightarrow X$ ponendo
\begin{gather}
    c(t) = \left\{ \begin{array}{cc}
         a(2t)& \text{ se } 0 \leq t \leq \frac{1}{2}  \\
         b(2t - 1)& \text{ se } \frac{1}{2} \leq t \leq 1 \notag
    \end{array} \right.
\end{gather}
dunque si ha che $c$ è continua e $c(0) = x$ e $c(1) = z$.\\ \\ \\
\textbf{Proposizione 1.4.3:} per ogni $p \in X$ si ha che $C_a(p)$ è il più grande sottoinsieme connesso per archi di $X$ contente $p$. \\

\textbf{Dimostrazione:} se $Y$ è un sottoinsieme connesso per archi di $X$ contenente $p$ e $q \in Y$, per definizione esiste $\gamma$ arco tale che $\gamma(0) = p$ e $\gamma(1) = q$, dunque $q \in C_a(p)$ e $Y \subset C_a(p)$. \\
Inoltre $C_a(p)$ è connesso per archi, infatti se $q \in C_a(p)$ esiste un arco $\alpha$ in $X$ tale che $\alpha(0) = p$ e $\alpha(1) = q$. Basta dimostrare che $\alpha(\mathbf{I}) \subset C_a(p)$. Sia $s \in \mathbf{I}$, definiamo $\beta_s : \mathbf{I} \longrightarrow X$ ponendo
\begin{gather}
    \beta_s(t) = \alpha(ts), \quad t \in \mathbf{I} \notag
\end{gather}
$\beta_s$ è un arco in $X$ di estremi $p$ e $\alpha(s)$, quindi $\alpha(s) \in C_a(p)$ per ogni $s \in \mathbf{I}$, pertanto $\alpha(\mathbf{I}) \subset C_a(p)$. \qed \\ \\ \\
\textbf{Teorema 1.4.1:} il prodotto topologico di una famiglia qualsiasi di spazi connessi per archi è connesso per archi. \\

\textbf{Dimostrazione:} sia $\{ X_\mu \}_{\mu \in M}$ una famiglia di spazi connessi per archi, siano $p, q \in X := \prod_{\mu \in M} X_\mu$. Per ipotesi di connessione di archi, per ogni $\mu \in M$ esiste un arco $\gamma_\mu$ tale che $\gamma_\mu(0) = p(\mu)$ e $\gamma_\mu(1) = q(\mu)$. Definiamo allora $\alpha : \mathbf{I} \longrightarrow X$ ponendo
\begin{gather}
    \alpha(t)(\mu) = \alpha_\mu(t) \notag
\end{gather}
$\alpha$ è un arco in $X$ tale che $\alpha(0) = p$ e $\alpha(1) = q$, quindi $X$ è connesso per archi. \qed
\subsection{Compattezza}
\textbf{Definizione 1.5.1:} uno spazio topologico $X$ è detto compatto se per ogni \textit{ricoprimento aperto} $\{ \mathcal{U}_i \}_{i \in I}$ esiste un \textit{sottoricoprimento finito} $U_1, \dots, U_n$ tale che $X = U_1 \cup \dots \cup U_n$. \\
Un sottospazio di $X$ si dice compatto se è compatto con la topologia relativa; ciò equivale a dire che un sottoinsieme $K$ di $X$ è compatto se ogni ricoprimento di $K$ con aperti di $X$ possiede un sottoricoprimento finito. \\
Ogni spazio topologico banale è compatto, infatti l'unico ricoprimento aperto è costituito da $X$ stesso. Inoltre è compatto ogni spazio finito e ogni spazio con la topologia cofinita. \\
$\mathbf{R}$ non è compatto, infatti la famiglia $\{ (-n, n) \}_{n \geq 1}$ non ammette sottoricoprimenti finiti. \\
Ogni intervallo aperto $(a, b)$ di $\mathbf{R}$ non è compatto, infatti la famiglia di intervalli
\begin{gather}
    \mathcal{U} := \left\{ (a + \varepsilon, b - \varepsilon) : 0 < \varepsilon < \frac{b - a}{2} \right\} \notag
\end{gather}
è un ricoprimento aperto di $(a, b)$, tuttavia ogni sottoricoprimento finito di $\mathcal{U}$ ha come unione un intervallo ancora appartenente ad $\mathcal{U}$, pertanto diverso da $(a, b)$. \\ \\ \\
\textbf{Teorema (di Heine-Borel):} ogni intervallo chiuso e limitato $[a, b]$ di $\mathbf{R}$ è compatto. \\

\textbf{Dimostrazione:} sia $\mathcal{U}$ un ricoprimento di $[a, b]$ con aperti di $\mathbf{R}$ e sia
\begin{gather}
    X := \{ x \in X : [a, x] \text{ è ricoperto da un numero finito di aperti} \} \notag
\end{gather}
basta dimostrare che $X = [a, b]$. Innanzitutto si ha che $X \neq \emptyset$ perché $a \in X$, infatti esiste $U \in \mathcal{U}$ tale che $a \in U$. \\
Notiamo inoltre che $X$ è un intervallo contenuto in $[a, b]$, infatti se $x \in X$, esistono finiti $U_1, \dots, U_n \in \mathcal{U}$ tali che $[a, x] \subset U_1 \cup \dots U_n$, pertanto si ha che $[a, t] \subset U_1 \cup \dots \cup U_n$ per ogni $a < t < x$. \\
$X$ è un aperto in $[a, b]$. Infatti, sia $x \in X$, per ipotesi esistono $U_1, \dots, U_n \in \mathcal{U}$ tali che $[a, x] \subset U_1 \cup \dots \cup U_n$, inoltre, poiché l'unione di una qualsiasi famiglia di aperti è aperta, si ha che esiste $r > 0$ tale che $(x - r, x+ r) \subset U_1 \cup \dots \cup U_n$, in particolare anche $[a, b] \cap (x - r, x + r) \subset U_1 \cup \dots \cup U_n$, ma allora
\begin{gather}
    [a, b] \cap (x - r, x + r) \subset X \notag
\end{gather}
per come è definito $X$, dunque per ogni $x \in X$ esiste un intorno di $x$ (nella topologia relativa) contenuto in $X$ stesso. \\
Sia $z := \sup X$, poiché $X \subset [a, b]$ che è chiuso e limitato, si ha $z \in [a, b]$, pertanto esiste $V \in \mathcal{U}$ tale che $z \in V \cap [a, b]$. Sia $s > 0$ tale che $(z - s, z] \subset V \cap [a, b]$, allora $z - s \in X$ (poiché $z$ è l'estremo superiore), pertanto esistono $V_1, \dots, V_m \in \mathcal{U}$ tali che $[a, z - s] \subset V_1 \cup \dots \cup V_m$, ma allora si ha
\begin{gather}
    [a, z] \subset V_1 \cup \dots \cup V_m \cup V \rightarrow z \in X \notag
\end{gather}
dunque $X$ è sia chiuso che aperto in $[a, b]$, tuttavia l'unico intervallo chiuso e aperto in $[a, b]$ è $[a, b]$ stesso, perciò $X = [a, b]$, la tesi segue. \qed \\ \\ \\
\textbf{Proposizione 1.5.1:} $a)$ Ogni sottoinsieme chiuso $K$ di uno spazio compatto $X$ è compatto. \\
$b)$ Ogni sottoinsieme compatto $K$ di uno spazio di Hausdorff $X$ è chiuso. \\

\textbf{Dimostrazione:} $a)$ Sia $K \subset X$ chiuso e $X$ compatto. Si ha che $X \setminus K$ è un aperto. Sia $\mathcal{A}$ un ricoprimento aperto di $K$, allora $\mathcal{A} \cup \{X \setminus K\}$ è un ricoprimento aperto di $X$, infatti
\begin{gather}
    \bigcup_{A \in \mathcal{A}} A \supset K \rightarrow (X \setminus K) \cup \bigcup_{A \in \mathcal{A}} A = X \notag
\end{gather}
Pertanto esiste un sottoricoprimento finito $\{A_i : 1 \leq i \leq n \} \cup \{ X \setminus K \}$ tale che
\begin{gather}
    X = A_1 \cup \dots \cup A_n \cup X \setminus K \notag
\end{gather}
ma allora
\begin{gather}
    K \subseteq A_1 \cup \dots \cup A_n \notag
\end{gather}
ovvero $K$ è compatto. \\ \\
$b)$ Sia $K \subset X$ compatto e $X$ di Hausdorff. Sia $p \in X \setminus K$, per ogni $q \in K$ esistono $V_q \in \mathcal{N}(p)$ e $W_q \in \mathcal{N}(q)$ aperti tali che $V_q \cap W_q = \emptyset$. Allora
\begin{gather}
    p \in X \setminus \bigcup_{q \in K} W_q \notag
\end{gather}
Inoltre $\{ W_q : q \in K \}$ è un ricoprimento aperto di $K$, pertanto esiste un sottoricoprimento aperto finito $\{ W_{q_\mu} : 1 \leq \mu \leq n \}$ tale che
\begin{gather}
    K \subseteq W_{q_1} \cup \dots \cup W_{q_n} \notag
\end{gather}
allora si ha che
\begin{gather}
    p \in X \setminus \bigcup_{\mu = 1}^n W_{q_\mu} \subset X \setminus K \notag
\end{gather}
e
\begin{gather}
    V := \bigcap_{\mu = 1}^n V_{q_\mu} \notag
\end{gather}
è un intorno aperto di $p$ contenuto in $X \setminus K$. Poiché ciò vale per ogni $p \in X \setminus K$ si ha che $X \setminus K$ è aperto, quindi $K$ è chiuso. \qed \\ \\ \\
\textbf{Corollario 1.5.1:} un sottoinsieme $K$ di $\mathbf{R}$ è compatto se e solo se è chiuso e limitato. \\

\textbf{Dimostrazione:} $(\rightarrow)$ Sia $K \subset \mathbf{R}$ compatto e sia $X$ di Hausdorff. Per la \textit{proposizione 1.5.1} si ha che $K$ è chiuso. Inoltre è sicuramente limitato, altrimenti $\{ (-n, n) \}_{n \geq 1}$ sarebbe un ricoprimento aperto che non ammette sottoricoprimento finito. \\ \\
$(\leftarrow)$ Sia $K$ chiuso e limitato. Poniamo
\begin{gather}
    a := \inf K, \quad b := \sup K \notag
\end{gather}
poiché $K$ è chiuso si ha $a, b \in K$, allora $K \subseteq [a, b]$ che è compatto perché è un intervallo chiuso di $\mathbf{R}$, inoltre $K$ è chiuso in $[a, b]$, pertanto $K$ è un sottoinsieme chiuso di uno spazio compatto, per la \textit{proposizione 1.5.1} $K$ è compatto. \qed \\ \\ \\
\textbf{Proposizione 1.5.2:} $a)$ Ogni sottoinsieme infinito $Z$ di uno spazio compatto $X$ possiede almeno un punto di accumulazione in $X$. \\
$b)$ Ogni successione di punti di uno spazio compatto e $T_1$ $X$ che soddisfa il primo assioma di numerabilità ammette una sottosuccessione convergente. \\

\textbf{Dimostrazione:} $a)$ Supponiamo che $Z$ non abbia punti di accumulazione, $Z$ è chiuso in $X$ e dunque compatto. Inoltre la topologia relativa di $Z$ è la topologia discreta. Tuttavia uno spazio infinito con la topologia discreta NON può essere compatto, perciò si ottiene un assurdo. Concludiamo che $Z$ ha almeno un punto di accumulazione in $X$. \\ \\
$b)$ Se la successione ha infiniti numeri uguali, essi formano una sottosuccessione convergente. Se la successione ha infiniti punti distinti, allora possiede un punto di accumulazione in $X$ e, per la \textit{proposizione 1.2.7}, ammette una sottosuccessione convergente. \qed \\ \\ \\
\textbf{Teorema (di Bolzano - Weierstrass):} ogni successione limitata $\{ x_n : n \in \mathbf{N} \}$ di numeri reali ammette una sottosuccessione convergente. \\

\textbf{Dimostrazione:} poiché $\{ x_n \}$ è limitata esiste $M > 0$ tale che $\{ x_n \} \subset [-M, M]$, ma $[-M, M]$ è un sottoinsieme compatto di $\mathbf{R}$, pertanto, per la \textit{proposizione 1.5.2}, $\{ x_n \}$ ammette una sottosuccessione convergente. \qed \\ \\ \\
\textbf{Proposizione 1.5.3:} sia $f : X \longrightarrow Y$ un'applicazione continua di spazio topologici e sia $X$ compatto, allora $f(X)$ è compatto in $Y$. \\

\textbf{Dimostrazione:} sia $\mathcal{U}$ un ricoprimento aperto di $f(X)$. La famiglia di sottoinsiemi aperti
\begin{gather}
    \mathcal{V} := \{ f^{-1}(U) : U \in \mathcal{U} \} \notag
\end{gather}
è un ricoprimento aperto di $X$, pertanto ammette un sottoricoprimento finito
\begin{gather}
    \{ f^{-1}(U_i) : U_i \in \mathcal{U}, 1 \leq i \leq n \} \notag
\end{gather}
tale che
\begin{gather}
    X = f^{-1}(U_1) \cup \dots \cup f^{-1}(U_n) = f^{-1}(U_1 \cup \dots \cup U_n) \notag
\end{gather}
Dunque si ha che
\begin{gather}
    f(X) = f(f^{-1}(U_1 \cup \dots \cup U_n)) \subset U_1 \cup \dots \cup U_n \notag
\end{gather}
(contenuto poiché $f$ potrebbe non essere suriettiva), ovvero $\{ U_1, \dots, U_n\}$ è un ricoprimento aperto finito di $f(X)$. $f(X)$ è compatto. \qed \\ \\
Dalla \textit{proposizione 1.5.3} segue immediatamente che la compattezza è una proprietà topologica, dunque invariante per omeomorfismi. \\ \\ \\
\textbf{Corollario 1.5.2:} sia $X$ uno spazio topologico e $f : X \longrightarrow \mathbf{R}$ continua. Se $X$ è compatto $f$ assume un valore massimo ed un valore minimo su $\mathbf{R}$. \\

\textbf{Dimostrazione:} dalla \textit{proposizione 1.5.3} segue che $f(X)$ è compatto, pertanto per il \textit{corollario 1.5.1} è chiuso e limitato, dunque $f(X)$ contiene il suo estremo superiore ed il suo estremo inferiore. \qed \\ \\ \\
\textbf{Proposizione 1.5.4:} siano $X$ compatto, $Y$ di Hausdorff, $f : X \longrightarrow Y$ continua. Allora $f$ è chiusa; inoltre se $f$ è biunivoca, $f$ è un omeomorfismo. \\

\textbf{Dimostrazione:} sia $C \subset X$ chiuso, allora è compatto (sottoinsieme chiuso di spazio compatto), per ipotesi di continuità $f(C) \subset Y$ è compatto e dunque chiuso (sottoinsieme compatto di spazio di Hausdorff), pertanto $f$ è chiusa. Se $f$ è anche biunivoca segue immediatamente che $f$ è omeomorfismo, poiché $f$ è continua, biunivoca e chiusa. \qed \\ \\ \\
\textbf{Proposizione 1.5.5:} condizione \textit{necessaria} e \textit{sufficiente} affinché uno spazio $X$ sia compatto è che ogni famiglia $\mathcal{F}$ di sottoinsiemi chiusi di $X$ che ha la proprietà dell'intersezione finita, ovvero $F_1 \cap \dots \cap F_n \neq \emptyset$ per ogni $F_1, \dots, F_n \in \mathcal{F}$, abbia intersezione non vuota, ovvero $\bigcap_{F \in \mathcal{F}} F \neq \emptyset$. \\

\textbf{Dimostrazione:} \textit{(necessità)} Supponiamo che esista una famiglia $\mathcal{F}$ di sottoinsiemi chiusi di $X$ con la proprietà dell'intersezione finita che ha intersezione nulla e dimostriamo che $X$ non può essere compatto. Consideriamo la famiglia di aperti
\begin{gather}
    \mathcal{A} := \{ X \setminus F : F \in \mathcal{F} \} \notag
\end{gather}
$\mathcal{A}$ è un ricoprimento aperto, infatti
\begin{gather}
    \bigcup_{F \in \mathcal{F}} X \setminus F = X \setminus \bigcap_{F \in \mathcal{F}} F = X \setminus \emptyset = X \notag
\end{gather}
Tuttavia $\mathcal{A}$ non ammette sottoricoprimento finito, infatti per ogni famiglia finita $\{ F_1, \dots, F_n \} \subset \mathcal{F}$ si ha
\begin{gather}
    \bigcup_{i = 1}^n X \setminus F_i = X \setminus \bigcap_{i = 1}^n F_i \neq X \notag 
\end{gather}
infatti $F_1 \cap \dots \cap F_n \neq \emptyset$ per ipotesi, pertanto $X$ non è compatto. \\ \\
\textit{(sufficienza)} Supponiamo che ogni famiglia $\mathcal{F}$ di sottoinsiemi chiusi di $X$ con la proprietà dell'intersezione finita abbia intersezione non nulla. Sia $\mathcal{A}$ un ricoprimento aperto di $X$, consideriamo la famiglia
\begin{gather}
    \mathcal{F} := \{ X \setminus A : A \in \mathcal{A} \} \notag
\end{gather}
di sottoinsiemi chiusi di $X$. Poiché $\mathcal{A}$ è un ricoprimento di $X$, si ha che
\begin{gather}
    \bigcap_{A \in \mathcal{A}} X \setminus A = X \setminus \bigcup_{A \in \mathcal{A}} A = \emptyset \notag
\end{gather}
ovvero $\mathcal{F}$ ha intersezione nulla, pertanto non ha la proprietà dell'intersezione finita. Segue che esiste una famiglia finita $\{ A_1, \dots, A_m \} \subset \mathcal{A}$ tale che
\begin{gather}
    \bigcap_{i = 1}^m X \setminus A_i = X \setminus \bigcup_{i = 1}^m A_i = \emptyset \notag
\end{gather}
ovvero $A_1 \cup \dots \cup A_m = X$, pertanto $\mathcal{A}$ ammette un sottoricoprimento finito, quindi $X$ è compatto. \qed \\ \\ \\
\textbf{Teorema (di Tichonov):} il prodotto topologico di una qualsiasi famiglia di spazi compatti è compatto. \\

\textbf{Dimostrazione:} partiamo dal prodotto topologico di due spazi compatti $X$ e $Y$. Sia $\{ A_j : j \in J \}$ un ricoprimento aperto di $X \times Y$ (topologia prodotto), allora
\begin{gather}
    A_j = \bigcup_{h \in h(j)} V_h \times W_h \notag
\end{gather}
con $V_h$ aperto di $X$ e $W_h$ aperto di $X$ per ogni $h \in h(j)$. La famiglia
\begin{gather}
    \{ V_h \times W_h : h \in H \}, \quad H = \bigcup_{j \in J} h(j) \notag
\end{gather}
è un ricoprimento aperto di $X \times Y$. Sia $x \in X$, lo spazio $\{ x \} \times Y$ è omeomorfo a $Y$, quindi è compatto, pertanto esiste un sottoinsieme finito $h(x) \subset H$ tale che
\begin{gather}
    \{ x \} \times Y \subset \bigcup_{h \in h(x)} V_h \times W_h \notag
\end{gather}
e tale che $x$ è contenuto in $V_h$ per ogni $h \in h(x)$, infatti se esiste $\overline{h} \in h(x)$ tale che $x \not\in V_{\overline{h}}$, basta rimuoverlo e si ha ancora un ricoprimento di $\{ x \} \times Y$. Per ogni $x \in X$ consideriamo l'intorno aperto
\begin{gather}
    V(x) := \bigcap_{h \in h(x)} V_h \notag
\end{gather}
allora $\{ V(x) : x \in X \}$ è un ricoprimento aperto di $X$, pertanto, per la compattezza di $X$, esistono $x_1, \dots, x_s \in X$ tali che
\begin{gather}
    X = V(x_1) \cup \dots \cup V(x_s) \notag
\end{gather}
allora
\begin{gather}
    \{ V(x_1) \times W_h : h \in h(x_1) \} \cup \dots \cup \{ V(x_s) \times W_h : h \in h(x_s) \} \notag
\end{gather}
è un ricoprimento aperto di $X \times Y$, inoltre per ogni $i = 1 \dots s$ e ogni $h \in h(x_i)$ si ha
\begin{gather}
    V(x_i) \times W_h \subset A_{j(i, h)}, \text{ per qualche } j(i, h) \in J \notag
\end{gather}
pertanto è un sottoricoprimento finito di $\{ A_j \}_{j \in J}$. Ragionando per induzione si ottiene che il prodotto topologico di una famiglia finita di spazi compatti è compatto. \\ \\
Dimostriamo ora il caso generale. Sia $\{ X_\mu : \mu \in M \}$ una famiglia di spazi compatti e sia
\begin{gather}
    X := \prod_{\mu \in M} X_\mu \notag
\end{gather}
Consideriamo una famiglia qualsiasi $\mathcal{B}$ di sottoinsiemi di $X$ con la proprietà dell'intersezione finita e dimostreremo che
\begin{gather}
    \bigcap_{B \in \mathcal{B}} \overline{B} \neq \emptyset \notag
\end{gather}
da ciò seguirà la compattezza dalla \textit{proposizione 1.5.5}. \\
Consideriamo la classe $\mathcal{I}$ di tutte le famiglie di sottoinsiemi di $X$ che hanno la proprietà dell'intersezione finita. $\mathcal{I}$ è \textit{parzialmente ordinata} rispetto all'inclusione, ovvero
\begin{gather}
    \mathcal{B} \subseteq \mathcal{B} \text{ per ogni } \mathcal{B} \in \mathcal{I} \notag \\
    \text{ se } \mathcal{B}_1 \subseteq \mathcal{B}_2 \text{ e } \mathcal{B}_2 \subseteq \mathcal{B}_1 \implies \mathcal{B}_1 = \mathcal{B}_2 \notag \\
    \text{ se } \mathcal{B}_1 \subseteq \mathcal{B}_2 \text{ e } \mathcal{B}_2 \subseteq \mathcal{B}_3 \implies \mathcal{B}_1 \subseteq \mathcal{B}_3 \notag
\end{gather}
Se $\Sigma$ è una \textit{catena} (famiglia \textit{totalmente ordinata}, ovvero ogni elemento è \textit{confrontabile}) in $\mathcal{I}$ si ha che
\begin{gather}
    F_\Sigma := \bigcup_{F \in \Sigma} F \notag
\end{gather}
è contenuto in $\mathcal{I}$ ed è un estremo superiore per $\Sigma$ (rispetto all'inclusione). Segue dal \textit{lemma di Zorn} che $\mathcal{I}$ ammette un \textit{elemento massimale}, ovverlo esiste $\mathcal{B}^* \in \mathcal{I}$ tale che
\begin{gather}
    \mathcal{B}^* \subseteq \mathcal{B} \implies \mathcal{B}^* = \mathcal{B} \notag
\end{gather}
Dato che siamo interessati al comportamento rispetto a operazioni di intersezione, non è restrittivo supporre che $\mathcal{B}$ sia proprio l'elemento massimale di $\mathcal{I}$. Pertanto segue che: \\

$(a)$ ogni sottoinsieme di $X$ che sia intersezione di un numero finito di elementi di $\mathcal{B}$ appartiene a $\mathcal{B}$.

$(b)$ se un sottoinsieme di $X$ ha intersezione non vuota con ogni elemento di $\mathcal{B}$ appartiene a $\mathcal{B}$ \\ \\
Sia $\mu \in M$ e $p_\mu : X \longrightarrow X_\mu$ la proiezione canonica, consideriamo la famiglia di sottoinsiemi di $X_\mu$ data da
\begin{gather}
    \mathcal{B}_\mu := \{ p_\mu(B) : B \in \mathcal{B} \} \notag
\end{gather}
$\mathcal{B}_\mu$ ha la proprietà dell'intersezione finita, infatti
\begin{gather}
    p_\mu(B_1) \cap \dots \cap p_\mu(B_n) = p_\mu(B_1 \cap \dots \cap B_n) \neq \emptyset \notag
\end{gather}
per ogni famiglia finita $\{ B_1, \dots, B_n \} \subset \mathcal{B}$. Inoltre, poiché $X_\mu$ è compatto si ha
\begin{gather}
    I_\mu := \bigcap_{B \in \mathcal{B}} \overline{p_\mu(B)} \neq \emptyset \notag
\end{gather}
Scegliamo un punto $x_\mu \in I_\mu$ per ogni $\mu \in M$, sia $x$ il punto di $X$ definito da
\begin{gather}
    x(\mu) = x_\mu, \quad \mu \in M \notag
\end{gather}
Dimostriamo che $x \in \overline{B}$ per ogni $B \in \mathcal{B}$. Scegliamo arbitrariamente $\mu \in M$ e consideriamo un aperto $U_\mu$ di $X_\mu$ contenente $x_\mu$. Allora si ha
\begin{gather}
    U_\mu \cap p_\mu(B) \neq \emptyset \text{ per ogni } B \in \mathcal{B} \notag
\end{gather}
Segue che
\begin{gather}
    p_\mu^{-1}(U_\mu) \cap B \neq \emptyset \text{ per ogni } B \in \mathcal{B} \notag
\end{gather}
dalla $(b)$ segue che $p_\mu^{-1}(U_\mu) \in \mathcal{B}$ per ogni $\mu \in M$. Dalla $(a)$ segue che ogni intersezione $U^*$ di una numero finito di insiemi $p_\mu^{-1}(U_\mu)$, con $\mu \in M$, appartiene a $\mathcal{B}$. Inoltre, dato che $x \in p_\mu^{-1}(U_\mu)$ per ogni $\mu \in M$, si ha che $U^*$ forma un sistema fondamentale di intorni di $x$ che interseca ogni $B \in \mathcal{B}$. Di conseguenza ogni intorno di $x$ interseca ogni $B \in \mathcal{B}$, segue che $x \in \overline{B}$ per ogni $B \in \mathcal{B}$. \qed \\ \\ \\
\textbf{Corollario 1.5.3:} un sottoinsieme $K$ di $\mathbf{R}^n$ è compatto se e solo se è chiuso e limitato. \\

\textbf{Dimostrazione:} sia $K$ compatto, poiché $\mathbf{R}^n$ è di Hausdorff, si ha che $K$ è chiuso, inoltre è limitato poiché altrimenti la famiglia dei dischi aperti centrati nell'origine di raggio crescente sarebbe un ricoprimento aperto che non ammette un sottoricoprimento finito. \\ \\
Sia $K$ chiuso e limitato, allora $K \subset \mathbf{I}_1 \times \dots \times \mathbf{I}_1$ ($n$ copie di $\mathbf{I}_1$) che è compatto per il \textit{teorema di Tichonov}, inoltre $K$ è chiuso in $\mathbf{I}_1 \times \dots \times \mathbf{I}_1$ poiché è chiuso in $\mathbf{R}^n$, pertanto $K$ è un sottoinsieme chiuso di uno spazio compatto, segue che $K$ è compatto. \qed
\subsubsection{Osservazioni e complementi}
\textbf{1.} Dal \textit{corollario 1.5.1} segue che gli unici intervalli compatti di $\mathbf{R}$ sono gli intervalli chiusi e limitati. Inoltre due intervalli chiusi e limitati qualsiasi sono omeomorfi, pertanto, poiché la compattezza è una proprietà topologica, si ha che se $\mathbf{J}_1, \mathbf{J}_2 \subset \mathbf{R}$ sono intervalli omeomorfi e $\mathbf{J}_1$ è compatto, anche $\mathbf{J}_2$ è compatto. \\ \\
\textbf{2.} Dal \textit{corollario 1.5.3} segue che $D_r(\mathbf{x})$, $\mathbf{x} \in \mathbf{R}^n$, non è compatto poiché è aperto. Un ricoprimento aperto che non ammette sottoricoprimento finito è
\begin{gather}
    \{ D_\rho (\mathbf{x}) : 0 < \rho < r \} \notag
\end{gather}
Mentre la $n$-cella $\mathbf{D}^n$ è compatta (chiusa e limitata) e dunque non è omeomorfa a nessun disco aperto. \\ \\
\textbf{3.} La sfera $\mathbf{S}^n$, $n \geq 1$, è compatta poiché chiusa e limitata in $\mathbf{R}^{n + 1}$. \\
Il \textit{toro} $n$-dimensionale
\begin{gather}
    \mathbf{T}^n = \mathbf{S}^1 \times \dots \times \mathbf{S}^1 \notag
\end{gather}
è compatto per il \textit{teorema di Tichonov}. \\ \\
\textbf{4.} Il gruppo ortogonale $\mathbf{O}_n(\mathbf{R})$ è compatto. Infatti la condizione che definisce le matrici $A = (a_{ij}) \in \mathbf{O}_n(\mathbf{R}^n)$ è $AA^t = \mathbf{I}_n$, cioè:
\begin{gather}
    \sum_h a_{hi}a_{jh} = \delta_{ij}, \quad 1 \leq i, j \leq n \notag
\end{gather}
In particolare
\begin{gather}
    \sum_h a^2_{hi} = 1, \quad i = 1, \dots, n \notag
\end{gather}
ovvero
\begin{gather}
    \sum_{h, i} a^2_{hi} = n \notag
\end{gather}
Pertanto $\mathbf{O}_n(\mathbf{R})$ è un sottoinsieme chiuso di $\mathbf{R}^{n^2}$ contenuto nel disco chiuso $\overline{D}_{\sqrt{n}}(0)$. Pertanto $\mathbf{O}_n(\mathbf{R})$ è chiuso e limitato, cioè compatto. In modo analogo si dimostra per $\mathbf{SO}_n(\mathbf{R})$. \\ \\
\textbf{5.} Ogni spazio compatto di Hausdorff $X$ è normale. Infatti, siano $A, B$ sottoinsiemi chiusi, non vuoti e disgiunti di $X$, sia $b \in B$, per ogni $a \in A$ esistono aperti disgiunti $V_a, W_a$ tali che $a \in V_a$ e $b \in W_a$. La famiglia
\begin{gather}
    \{ V_a : a \in A \} \notag
\end{gather}
è un ricoprimento aperto di $A$, che è compatto, pertanto esistono $a_1, \dots, a_s \in A$ tali che
\begin{gather}
    A \subset V_{a_1} \cup \dots \cup V_{a_s} \notag
\end{gather}
e inoltre, posti $V_b := V_{a_1} \cup \dots \cup V_{a_s}$ e $W_b := W_{a_1} \cap \dots \cap W_{a_s}$, si ha
\begin{gather}
    A \subset V_b, \quad b \in W_b, \quad V_b \cap W_b = \emptyset \notag
\end{gather}
Analogamente, la famiglia
\begin{gather}
    \{ W_b : b \in B \} \notag
\end{gather}
è un ricoprimento aperto di $B$, che è compatto, dunque esistono $b_1, \dots, b_k \in B$ tali che
\begin{gather}
    B \subset W_{b_1} \cup \dots \cup W_{b_k} \notag
\end{gather}
Pertanto, posti $V := V_{b_1} \cap \dots \cap V_{b_k}$ e $W := W_{b_1} \cup \dots \cup W_{b_k}$, si ha
\begin{gather}
    A \subset V, \quad B \subset W, \quad V \cap W = \emptyset \notag
\end{gather}
ovvero $X$ è uno spazio normale.
\subsubsection{Esercizi}
\textbf{1.} Dimostrare che un sottoinsieme $K$ chiuso e discreto di uno spazio compatto $X$ è finito. \\

\textbf{Dimostrazione:} $K$ è compatto poiché sottoinsieme chiuso di uno spazio compatto. Inoltre $K$ è sicuramente finito, altrimenti $\{ \{ x \} : x \in K \}$ sarebbe un ricoprimento aperto che non ammette sottoricoprimento finito. \\ \\ \\
\textbf{2.} Siano $X$ uno spazio topologico e $K_1, \dots, K_s$ sottoinsiemi compatti. Dimostrare che $K := K_1 \cup \dots \cup K_s$ è compatto. \\

\textbf{Dimostrazione:} sia $\{ A_i : i \in I \}$ un ricoprimento aperto di $K$, ovviamente è un ricoprimento aperto anche per $K_j$, per ogni $j = 1 \dots s$. Pertanto esiste finito $h(j) \subset I$ tale che
\begin{gather}
    K_j \subset \bigcup_{h \in h(j)} A_h \notag
\end{gather}
segue che
\begin{gather}
    K \subset \bigcup_{h \in H} A_h, \quad H = \bigcup_{j = 1}^s h(j) \notag
\end{gather}
e $\{ A_h : h \in H \}$ è un sottoricoprimento finito contenuto in $\{ A_i : i \in I \}$. \\ \\ \\
\textbf{3.} I seguenti sottoinsiemi di $\mathbf{R}^n$, $n \geq 2$, non sono compatti:
\begin{gather}
    \mathbf{D}^n \setminus \{ 0 \}, \quad \mathbf{S}^{n - 1} \setminus \{ (1, 0, \dots, 0) \}, \quad \{ x \in \mathbf{R}^n : x_n = 0 \}, \quad \mathring{\mathbf{I}}^n \notag
\end{gather}
Esibire per ognuno un ricoprimento aperto che non ammette sottoricoprimento finito. \\

\textbf{Soluzione:} $\mathbf{R}^n$ è uno spazio di Hausdorff. Per ogni $\mathbf{x} \in \mathbf{D}^n \setminus \{ 0 \}$ scegliamo $V_\mathbf{x}, W_\mathbf{x}$ aperti disgiunti tali che
\begin{gather}
    0 \in V_\mathbf{x}, \quad \mathbf{x} \in W_\mathbf{x} \notag
\end{gather}
Si ha che $\{ W_\mathbf{x} : \mathbf{x} \in \mathbf{D}^n \setminus \{ 0 \} \}$ è un ricoprimento aperto di $\mathbf{D}^n \setminus \{ 0 \}$. Supponiamo che ammetta sottoricoprimento finito, ovvero che esistano $\mathbf{x}_1, \dots, \mathbf{x}_s$ tali che
\begin{gather}
    \mathbf{D}^n \setminus \{ 0 \} \subset W_{\mathbf{x}_1} \cup \dots \cup W_{\mathbf{x}_s} \notag
\end{gather}
Posti
\begin{gather}
    V := V_{\mathbf{x}_1} \cap \dots \cap V_{\mathbf{x}_s}, \quad W := W_{\mathbf{x}_1} \cup \dots \cup W_{\mathbf{x}_s} \notag
\end{gather}
si ottiene che $V$ è un intorno aperto di $\{ 0 \}$ tale che
\begin{gather}
    V \cap \mathbf{D}^n \setminus \{ 0 \} \neq \emptyset, \quad V \cap W = \emptyset \notag
\end{gather}
tuttavia ciò è un assurdo poiché $\mathbf{D}^n \setminus \{ 0 \} \subset W$. \\ \\
Facendo un ragionamento del tutto analogo è possibile costruire un ricoprimento aperto di $\mathbf{S}^{n - 1} \setminus \{ (1, 0, \dots, 0) \}$ che non ammette sottoricoprimento finito. \\ \\
Sia $\Sigma := \{ x \in \mathbf{R}^n : x_n = 0 \}$. Si ha che
\begin{gather}
    \{ D_r(0) \cap \Sigma : r > 0 \} \notag
\end{gather}
sono le proiezioni delle sfere aperte di raggio $r$ centrate nell'origine su $\Sigma$, ovvero sono dischi $(n - 1)$-dimensionali di raggio $r$ centrati nell'origine e giacenti su $\Sigma$ stesso. Si nota subito che esse formano un ricoprimento aperto di $\Sigma$, infatti è come considerare la famiglia dei dischi aperti $(n - 1)$-dimensionali di raggio $r > 0$ centati nell'origine in $\mathbf{R}^{n - 1}$. Tuttavia non ammette un sottoricoprimento finito, infatti per ogni famiglia finita $r_1, \dots, r_s \in \mathbf{R}$, posto $\overline{r} := \sup \{ r_j : j = 1, \dots, s \}$, si ha che
\begin{gather}
    (\overline{r}, 0, \dots, 0) \not\in \bigcup_{i = 1}^s (D_{r_i}(0) \cap \Sigma) \notag
\end{gather}
ma
\begin{gather}
    (\overline{r}, 0, \dots, 0) \in \Sigma \notag
\end{gather}
per definizione. \\ \\
Si ha che $\{ (0 + \varepsilon, 1 - \varepsilon)^n : 0 < \varepsilon < \frac{1}{2} \}$ è un ricoprimento aperto di $\mathring{\mathbf{I}}^n$, tuttavia l'unione di ogni sottofamiglia finita è ancora un elemento di $\{ (0 + \varepsilon, 1 - \varepsilon)^n : 0 < \varepsilon < \frac{1}{2} \}$, pertanto non ricopre $\mathring{\mathbf{I}}^n$. \\ \\ \\
\textbf{4.} Sia $X$ uno spazio metrizzabile, siano $A, B$ sottoinsieme chiusi, non vuoti di $X$. Posto $d(A, B) := \inf \{ d(a, b) : a \in A, b \in B \}$, dimostrare che se $A$ (oppure $B$) è compatto si ha $d(A,B) = 0$ se e solo se $A \cap B \neq \emptyset$. \\

\textbf{Dimostrazione:} fissiamo $A$ compatto. $(\rightarrow)$ Sia $d(A, B) = 0$, per ogni $a \in A$ poniamo
\begin{gather}
    d_B(a) = \inf \{ d(a, b) : b \in B \} \notag
\end{gather}
$d_B$ è continua e $A$ è compatto, per il \textit{teorema di Weierstrass} esiste $a_0 \in A$ tale che 
\begin{align}
    d_B(a_0) \leq d_B(a) \notag
\end{align}
per ogni $a \in A$. Segue che 
\begin{align}
    d_B(a_0) = \inf\{d(a, b) : a \in A, b \in B\} = d(A, B) = 0 \notag
\end{align}
Mostriamo che $a \in A \cap B$. Supponiamo per assurdo che $A \cap B = \emptyset$, allora $A \subset X \setminus B$, dove $X \setminus B$ è aperto dato che $B$ è chiuso per ipotesi. Segue che esiste $\varepsilon > 0$ tale che $D_\varepsilon(a) \subset X \setminus B$ e quindi $d_B(a) > \varepsilon$, che è assurdo. \\ \\
$(\leftarrow)$ Sia $A \cap B \neq \emptyset$, allora esistono $a \in A$ e $b \in B$ tali che $a = b$, pertanto $d(a, b) = 0$ per definizione, segue $d(A, B) = 0$. \\ \\ \\
\textbf{5.} Sia $G$ un gruppo di matrici compatto. Allora $|\det A| = 1 \ \forall A \in G$. \\

\textbf{Dimostrazione:} siano $A,\, B \in G$, poiché per ipotesi $G$ è un gruppo di matrici si ha che $A \cdot B \in G$, dove $\cdot : G \times G \longrightarrow G$ è il prodotto di matrici, inoltre per il teorema di Binet si ha che $\det (A \cdot B) = (\det A)(\det B)$. \\
Supponiamo per assurdo che esiste $A \in G$ tale che $|\det A| = a > 1$, ciò implica che
\begin{gather}
    \det G := \{ \det A : A \in G \} \subseteq \mathbf{R} \notag
\end{gather}
non è compatto. Infatti per ogni $M \in \mathbf{R}$ esiste $n \in \mathbf{N}$ tale che
\begin{gather}
    |\det A^n| = a^n > M \notag
\end{gather}
basta scegliere $n > \log_a M$, perciò $\det G$ non è limitato e per il teorema di Heine-Borel non è compatto. \\
Tuttavia $\det : \mathbb{M}_n \longrightarrow \mathbf{R}$ è continua, dunque per la \textit{proposizione 1.5.3} si ha che $\det(G)$ è compatto, ma ciò è un assurdo. \\ \\
Supponiamo ora che esiste $A \in G$ tale che $|\det A| < 1$, poiché per ipotesi $G$ è un gruppo si ha che $A^{-1} \in G$ e, conseguenza del teorema di Binet, $|\det A^{-1}| = \frac{1}{|\det A|} > 1$, pertanto, procedendo con un ragionamento analogo al precedente, si ottiene la tesi. \\ \\
\textbf{6.} Sia $X$ uno spazio topologico compatto, $U \subset X$ un aperto e
\begin{gather}
    \{ C_i : i \in I \} \notag
\end{gather}
una famiglia di sottospazi chiusi di $X$. Dimostrare che esiste un insieme finito di indici $i_1, 1_2, \dots, i_n \in I$ tali che
\begin{gather}
    C_{i_1} \cap C_{i_2} \cap \dots C_{i_n} \subset U \notag
\end{gather} \\

\textbf{Dimostrazione:} consideriamo il ricoprimento aperto
\begin{gather}
    U \cup \{ X \setminus C_i : i \in I \} \notag
\end{gather}
esso ammette sottoricoprimento finito tale che
\begin{gather}
    U \cup X \setminus (C_{i_1} \cap \dots \cap C_{i_n}) = X \notag
\end{gather}
Allora si ottiene $C_{i_1} \cap \dots \cap C_{i_n} \subset U$.

\subsection{Generalizzazioni della compattezza}
\textbf{Definizione 1.6.1:} uno spazio topologico $X$ è detto \textit{localmente compatto} se ogni elemento $p \in X$ ammette un intorno compatto. \\ \\
Esempi di insiemi localmente compatti sono $\mathbf{R}^n$ per ogni $n \in \mathbf{N}$, gli aperti di $\mathbf{R}^n$ e in generale le varietà topologiche, mentre $\mathbf{Q}^n$ non è localmente compatto per ogni $n \in \mathbf{N}$. \\ \\ \\
\textbf{Proposizione 1.6.1:} $(a)$ ogni sottospazio chiuso $C$ di uno spazio localmente compatto $X$ è localmente compatto. \\ \\
$(b)$ sia $X$ uno spazio localmente compatto di Hausdorff. Ogni punto $p \in X$ possiede un sistema fondamentale di intorni relativamente compatti. \\

\textbf{Dimostrazione:} $(a)$ sia $p \in C$, sia $K \subset X$ un intorno compatto di $p$, si ha che $K \cap C$ è un intorno di $p$ in $C$ ed è compatto. \\ \\
$(b)$ Sia $p \in X$, per ipotesi $p$ possiede un intorno compatto $K \subset X$. Sia $U$ un aperto contenente $p$, si ha che $V := \mathring{K} \cap U$ è un aperto contenente $p$, inoltre $\overline{V} \subset \overline{K} = K$, ovvero $\overline{V}$ è un chiuso contenuto in insieme compatto $K$ di uno spazio di Hausdorff e quindi $\overline{V}$ è compatto. \qed \\ \\ \\
\textbf{Proposizione 1.6.2:} sia $X$ uno spazio topologico localmente compatto che soddisfa il secondo assioma di numerabilità, allora $X$ è unione di una famiglia numerabile di sottoinsiemi compatti $\{K_i : i \in \mathbf{N}\}$ tale che $K_i \subset \mathring{K}_{i + 1}$ per ogni $i \in \mathbf{N}$. \\

\textbf{Dimostrazione:} sia $x \in X$, sia $U$ un aperto contenente $x$ e sia $K$ un intorno compatto di $x$. Per ipotesi $X$ possiede una base numerabile di aperti $\mathcal{B}$, allora esiste $U_i \in \mathcal{B}$ tale che $U_i \subset \mathring{K} \cap U$. Inoltre per il punto $(b)$ della \textit{proposizione 1.6.1} si ha che $U_i$ è relativamente compatto in $\mathring{K} \cap U$. Dunque per ogni $x \in X$ e ogni aperto $U$ contenente $x$ esiste $U_i \in \mathcal{B}$ relativamente compatto tale che $x \in U_i \subset U$. Allora $\{U_i : i \in \mathbf{N}\}$ è una base numerabile di aperti relativamente compatti di $X$. Per ogni $i \in \mathbf{N}$ si ha che $\overline{U}_i$ è ricoperto da un numero finito di $U_j$, la cui unione ha chiusura compatta $H_i$ dato che si tratta di insiemi relativamente compatti. Pertanto
\begin{align}
    \overline{U}_i \subset \mathring{H}_i \notag
\end{align}
per ogni $i \in \mathbf{N}$ e quindi
\begin{align}
    X = \bigcup_{i \in \mathbf{N}}\mathring{H}_i \notag
\end{align}
A questo punto definiamo induttivamente
\begin{align}
    K_1 := H_1 \qquad K_{i + 1} := \bigcup_{j \in J(i)}H_j \notag
\end{align}
dove $J(i)$ è un sottoinsieme di $\mathbf{N}$ tale che
\begin{align}
    K_i \subset \bigcup_{j \in J(i)}\mathring{H}_j \notag
\end{align}
La famiglia di sottoinsiemi compatti $\{K_i\}$ soddisfa le proprietà ricercate. \\ \\
L'esistenza di una famiglia di compatti con le proprietà enunciate nella \textit{proposizione 1.6.2} si può dimostrare direttamente nel caso $X = \mathbf{R}^n$, o più generale nel caso $X = S \subset \mathbf{R}^n$, con $S$ aperto o chiuso. Infatti se $X$ è chiuso basta prendere
\begin{align}
    K_i = X \cap \overline{D}_i(p) \notag
\end{align}
fissato un punto $p$ di $X$. Mentre se $X$ è aperto basta prendere
\begin{align}
    K_i = \{x \in X : \|x\| \leq i, \, d(x, \mathbf{R}^n \setminus X) \geq 1/i\} \notag
\end{align}
notiamo che in questo caso $K_i$ è chiuso nella topologia relativa ed è compatto poiché chiuso e limitato in $\mathbf{R}^n$. \qed \\ \\ \\
\textbf{Definizione 1.6.2:} sia $X$ uno spazio topologico e siano $\mathcal{U}, \mathcal{V}$ ricoprimenti aperti di $X$. $\mathcal{V}$ è detto \textit{raffinamento} di $\mathcal{U}$, o \textit{più fine} di $\mathcal{U}$, se ogni $V \in \mathcal{V}$ è contenuto in qualche aperto di $\mathcal{U}$. \\ \\ \\
\textbf{Definizione 1.6.3:} una famiglia $\mathcal{F}$ di sottoinsiemi di $X$ si dice \textit{localmente finita} se ogni $x \in X$ possiede un intorno $U_x$ che interseca solo con un numero finito di elementi di $\mathcal{F}$. \\ \\ \\
\textbf{Definizione 1.6.4:} uno spazio topologico $X$ si dice \textit{paracompatto} se ogni ricoprimento $\mathcal{U}$ di $X$ possiede un raffinamento $\mathcal{V}$ localmente finito. \\ \\
Ogni spazio $X$ compatto è anche paracompatto, infatti se $X$ è compatto ogni ricoprimento $\mathcal{U}$ di $X$ possiede un sottoricoprimento $\mathcal{V}$ finito, allora $\mathcal{V}$ è più fine di $\mathcal{U}$ (banale) e per ogni $x \in X$ si ha che ogni intorno di $x$ interseca solamente con un numero finito di elementi di $\mathcal{V}$, proprio perché $\mathcal{V}$ è un insieme finito, pertanto $\mathcal{V}$ è localmente finito. \\ \\ \\
\textbf{Proposizione 1.6.3:} sia $X$ uno spazio topologico localmente compatto di Hausdorff che soddisfa il secondo assioma di numerabilità. Allora ogni ricoprimento aperto $\mathcal{U}$ di $X$ possiede un raffinamento $\mathcal{V}$ localmente finito costituito al più da un'infinità numerabile di aperti. In particolare $X$ è paracompatto. \\

\textbf{Dimostrazione:} sia $\mathcal{U}$ un ricoprimento aperto di $X$, per la \textit{proposizione 1.6.2} $X$ è unione di una famiglia numerabile $\{K_i\}_{i \in I}$ di insiemi compatti tali che $K_i \subset \mathring{K}_{i + 1}$. Per ogni $i \in I$ la famiglia 
\begin{align}
    \mathcal{V}_i = \{[\mathring{K}_{i + 2} \setminus K_{i - 1}] \cap U : U \in \mathcal{U}\} \notag
\end{align}
dove si pone $K_j = \emptyset$ per $j \leq 0$, costituisce un ricoprimento aperto di $K_{i + 1} \setminus \mathring{K}_i$ che è chiuso in $K_{i + 1}$ e quindi compatto. Allora esiste un sottoricoprimento finito $V_{i, 1}, \dots, V_{i, \alpha(i)}$ tale che 
\begin{align}
    K_{i + 1} \setminus \mathring{K}_i \subset V_{i, 1} \cup \dots \cup V_{i, \alpha(i)} \notag
\end{align}
Poiché 
\begin{align}
    \bigcup_{i \in I}(K_{i + 1} \setminus \mathring{K}_i) = X \notag
\end{align}
si ha che la famiglia $\mathcal{V} = \{V_{i, j} : i \in I, j = 1, \dots, \alpha(i)\}$ è un ricoprimento aperto di $X$ composta da al più un'infinità numerabile di elementi. Inoltre $\mathcal{V}$ è più fine di $\mathcal{U}$ per costruzione. Sia $x \in X$, esiste $i \in I$ tale che $x \in K_i$, l'aperto $\mathring{K}_{i + 1}$ è un intorno aperto di $x$ tale che 
\begin{align}
    \mathring{K}_{i + 1} \cap V_{j, h} = \emptyset \notag
\end{align}
per ogni $j \geq  i + 2$, ovvero $\mathring{V}$ paracompatto. \qed \\ \\ \\
\textbf{Corollario 1.6.1:} ogni aperto o chiuso di $\mathbf{R}^n$ è paracompatto, ogni varietà topologica è paracompatta. \\ \\ \\
\textbf{Definizione 1.6.5:} uno spazio topologico $X$ è detto \textit{numerabilmente compatto} se ogni sottoinsieme infinito $Z \subset X$ possiede un punto di accumulazione. Uno spazio topologico $X$ è detto \textit{compatto per successioni} se ogni successione di punti in $X$ ammette una sottosuccessione convergente in $X$. \\ \\ \\
Ogni spazio compatto per successioni è anche numerabilmente compatto, inoltre per la \textit{proposizione 1.5.2} ogni spazio compatto è anche numerabilmente compatto. \\ \\ \\
\textbf{Definizione 1.6.6:} sia $X$ uno spazio metrizzabile e sia $\mathcal{U}$ un ricoprimento aperto di $X$. Se esiste $\varepsilon > 0$ tale che il ricoprimento aperto costituito dai dischi aperti di raggio $\varepsilon$ è un raffinamento di $\mathcal{U}$, $\varepsilon$ è detto numero di Lebesgue del ricoprimento $\mathcal{U}$. \\ \\ \\
\textbf{Proposizione 1.6.4:} sia $X$ uno spazio metrizzabile compatto per successioni, allora ogni ricoprimento aperto $\mathcal{U}$ di $X$ ammette numero di Lebesgue. \\

\textbf{Dimostrazione:} sia $\mathcal{U}$ un ricoprimento aperto di $X$. Per ogni $x \in X$ poniamo 
\begin{align}
    \varepsilon(x) = \sup\{\delta > 0 : D_\delta(x) \subset U, U \in \mathcal{U}\} \notag
\end{align}
inoltre sia 
\begin{align}
    \varepsilon_0 = \inf_{x \in X}\varepsilon(x) \notag
\end{align}
mostriamo che $\varepsilon_0 > 0$. Supponiamo per assurdo $\varepsilon_0 = 0$, allora esiste una successione $\{x_n : n  \in \mathbf{N}\}$ tale che 
\begin{align}
    \lim_{n \rightarrow \infty}\varepsilon(x_n) = 0 \notag
\end{align}
Per ipotesi di compattezza per successioni, a meno di sostituire $\{x_n\}$ con una sottosuccessione, esiste $x_0 \in X$ tale che 
\begin{align}
    x_0 = \lim_{n \rightarrow \infty}x_n \notag
\end{align}
Per definizione esiste $N \in \mathbf{N}$ tale che $d(x_0, x_n) < \frac{1}{4}\varepsilon(x_0)$ per ogni $n > N$. Allora si ha 
\begin{align}
    D_{\varepsilon(x_0)/4}(x_n) \subset D_{\varepsilon(x_0)/2}(x_0) \notag
\end{align}
per ogni $n > N$, quindi $\varepsilon(x_n) \geq \frac{1}{4}\varepsilon(x_0)$ per ogni $n > N$ (per definizione di numero di Lebesgue). Allora 
\begin{align}
    \lim_{n \rightarrow \infty}\varepsilon(x_n) \geq \frac{1}{4}\varepsilon(x_0) > 0 \notag
\end{align}
pertanto si ottiene un assurdo. \qed \\ \\ \\
\textbf{Teorema 1.6.1:} sia $X$ uno spazio metrizzabile, allora le seguenti sono equivalenti: \\
\textbf{1.} $X$ è compatto. \\
\textbf{2.} $X$ è numerabilmente compatto. \\
\textbf{3.} $X$ è compatto per successioni. \\

\textbf{Dimostrazione:} $(1 \rightarrow 2)$ segue dalla \textit{proposizione 1.5.2}. \\ \\
$(2 \rightarrow 3)$ Poiché $X$ è uno spazio metrizzabile è anche $T_1$ e soddisfa il primo assioma di numerabilità, pertanto per la \textit{proposizione 1.2.7} ogni successione che ha un punto di accumulazione ammette una sottosuccessione convergente, segue che $X$ è compatto per successioni. \\ \\
$(3 \rightarrow 1)$ Sia $\mathcal{U}$ un ricoprimento aperto di $X$, per la \textit{proposizione 1.6.4} $\mathcal{U}$ ammette un numero di Lebesgue $\varepsilon > 0$ tale che la famiglia $\mathcal{V}_\varepsilon$ dei dischi aperti di raggio $\varepsilon$ è un raffinamento di $\mathcal{U}$. Supponiamo per assurdo che $\mathcal{U}$ non ammetta sottoricoprimento finito, allora neanche $\mathcal{V}_\varepsilon$ ammette sottoricoprimento finito. Sia $x_1 \in X$, allora esiste $x_2 \in X$ tale che $x_2 \not\in D_\varepsilon(x_1)$. Analogamente, dato che $\{D_\varepsilon(x_1), D_\varepsilon(x_2)\}$ non è un ricoprimento, esiste $x_3 \in X$ tale che $x_3 \not\in D_\varepsilon(x_2) \cup D_\varepsilon(x_1)$. Ragionando per induzione, esiste $x_k \in X$ tale che
\begin{align}
    x_k \not\in \bigcup_{i = 1}^{k - 1}D_\varepsilon(x_k) \notag
\end{align}
Abbiamo costruito una successione di punti $\{x_n\}$ in $X$ tale che $d(x_k, x_j) > \varepsilon$ per ogni $k \neq j$. Per ipotesi $X$ è compatto per successioni, quindi $\{x_n\}$ ammette una sottosuccesione $\{x_{\nu(k)}\}$ convergente ad un opportuno punto $x \in X$, segue che esiste $N_{\varepsilon/2} \in \mathbf{N}$ tale che $d(x_{\nu(k)}, x) < \varepsilon$ per ogni $k > N_{\varepsilon/2}$. Allora, fissati $k_1, k_2 > N_{\varepsilon/2}$ si ha 
\begin{align}
    d(x_{\nu(k_1)}, x_{\nu(k_2)}) \leq d(x_{\nu(k_1)}, x) + d(x, x_{\nu(k_2)}) < \varepsilon \notag
\end{align}
ma ciò è assurdo. \qed \\ \\ \\
Un'altra caratterizzazione degli spazi compatti viene data attraverso la proprietà di completezza degli spazi metrici. \\ \\
\textbf{Definizione 1.6.7:} sia $X$ uno spazio metrico, una successione $\{x_n\}_{n \in \mathbf{N}}$ è detta di Cauchy se per ogni $\varepsilon > 0$ esiste $N_\varepsilon \in \mathbf{N}$ tale che $d(x_n, x_m) < \varepsilon$ per ogni $n, m > N_\varepsilon$. \\ \\ \\
\textbf{Osservazione:} ogni successione convergente è anche di Cauchy. Infatti data $\{x_n\}_{n \in \mathbf{N}}$ convergente a $x \in X$, per ogni $\varepsilon > 0$ esiste $N_\varepsilon \in \mathbf{N}$ tale che $d(x_n, x_m) \leq d(x_n, x) + d(x, x_m) < 2\varepsilon$ per ogni $n, m > N_\varepsilon$, per la disuguaglianza triangolare. \\ \\ \\
\textbf{Definizione 1.6.8:} uno spazio metrico è detto \textit{completo} se ogni successione di Cauchy è anche convergente. \\ \\ \\
Dai corsi di Analisi Matematica sappiamo che $\mathbf{R}^n, \mathbf{C}^n$ sono spazi metrici completi rispetto alla metrica indotta dalla norma per ogni $n \geq 1$ (partendo dal caso $n = 1$ e usando successioni a variazione finita, poi generalizzando al caso $n \geq 1$). \\ \\
La completezza non è una proprietà topologica, ovvero non si conserva rispetto a omeomorfismi. Ad esempio $\mathbf{R} \cong (0, 1)$, tuttavia $\mathbf{R}$ è completo, mentre $(0, 1)$ no, infatti $\{1/n\}_{n \in \mathbf{N}}$ è una successione di Cauchy in $(0, 1)$ ma non è ivi convergente. \\ \\ \\
\textbf{Definizione 1.6.9:} uno spazio metrico è detto \textit{totalmente limitato} se per ogni $\varepsilon > 0$ esiste un ricoprimento finito costituito da dischi chiusi di raggio $\varepsilon$. \\ \\ \\
Ovviamente ogni spazio metrico totalmente limitato è anche limitato dato che l'unione finita di insiemi limitati è ancora limitato. Tuttavia in generale il viceversa non è vero, esistono spazi metrici limitati che non sono totalmente limitati. \\ \\ \\
\textbf{Teorema 1.6.2:} uno spazio metrico è compatto se e solo se è completo e totalmente limitato. \\

\textbf{Dimostrazione:} $(\rightarrow)$ sia $(X, d)$ uno spazio metrico compatto. Per il \textit{teorema 1.6.1} $(X, d)$ è compatto per successioni, pertanto ogni successione in $X$ ammette una sottosuccessione convergente ad un opportuno punto di $X$. Sia $\{x_n\}_{n \in \mathbf{N}}$ una successione di Cauchy in $X$. $\{x_n\}_{n \in \mathbf{N}}$ ammette una sottosuccessione $\{x_{\nu(k)}\}_{k \in \mathbf{N}}$ convergente ad un opportuno $x \in X$. Fissato $k$ abbastanza grande segue che 
\begin{align}
    d(x_n, x) \leq d(x_n, x_{\nu(k)}) + d(x_{\nu(k)}, x) < 2\varepsilon \notag
\end{align}
per ogni $n > N_\varepsilon$ (per ipotesi di Cauchy), ovvero $\{x_n\}_{n \in \mathbf{N}}$ converge a $x$, pertanto $(X, d)$ è completo. Inoltre $X$ è totalmente limitato, infatti ogni ricoprimento aperto ammette sottoricoprimento finito, in particolare anche il ricoprimento aperto costituito dai dischi di raggio $\varepsilon$ per ogni $\varepsilon > 0$. \\ \\
$(\leftarrow)$ Sia $(X, d)$ completo e totalmente limitato. Dimostriamo che $X$ è compatto per successioni e dal \textit{teorema 1.6.1} seguirà la tesi. Sia $\{x_n\}_{n \in \mathbf{N}}$ una successione di punti in $X$. Se $\{x_n\}_{n \in \mathbf{N}}$ ha un numero finito di punti distinti esiste certamente una sottosuccessione convergente. Supponiamo allora che $\{x_n\}_{n \in \mathbf{N}}$ abbia un numero infinito di punti distinti e costruiamo per induzione una sottosuccessione convergente. Per ogni $k \in \mathbf{N}$ consideriamo un ricoprimento chiuso di $X$ composto da un numero finito di dischi di raggio $\frac{1}{2^k}$. Sia $k = 1$, per ipotesi uno dei dischi contiene un numero infinito di termini della successione, chiamiamo $D_1$ tale disco ed indichiamo con $I_1$ l'insieme (infinito) degli indici che identificano i termini della successione contenuti in $D_1$. Sia $k = 2$, poiché $\{x_n\}_{n \in I_1}$ ha sempre infiniti termini distinti esiste un disco, che chiamiamo $D_2$, che contiene un numero infinito di termini della successione ed indichiamo con $I_2$ l'insieme degli indici che li identificano, ovviamente si ha $I_2 \subset I_1$. Per ogni $k \geq 2$ possiamo costruire induttivamente $I_k$ in modo analogo. Costruiamo una sottosuccessione $\{x_{n_k}\}_{k \in \mathbf{N}}$ con $n_1 < n_2 < \dots$ tale che $n_k \in I_k$ per ogni $k \in \mathbf{N}$. Allora per ogni $i, j > k$ si ha 
\begin{align}
    d(x_{n_i}, x_{n_j}) < \frac{1}{2^{k - 1}} \notag
\end{align}
ovvero $\{x_{n_k}\}_{k \in \mathbf{N}}$ è di Cauchy. Per la completezza di $X$ la sottosuccessione $\{x_{n_k}\}$ convergente ad un opportuno punto in $X$ e quindi $X$ è compatto per successioni. \qed \\ \\ \\
\textbf{Teorema 1.6.3:} sia $(X, d)$ uno spazio metrico completo. Per ogni successione $\{D_n = \overline{{D}_{r_n}(x_n)}\}$ tale che $D_1 \supset D_2 \supset \dots$ e $\lim_n r_n = 0$ si ha 
\begin{align}
    \bigcap_{n}D_n \neq \emptyset \notag
\end{align} \\

\textbf{Dimostrazione:} poiché per ipotesi $D_1 \supset D_2 \supset \dots$ si ha che $d(x_n, x_m) < |r_n - r_m|$ per ogni $n, m \in \mathbf{N}$. Allora $\{x_n\}_{n \in \mathbf{N}}$ è di Cauchy e quindi convergente, pertanto esiste $x \in X$ tale che $x_n \rightarrow x$ per $n \rightarrow \infty$. \\
Mostriamo che $x \in \bigcap_n D_n$. Per ogni $n \geq 1$ si ha che $x_n, x_{n + 1}, x_{n + 2}, \dots \in D_n$. Poiché $x$ è punto di accumulazione di $\{x_n\}$ si ha che $x \in D_n$ per ogni $n \geq 1$ dato che $D_n$ è chiuso. Allora $x \in \bigcap_n D_n$. \qed \\ \\ \\
\textbf{Definizione 1.6.10:} uno spazio topologico è detto di \textit{prima categoria} se è unione numerabile di sottoinsiemi chiuso aventi interno vuoto. \\
Altrimenti è detto di \textit{seconda categoria}. \\ \\ \\
\textbf{Teorema (di Baire):} ogni spazio completo è di seconda categoria. \\

\textbf{Dimostrazione:} sia $(X, d)$ uno spazio metrico completo. Supponiamo per assurdo che $X$ è di prima categoria. \\
Per ipotesi esiste una famiglia di sottoinsiemi chiusi $\{C_n\}_{n \in \mathbf{N}}$ di $X$ tali che $\mathring{C}_n = \emptyset$ per ogni $n \in \mathbf{N}$ e $X = \bigcup_n C_n$.



\newpage
\section{Gruppi di matrici}
\subsection{Appendice sui gruppi}
\textbf{Definizione A.1:} sia $G$ un insieme, un'operazione binaria è una funzione
\begin{gather}
    \cdot : G \times G \longrightarrow G \notag
\end{gather}
Siano $a, b \in G$, se $a \cdot b = c$ si dice che $c$ è il \textit{prodotto} di $a$ e $b$. \\ \\ \\
\textbf{Definizione A.2:} $G$ è detto \textit{gruppo} se è dotato di un'operazione binaria \\ $\cdot : G \times G \longrightarrow G$ tale che: \\
\textbf{1.} $a \cdot (b \cdot c) = (a \cdot b) \cdot c$ per ogni $a, b, c \in G$. \\
\textbf{2.} esiste $e \in G$ tale che $a \cdot e = e \cdot a = a$ per ogni $a \in G$. \\
\textbf{3.} esiste $a^{-1} \in G$ tale che $a \cdot a^{-1} = a^{-1} \cdot a = e$ per ogni $a \in G$. \\ \\ \\
\textbf{Osservazione:} l'elemento neutro e l'elemento inverso sono unici. \\ \\ \\
\textbf{Definizione A.3:} sia $G$ un gruppo, $S \subset G$ è detto \textit{sottogruppo} (scritto $S \leq G$) se $S$ è un gruppo rispetto alla stessa operazione binaria definita per $G$. \\ \\ \\
\textbf{Proposizione A.1:} sia $G$ un gruppo, $S \subset G$ è un sottogruppo se e solo se: \\
\textbf{1.} $x \cdot y \in S$ per ogni $x, y \in S$. \\
\textbf{2.} $e \in S$. \\
\textbf{3.} $x \in S$ implica $x^{-1} \in S$. \\

\textbf{Dimostrazione:} segue immediatamente dalla definizione di gruppo. \qed \\ \\ \\
\textbf{Definizione A.4:} siano $G, H$ gruppi, una funzione $f : G \longrightarrow H$ è detta \textit{omomorfismo} se $f(a \cdot b) = f(a) \cdot f(b)$ per ogni $a, b \in G$, ovvero se e solo se $f$ \textit{preserva la struttura di gruppo}. Se $f$ è anche biunivoca $f$ è detta \textit{isomorfismo}. \\ \\ \\
\textbf{Definizione A.5:} sia $n \in \mathbf{N}$. Il gruppo simmetrico (di rango $n$) è
\begin{gather}
    S_n := \{ f : [n] \longrightarrow [n] : f \text{ biunivoca} \} \notag
\end{gather}
con $[n] = \{1, 2, \dots, n \}$. Ovvero $S_n$ è l'insieme delle permutazioni definite su un insieme ordinato contenente $n$ elementi. \\
Se $f \in S_n$ scriviamo $f = a_1a_2a_3 \dots a_n$ per dire $f(i) = a_i$ per ogni $i = 1, \dots, n$. \\ \\ \\
\textbf{Definizione A.6:} siano $f, g \in S_n$, la composizione di $f$ e $g$ è
\begin{gather}
    f \cdot g : [n] \longrightarrow [n] \notag
\end{gather}
dove $(f \cdot g) (i) := f(g(i))$ per ogni $i \in [n]$. \\ \\ \\
\textbf{Osservazione:} sfruttando la notazione definita si ha che se $f, g \in S_n$, \\ $f = a_1 \dots a_n$ e $g = b_1 \dots b_n$ si ha che $f \cdot g = a_{b_1} \dots a_{b_n}$. \\ \\ \\
\textbf{Teorema A.1:} $S_n$ è un gruppo rispetto alla composizione. \\

\textbf{Dimostrazione:} siano $f, g \in S_n$, poiché sono entrambe biunivoche, anche la composizione è biunivoca, pertanto $f \cdot g \in S_n$. \\
Siano $f, g, h \in S_n$, si ha che
\begin{gather}
    f \cdot (g \cdot h)(i) = f((g \cdot h)(i)) = f(g(h(i))) \notag \\
    (f \cdot g) \cdot h(i) = (f \cdot g)(h(i)) = f(g(h(i))) \notag
\end{gather}
pertanto soddisfa l'associatività. \\
Sia $\mathbf{1}_n$ definita pondendo $\mathbf{1}_n(i) = i$ per ogni $i \in [n]$. $\mathbf{1}_n$ è biunivoca, pertanto $\mathbf{1}_n \in S_n$. Inoltre
\begin{gather}
    f \cdot \mathbf{1}_n (i) = f(\mathbf{1}_n(i)) = f(i) = \mathbf{1}_n (f(i)) = \mathbf{1}_n \cdot f (i) \notag
\end{gather}
quindi esiste l'elemento neutro. \\
Infine, sia $f \in S_n$, $f = a_1 \dots a_n$, e sia $g := [n] \longrightarrow [n]$ definita ponendo
\begin{gather}
    g(a) = b \text{ tale che } f(b) = a \notag
\end{gather}
Dunque, $f \cdot g (i) = f(g(i)) = i$ e $g \cdot f (i) = g(f(i)) = i$ per definizione, pertanto esiste l'inverso per ogni elemento di $S_n$. \qed \\ \\ \\
\textbf{Osservazione:} sia $f = a_1 \dots a_n$, allora $f^{-1} = b_1 \dots b_n$ dove $b_i$ è la posizione (numerata) di $i$ nella seguenza $a_1 \dots a_n$. \\
Ad esempio se $f = 614587329$ si ha $f^{-1} = 287341659$. \\ \\ \\
\textbf{Definizione A.7:} un gruppo $G$ è detto \textit{abeliano} (o \textit{commutativo}) se $ab = ba$ per ogni $a, b \in G$. \\ \\ \\
\textbf{Definizione A.8:} due gruppi $G, H$ sono detti \textit{isomorfi}, ovvero $G \cong H$, se esiste un isomorfismo $f : G \longrightarrow H$. \\ \\ \\
\textbf{Teorema (di Cayley):} sia $G$ un gruppo finito, allora esiste $n \in \mathbf{N}$ e $S \leq S_n$ tali che $G$ è isomorfo a $S$. \\

\textbf{Dimostrazione:} omessa. \qed
\newpage
\subsection{Gruppi ortogonali}
Tratteremo i vettori come vettori riga, ovvero
\begin{gather}
    x = (x_1, \dots, x_n) \notag
\end{gather}
dove $x_i$ per $i = 1, \dots, n$ sono le componenti di $x$. Pertanto la moltiplicazione per matrici sarà sinistra, ovvero
\begin{gather}
    xA = \begin{pmatrix}
        \langle x, A_1 \rangle \\
        . \\
        . \\
        . \\
        \langle x, A_n \rangle
    \end{pmatrix} \notag
\end{gather}
dove $\langle \cdot, \cdot \rangle : \mathbf{R}^n \times \mathbf{R}^n \longrightarrow \mathbf{R}$ è il prodotto scalare euclideo. \\ \\ \\
\textbf{Definizione 2.2.1:} sia $n \in \mathbf{N}$, poniamo
\begin{gather}
    \mathbf{M}_n(\mathbf{R}) := \{ \text{ matrici $n \times n$ a coefficienti reali} \} \notag
\end{gather}
Se $A \in \mathbf{M}_n(\mathbf{R})$ si ha che $a_{ij}$ è l'elemento in riga $i$ e colonna $j$; $A_i$ è l'$i$-esima riga. \\ \\ \\
\textbf{Definizione 2.2.2:} il \textit{gruppo generale lineare} (di rango $n$) è
\begin{gather}
    \mathbf{GL}_n(\mathbf{R}) := \{ A \in \mathbf{M}_n(\mathbf{R}) : A \text{ invertibile} \} \notag
\end{gather}
che equivale a
\begin{gather}
    \mathbf{GL}_n(\mathbf{R}) := \{ A \in \mathbf{M}_n(\mathbf{R}) : \det A \neq 0 \} \notag
\end{gather} \\ \\
\textbf{Definizione 2.2.3:} il \textit{gruppo ortogonale} (di rango $n$) è
\begin{gather}
    \mathbf{O}_n(\mathbf{R}) := \{ A \in \mathbf{GL}_n(\mathbf{R}) : \langle x, y \rangle = \langle xA, yA \rangle \text{ per ogni } x, y \in \mathbf{R}^n \} \notag
\end{gather} \\ \\
\textbf{Teorema 2.2.1:} sia $A \in \mathbf{GL}_n(\mathbf{R})$, i seguenti punti sono equivalenti: \\
\textbf{1.} $A \in \mathbf{O}_n(\mathbf{R})$. \\
\textbf{2.} $\{ v_1A, \dots, v_nA \}$ è una base ortonormale se $\{ v_1, \dots, v_n \}$ lo è. \\
\textbf{3.} $\{ A_1, \dots, A_n \}$ (righe di $A$) formano una base ortonormale. \\
\textbf{4.} $AA^t = \mathbf{I}_n$. \\

\textbf{Dimostrazione:} $(1 \implies 2)$ Si ha che
\begin{gather}
    \langle v_i, v_j \rangle = 0 \text{ per ogni } i \neq j \notag
\end{gather}
per ipotesi segue che
\begin{gather}
    \langle v_iA, v_jA \rangle = 0 \text{ per ogni } i \neq j \notag
\end{gather}
dunque anche $\{ v_1A, \dots, v_nA \}$ è una base ortonormale. \\ \\
$(2 \implies 3)$ Consideriamo la base canonica di $\{ \mathbf{e}_1, \dots, \mathbf{e}_n \}$, si ha che
\begin{gather}
    A_i = \mathbf{e}_iA \text{ per ogni } i \in [n] \notag
\end{gather}
Allora $\{ A_1, \dots, A_n \} = \{ \mathbf{e}_1A, \dots, \mathbf{e}_nA \}$ è una base ortonormale. \\ \\
$(3 \implies 4)$ Si ha che
\begin{gather}
    (AA^t)_{ij} = \langle A_i, A_j \rangle  = \delta_{ij} \notag
\end{gather}
ovvero $AA^t = \mathbf{I}_n$. \\ \\
$(4 \implies 1)$ si ha che
\begin{gather}
    \langle vA, wA \rangle = \sum_i \sum_j \lambda_i\mu_j \langle \mathbf{e}_iA, \mathbf{e}_jA \rangle = \sum_i \sum_j \lambda_i\mu_j \langle A_i, A_j \rangle = \sum_i\sum_j \lambda_i\mu_j \delta_{ij} = \notag \\
    = \langle \sum_i \lambda_i\mathbf{e_i}, \sum_j \mu_i\mathbf{e}_j \rangle = \langle v, w \rangle \notag
\end{gather} \qed \\
\textbf{Osservazioni:} $\mathbf{GL}_n(\mathbf{R})$ è un gruppo rispetto al prodotto tra matrici, inoltre $\mathbf{O}_n(\mathbf{R}) \leq \mathbf{GL}_n(\mathbf{R})$. \\ \\ \\
\textbf{Proposizione 2.2.1:} $\mathbf{O}_n(\mathbf{R}) = \{ A \in \mathbf{GL}_n(\mathbf{R}) : \| x \| = \| xA \| \text{ per ogni } x \in \mathbf{R}^n \}$ \\

\textbf{Dimostrazione:} sia $A \in \mathbf{O}_n(\mathbf{R})$, si ha che
\begin{gather}
    \| x \| = \sqrt{\langle x, x \rangle} = \sqrt{\langle xA, xA \rangle} = \| xA \| \notag
\end{gather}
per definizione. \qed \\ \\ \\
\textbf{Definizione 2.2.4:} il \textit{gruppo speciale lineare} è
\begin{gather}
    \mathbf{SL}_n(\mathbf{R}) := \{ A \in \mathbf{M}_n(\mathbf{R}) : \det A = 1 \} \notag
\end{gather}
Notiamo che si ha $\mathbf{SL}_n(\mathbf{R}) \leq \mathbf{GL}_n(\mathbf{R})$. \\ \\ \\
\textbf{Proposizione 2.2.2:} sia $A \in \mathbf{O}_n(\mathbf{R})$, allora $|\det A| = 1$. \\

\textbf{Dimostrazione:} vale $\det \mathbf{I}_n = \det (AA^t) = (\det A)^2$. \qed \\ \\ \\
\textbf{Definizione 2.2.5:} il \textit{gruppo speciale ortogonale} è
\begin{gather}
    \mathbf{SO}_n(\mathbf{R}) := \{ A \in \mathbf{O}_n(\mathbf{R}) : \det A = 1 \} \notag
\end{gather}
Notiamo subito che $\mathbf{SO}_n(\mathbf{R}) \leq \mathbf{O}_n(\mathbf{R})$. \\ \\ \\
\textbf{Definizione 2.2.6:} $f : \mathbf{R}^n \longrightarrow \mathbf{R}^n$ si dice \textit{isometria} se $\| x - y \| = \| f(x) - f(y) \|$ per ogni $x, y \in \mathbf{R}^n$. \\ \\ \\
\textbf{Teorema 2.2.2:} sia $B \in \mathbf{O}_n(\mathbf{R})$ e sia $f$ un'isometria tale che $f(0) = 0$ (fissa l'origine). Allora: \\
\textbf{1.} $x \hookrightarrow xB$ è un'isometria. \\
\textbf{2.} esiste $A \in \mathbf{O}_n(\mathbf{R})$ tale che $f(x) = xA$ per ogni $x \in \mathbf{R}^n$. \\

\textbf{Dimostrazione:} $(1)$ si ha che
\begin{gather}
    \| x - y \| = \| (x - y)B \| = \| xB - yB \| \notag
\end{gather}
dunque $B$ è un'isometria. \\ \\
$(2)$ sia
\begin{gather}
    A = \begin{pmatrix}
        f(\mathbf{e}_1) \\
        . \\
        . \\
        . \\
        f(\mathbf{e}_n)
    \end{pmatrix} \notag
\end{gather}
$A \in \mathbf{O}_n(\mathbf{R})$ perché le righe di $A$ sono una base ortonormale, infatti
\begin{gather}
    \langle f(\mathbf{e}_i), f(\mathbf{e}_j) \rangle = \langle \mathbf{e}_i, \mathbf{e}_j \rangle = \delta_{ij} \notag
\end{gather}
Inoltre per ogni $x \in \mathbf{R}^n$ si ha $x = \sum \lambda_i \mathbf{e}_i$, dunque
\begin{gather}
    xA = (\sum_i \lambda_i \mathbf{e}_i)A = \sum_i \lambda_i A_i = \sum_i \lambda_i f(\mathbf{e}_i) \notag
\end{gather}
Sia $g$ un'isometria definita ponendo $g(x) = f(x)A^{-1}$ per ogni $x \in \mathbf{R}^n$, mostriamo che $g = \mathbf{I}_n$. Poiché $g$ conserva il prodotto scalare cartesiano si ha che
\begin{gather}
    g(\mathbf{e}_i) = f(\mathbf{e}_i)A^{-1} = \mathbf{e}_i \notag
\end{gather}
Dunque
\begin{gather}
    g(x) = g(\sum_i \lambda_i \mathbf{e}_i) = \sum_i \mu_i \mathbf{e}_i \notag
\end{gather}
con $\mu_i = \langle g(x), \mathbf{e}_i \rangle$, inoltre
\begin{gather}
    \mu_i = \langle g(x), \mathbf{e}_i \rangle = \langle g(x), g(\mathbf{e}_i) \rangle = \langle x, \mathbf{e}_i \rangle = \lambda_i \notag 
\end{gather}
ovvero $g(x) = x$ per ogni $x \in \mathbf{R}^n$, dunque $g = \mathbf{I}_n$. \qed \\ \\ \\
\textbf{Definizione 2.2.7:} una base ordinata ortonormale $\{ v_1, v_2, v_3 \}$ di $\mathbf{R}^3$ è detta \textit{destra} se $v_3 = v_1 \wedge v_2$. \\ \\ \\
\textbf{Proposizione 2.2.3:} sia $A \in \mathbf{O}_3(\mathbf{R})$, se $\{ A_1, A_2, A_3 \}$ è una base destra di $\mathbf{R}^3$ si ha che $A \in \mathbf{SO}_n(\mathbf{R})$. \\

\textbf{Dimostrazione:} immediata. \qed \\ \\
Notiamo che gli elementi di $\mathbf{SO}_3(\mathbf{R})$ sono le isometrie che fissano l'origine e mandano la base canonica di $\mathbf{R}^3$ in una base destra. \\ \\ \\
Poniamo $Isom(\mathbf{R}^n) := \{ f : \mathbf{R}^n \longrightarrow \mathbf{R}^n : f \text{ isometria} \}$. \\ \\
Osserviamo che $Isom(\mathbf{R}^n)$ è un gruppo rispetto alla composizione di funzioni. Inoltre $\mathbf{O}_(\mathbf{R}) \leq Isom(\mathbf{R}^n)$, infatti il gruppo ortogonale è composto dalle isometrie che fissano l'origine.  \\
Tuttavia ponendo $v := f(0)$ si ha che $g(x) := f(x) - v$ è un isometria che fissa l'origine per ogni isometria $f$, pertanto, per il \textit{teorema 2.2.2} esiste $A \in \mathbf{O}_n(\mathbf{R})$ tale che $g(x) = xA$ per ogni $x \in \mathbf{R}^n$, cioè $f(x) = xA + v$. \\
Pertanto, aggiungendo una dimesione, si ottiene
\begin{gather}
    (f(x), 1) = (x, 1) \begin{pmatrix}
        A & 0 \\
        v & 1 
    \end{pmatrix}
    = (xA + v, 1) \notag
\end{gather} \\ \\
\textbf{Teorema 2.2.3:} sia $A \in \mathbf{SO}_3(\mathbf{R})$. Allora esiste une retta $r$ passante per l'origine tale che $x \hookrightarrow xA$ è una rotazione attorno a $r$. \\

\textbf{Dimostrazione:} siano $\lambda_1, \lambda_2, \lambda_3 \in \mathbf{C}$ gli autovalori di $A$, mostriamo che almeno uno di loro è uguale a $1$. Poiché $A$ è a coefficienti reali almeno uno degli autovalori è reale, supponiamo, per fissare le idee, che $\lambda_1 \in \mathbf{R}$. \\
Notiamo che se $\lambda_1 \in \mathbf{R}$ si ha $|\lambda_1| = 1$ perché si conservano le distanze e i moduli. Infatti, sia $v_1 \in \mathbf{R}^3 \setminus \{ 0 \}$ autovettore corrispondente all'autovalore $\lambda_1$, si ha che
\begin{gather}
    \| v_1 \| = \| v_1A \| = \| \lambda_1 v_1 \| = |\lambda_1|\| v_1 \| \notag
\end{gather}
perciò $|\lambda_1| = 1$. Inoltre se $\lambda_2 \in \mathbf{C} \setminus \mathbf{R}$ si ha che $\lambda_3 = \overline{\lambda_2}$, quindi $\lambda_2 \cdot \lambda_3 \in \mathbf{R}$. Pertanto si ottiene
\begin{gather}
    \det A = \lambda_1 \cdot \lambda_2 \cdot \lambda_3 = \lambda_1 \cdot |\lambda_2|^2 = 1 \notag
\end{gather}
ovvero $\lambda_1 > 0 \implies \lambda_1 = 1$. \\
Consideriamo la retta $r := \langle v_1 \rangle$, dove $v_1$ è sempre l'autovettore corrispondente all'autovalore $\lambda_1 = 1$. Si ha che $vA = v$ per ogni $v \in \langle v_1 \rangle$, ovvero $x \hookrightarrow xA$ lascia fissi tutti i punti di $r$. Ciò comporta che $x \hookrightarrow xA$ lascia fissi anche tutti i punti di $r_\perp$, infatti se $u \in r_\perp$ si ha
\begin{gather}
    0 = \langle u, v \rangle = \langle uA, vA \rangle = \langle uA, v \rangle \notag
\end{gather}
per ogni $v \in r$. Pertanto $A$ è un'isometria che fissa l'origine e lascia invariate $r$ e $r_\perp$, ovvero $A$ è una rotazione attorno a $r$. \qed \\ \\ \\
\textbf{Definizione 2.2.8:} sia $X \subset \mathbf{R}^n$, il \textit{gruppo delle simmetrie di $X$} è
\begin{gather}
    Symm(X) := \{ f \in Isom(\mathbf{R}^n) : f(X) = X \} \notag
\end{gather}
Ad esempio
\begin{gather}
    Symm(\mathbf{S}^{n - 1}) = \{ f \in Isom(\mathbf{R}^n) : f(0) = 0 \} = \mathbf{O}_n(\mathbf{R}) \notag
\end{gather}
ovvero il gruppo delle simmetrie della sfera centrata nell'origine in $\mathbf{R}^n$ è composto da tutti i movimenti rigidi che fissano l'origine. \\ \\ \\
Poniamo
\begin{gather}
    Symm^+(X) := \{ f \in Symm(X) : \det (f) = 1 \} \notag
\end{gather}
e analogamente
\begin{gather}
    Symm^-(X) := \{ f \in Symm(X) : \det (f) = -1 \} \notag
\end{gather}
dove
\begin{gather}
    \det (f) = \det \begin{pmatrix}
        A & 0 \\
        v & 1
    \end{pmatrix}, \text{ se } f(x) = xA + v \notag
\end{gather} \\ \\
\textbf{Definizione 2.2.9:} sia $m \in \mathbf{N}$, $m \geq 3$, il \textit{gruppo diedrale} di ordine $2m$ è
\begin{gather}
    D_{2m} := Symm(m\text{-gono regolare}) \notag
\end{gather}
anche scritto $I_2(m)$. \\ \\ \\
\textbf{Definizione 2.2.10:} il \textit{gruppo ciclico} di ordine $m$ è
\begin{gather}
    \mathbf{Z}_m := Symm^+(m\text{-gono regolare}) \notag
\end{gather} \\
Con regolare si intende lati e angoli tutti uguali. \\ \\ \\
\textbf{Teorema (di Leonardo da Vinci):} sia $X \subset \mathbf{R}^2$ tale che $|Symm(X)| < +\infty$. Allora $Symm(X) \cong \mathbf{Z}_m$ o $Symm(X) \cong I_2(m)$ per qualche $m \geq 2$. \\

\textbf{Dimostrazione:} omessa. \qed \\ \\ \\
\textbf{Esempio:} $I_2(2)$ sono le simmetrie di un segmento di $\mathbf{R}^2$. \\ \\ \\
\textbf{Proposizione 2.2.4:} sia $X \subseteq \mathbf{R}^n$, si ha che
\begin{gather}
    Symm^+(X) = Symm(X) \notag
\end{gather}
oppure
\begin{gather}
    Symm(X) = Symm^+(X) \cup \left\{ \begin{pmatrix}
        -1 & 0 \\
        0 & \mathbf{I}_{n - 1}
    \end{pmatrix} A : A \in Symm^+(X) \right\} \notag
\end{gather} \\

\textbf{Dimostrazione:} omessa. \qed \\ \\ \\
\textbf{Teorema 2.2.3:} sia $X \subseteq \mathbf{R}^3$ tale che $|Symm^+(X)| < +\infty$, si ha che $Symm^+(X) \cong S_4$ oppure $A_4$ o $A_5$ o $\mathbf{Z}_m$ o $I_2(m)$ per qualche $m \geq 2$, con $A_n := \{ \sigma \in S_n : \sigma \text{ è pari} \}$. \\

\textbf{Dimostrazione:} omessa. \qed \\ \\ \\
\textbf{Definizione 2.2.11:} sia $X$ un poliedro in $\mathbf{R}^3$. $X$ è detto \textit{regolare} se valgono: \\
\textbf{1.} Tutte le facce sono congruenti. \\
\textbf{2.} In ogni vertice si incontrano lo stesso numero di spigoli o di lati. \\ \\ \\
\textbf{Teorema 2.2.4:} sia $X$ un poliedro regolare, si ha che
\begin{gather}
    X \in \{ \text{cubo, tetraedro, ottaedro, dodecaedro, icosaedro} \} \notag
\end{gather}
Inoltre $Symm^+(cubo) \cong S_4$, $Symm^+(tetraedro) \cong A_4$ e $Symm^+(icosaedro) \cong A_5$. \\

\textbf{Dimostrazione:} omessa. \qed \\ \\ \\
Dato $X \subset \mathbf{R}^3$ scriviamo un algoritmo per calcolare $Symm^+(X)$. \\ \\
\textbf{1.} Posizioniamo $X$ in modo tale che il baricentro corrisponda con l'origine degli assi. Perciò se $f \in Symm^+(X)$ si ha $f(0) = 0$, ovvero $Symm^+(X) \leq \mathbf{SO}_3(\mathbf{R})$. \\ \\
\textbf{2.} Per il \textit{teorema 2.2.3} si ha che gli elementi di $\mathbf{SO}_3(\mathbf{R})$ sono rotazioni attorno ad una retta $r$ passante per l'origine. Pertanto gli elementi di $Symm^+(X)$ sono rotazioni attorno a qualche retta $r$ passante per l'origine $\underline{0}$ che mandano $X$ in $X$. \\ \\
\textbf{3.} Determinare quante coppie $(\hat{u}, \theta) \in \mathbf{S}^2 \times \mathbf{S}^1$ ci sono tali che una rotazione di $\theta$ attorno a $\hat{u}$ manda $X$ in $X$. \\ \\
\textbf{4.} Se le coppie sono infinite si ha che $Symm^+(X)$ è un gruppo di matrici. L'analisi in questo caso viene omessa poiché complessa. \\ \\
\textbf{5.} Se le coppie sono finite, determinare il numero $n$ di coppie tali che $\theta \neq 0$ e dividerlo per due (infatti una rotazione di $\theta$ attorno a $\hat{u}$ è identica ad una rotazione di $2\pi - \theta$ attorno a $-\hat{u}$). \\ \\
\textbf{6.} Si ha che $n/2$ può essere $60, 24, 12$ o un numero diverso dai precedenti. Se $n/2 = 60$ si ha $Symm^+(X) \cong \mathbf{Z}_{60}$ (oppure $\cong I_2(30)$, oppure $\cong A_5$). Se $n/2 = 24$ si ha $Symm^+(X) \cong \mathbf{Z}_{24}$ (oppure $\cong I_2(12)$, oppure $\cong S_4$). Se $n/2 = 12$ si ha $Symm^+(X) \cong \mathbf{Z}_{12}$ (oppure $\cong I_2(6)$, oppure $\cong A_4$). Se $n/2 = m$ diverso dai numeri precedenti si ha $Symm^+(X) \cong \mathbf{Z}_m$ (oppure $\cong I_2(m/2)$). \\ \\
\textbf{7.} Determinare quante coppie $(\hat{u}, \theta)$ sono tali che $\theta = \pi$ (rotazioni che eseguite due volte danno l'identità). Se $n/2 = 60$ possono essere $1 \implies Symm^+(X) \cong \mathbf{Z}_{60}$, $31 \implies Symm^+(X) \cong I_2(30)$ oppure $15 \implies Symm^+(X) \cong A_5$. \\
Se $n/2 = 24$ possono essere $1 \implies Symm^+(X) \cong \mathbf{Z}_{24}$, $13 \implies Symm^+(X) \cong I_2(12)$ oppure $9 \implies Symm^+(X) \cong S_4$. \\
Se $n/2 = 12$ possono essere $1 \implies Symm^+(X) \cong \mathbf{Z}_{12}$, $7 \implies Symm^+(X) \cong I_2(6)$ oppure $3 \implies Symm^+(X) \cong A_4$. \\
Se $n/2 = m$ possono essere $\leq 1 \implies Symm^+(X) \cong \mathbf{Z}_m$ oppure $\geq 2 \implies Symm^+(X) \cong I_2(m/2)$.

\newpage
\subsection{Topologia dei gruppi di matrici}
Notiamo che $\mathbf{M}_n(\mathbf{R}) \cong \mathbf{R}^{n^2}$ pertanto possiamo trattare lo spazio delle matrici $n \times n$ come uno spazio topologico euclideo. \\ \\ \\
\textbf{Definizione 2.3.1:} un \textit{gruppo di matrici} è un sottogruppo $G$ di $\mathbf{GL}_n(\mathbf{R})$ tale che $G$ è chiuso in $\mathbf{GL}_n(\mathbf{R})$ (topologia relativa). \\ \\
Ovviamente $\mathbf{GL}_n(\mathbf{R})$ è chiuso in se stesso e quindi è un gruppo di matrici. \\ \\ \\
\textbf{Proposizione 2.3.1:} $\mathbf{O}_n(\mathbf{R}), \mathbf{SO}_n(\mathbf{R})$ e $\mathbf{SL}_n(\mathbf{R})$ sono gruppi di matrici. \\

\textbf{Dimostrazione:} sappiamo che $\mathbf{O}_n(\mathbf{R})$ è un sottogruppo di $\mathbf{GL}_n(\mathbf{R})$, mostriamo che è chiuso nella topologia relativa. Consideriamo la funzione
\begin{gather}
    f : \mathbf{M}_n(\mathbf{R}) \longrightarrow \mathbf{M}_n(\mathbf{R}), \quad f(A) = AA^t \notag
\end{gather}
ovviamente $f$ è continua (è un polinomio), inoltre $\{ \mathbf{I}_n \}$ è chiuso (per le proprietà di separazione) e
\begin{gather}
    f^{-1}(\{ \mathbf{I}_n \}) = \mathbf{O}_n(\mathbf{R}) \notag
\end{gather}
dunque $\mathbf{O}_n(\mathbf{R})$ è chiuso in $\mathbf{R}^{n^2}$, pertanto è chiuso anche in $\mathbf{GL}_n(\mathbf{R})$, nella topologia relativa. \\ \\
Analogamente $\det : \mathbf{M}_n(\mathbf{R}) \longrightarrow \mathbf{R}$ è continua (anch'essa è un polinomio) e
\begin{gather}
    \mathrm{det}^{-1}(\{ 1 \}) = \mathbf{SL}_n(\mathbf{R}) \notag
\end{gather}
pertanto anche $\mathbf{SL}_n(\mathbf{R})$ è chiuso in $\mathbf{R}^{n^2}$ e quindi anche nella topologia relativa su $\mathbf{GL}_n(\mathbf{R})$. \\ \\
Infine si ha che $\mathbf{SO}_n(\mathbf{R}) = \mathbf{SL}_n(\mathbf{R}) \cap \mathbf{O}_n(\mathbf{R})$, dunque intersezione di due chiusi e quindi chiuso. \qed \\ \\ \\
\textbf{Proposizione:} $\mathbf{O}_n(\mathbf{R})$ e $\mathbf{SO}_n(\mathbf{R})$ sono compatti. \\

\textbf{Dimostrazione:} sia $A \in \mathbf{O}_n(\mathbf{R})$, si ha che $AA^t = \mathbf{I}_n$, pertanto
\begin{gather}
    AA^t = \mathbf{I}_n \implies \sum_h a_{ih}a_{hj} = \delta_{ij} \notag
\end{gather}
ovvero per $i = j$ si ha
\begin{gather}
    \sum_h a^2_{ih} = 1, \quad i = 1, \dots, n \notag
\end{gather}
quindi si ottiene
\begin{gather}
    \sum_{h, i} a^2_{ih} = n \notag
\end{gather}
cioè $\mathbf{O}_n(\mathbf{R}) \subset \overline{D_{\sqrt{n}}(0)}$, allora $\mathbf{O}_n(\mathbf{R})$ è un sottoinsieme chiuso e limitato di $\mathbf{R}^{n^2}$ e per il \textit{corollario 1.5.3} è compatto. Inoltre $\mathbf{SO}_n(\mathbf{R})$ è un sottoinsieme chiuso di un insieme compatto e pertanto compatto. \qed

\subsubsection{Esercizi}
\textbf{1.} Dimostrare se $\mathbf{O}_n(\mathbf{R}) \leq \mathbf{O}_{n + 1}(\mathbf{R})$, cioè se esiste $S \leq \mathbf{O}_{n + 1}(\mathbf{R})$ tale che $\mathbf{O}_n(\mathbf{R}) \cong S$. \\

\textbf{Dimostrazione:} consideriamo $f : \mathbf{O}_n(\mathbf{R}) \longrightarrow S$ definita ponendo
\begin{gather}
    f(A) = \begin{pmatrix}
        A & 0 \\
        0 & 1
    \end{pmatrix} \text{ per ogni } A \in \mathbf{O}_n(\mathbf{R}) \notag
\end{gather}
con $S := f(\mathbf{O}_n(\mathbf{R}))$. Notiamo immediatamente che $f$ è biettiva e continua e anche $f^{-1}$ lo è. Inoltre $S \leq \mathbf{O}_{n + 1}(\mathbf{R})$, infatti
\begin{gather}
    \begin{pmatrix}
        A & 0 \\
        0 & 1
    \end{pmatrix}
    \begin{pmatrix}
        A^t & 0 \\
        0 & 1
    \end{pmatrix} = 
    \begin{pmatrix}
        AA^t & 0 \\
        0 & 1
    \end{pmatrix} = 
    \begin{pmatrix}
        \mathbf{I}_n & 0 \\
        0 & 1
    \end{pmatrix} = \mathbf{I}_{n + 1} \notag
\end{gather}
dunque se $M \in S$ si ha $M \in \mathbf{O}_{n + 1}(\mathbf{R})$, inoltre $\mathbf{I}_{n + 1} \in S$, $S$ è chiuso rispetto al prodotto di matrici e per ogni $M \in S$ esiste $M^{-1} \in S$ tale che $MM^{-1} = \mathbf{I}_{n + 1}$, infatti $A^{-1} \in \mathbf{O}_n(\mathbf{R})$, e
\begin{gather}
    \begin{pmatrix}
        A & 0 \\
        0 & 1
    \end{pmatrix}
    \begin{pmatrix}
        A^{-1} & 0 \\
        0 & 1
    \end{pmatrix} = 
    \mathbf{I}_{n + 1} \notag
\end{gather}
per quanto visto prima. \\ \\ \\
\textbf{2.} Sia
\begin{gather}
    \textit{Aff}_n(\mathbf{R}) := \left\{ \begin{pmatrix}
        A & 0 \\
        v & 1
    \end{pmatrix} : A \in \mathbf{O}_n(\mathbf{R}), v \in \mathbf{R} \right\} \notag
\end{gather}
detto \textit{gruppo affine} di rango $n$. Dimostrare se $\textit{Aff}_n(\mathbf{R}) \leq \mathbf{GL}_{n + 1}(\mathbf{R})$. \\

\textbf{Dimostrazione:} come nell'\textit{esercizio 1}, dimostriamo che esiste un sottogruppo $S$ di $\mathbf{GL}_{n + 1}(\mathbf{R})$ tale che $\textit{Aff}_n(\mathbf{R}) \cong S$. Sia $\mathbf{1}_{n + 1} : \textit{Aff}_n(\mathbf{R}) \longrightarrow S$ la funzione identità. Ovviamente è un omeomorfismo (ponendo $S = \mathbf{1}_{n + 1}(\textit{Aff}_n(\mathbf{R})$), mostriamo ora che $S$ è un sottogruppo di $\mathbf{GL}_{n + 1}(\mathbf{R})$. Si ha che
\begin{gather}
    \det \begin{pmatrix}
        A & 0 \\
        v & 1
    \end{pmatrix} = \det A \neq 0 \text{ per ogni } A \in \mathbf{O}_n(\mathbf{R}) \notag
\end{gather}
dunque $S \subseteq \mathbf{GL}_{n + 1}(\mathbf{R})$. Inoltre $\mathbf{I}_{n + 1} \in S$, perché $\mathbf{I}_n \in \mathbf{O}_n(\mathbf{R})$, e per ogni elemento di $S$ esiste l'inverso in $S$ stesso, infatti
\begin{gather}
    \begin{pmatrix}
        A & 0 \\
        v & 1
    \end{pmatrix}
    \begin{pmatrix}
        B & 0 \\
        w & 1
    \end{pmatrix} = 
    \begin{pmatrix}
        AB & 0 \\
        vB + w & 1
    \end{pmatrix} = \mathbf{I}_{n + 1} \notag
\end{gather}
se e solo se
\begin{gather}
    B = A^t = A^{-1} \in \mathbf{O}_n(\mathbf{R}), \quad w = -vB \in \mathbf{R}^n \notag
\end{gather}
pertanto $S \leq \mathbf{GL}_{n + 1}(\mathbf{R})$. \\ \\ \\
\textbf{3.} Siano $G$ e $H$ gruppi, $f : G \longrightarrow H$ un omomorfismo e $S \leq G$. Allora $f(S) \leq H$. \\

\textbf{Dimostrazione:} sia $e$ l'elemento neutro di $G$, si ha che $e \in S$, inoltre
\begin{gather}
    f(a) = f(a \cdot e) = f(a) \cdot f(e) = f(e) \cdot f(a) = f(e \cdot a) = f(a) \notag
\end{gather}
per ogni $a \in G$, ovvero $f(e)$ è l'elemento neutro di $H$ e $f(e) \in f(S)$. Inoltre, sia $a^{-1} \in G$ tale che $aa^{-1} = e$, si ha che
\begin{gather}
    f(e) = f(a \cdot a^{-1}) = f(a) \cdot f(a^{-1}) = f(a^{-1}) \cdot f(a) = f(a^{-1} \cdot a) = f(e) \notag
\end{gather}
ovvero $f(a^{-1}) = f(a)^{-1}$, quindi $f(S)$ contiene l'inverso di ogni suo elemento. Infine siano $a, b \in S$, si ha che
\begin{gather}
    f(a) \cdot f(b) = f(a \cdot b) = f(c) \notag
\end{gather}
con $c$ prodotto di $a$ e $b$, $c \in S$ per ipotesi. Pertanto concludiamo che $f(S)$ è un sottogruppo di $H$. \\ \\ \\
\textbf{4.} Dimostrare se $\textit{Aff}_n(\mathbf{R})$ è un gruppo di matrici abeliano. \\

\textbf{Dimostrazione:} si ha che
\begin{gather}
    \begin{pmatrix}
        A & 0 \\
        v & 1
    \end{pmatrix}
    \begin{pmatrix}
        B & 0 \\
        w & 1
    \end{pmatrix} = 
    \begin{pmatrix}
        AB & 0 \\
        vB + w & 1
    \end{pmatrix} \notag
\end{gather}
mentre
\begin{gather}
    \begin{pmatrix}
        B & 0 \\
        w & 1
    \end{pmatrix}
    \begin{pmatrix}
        A & 0 \\
        v & 1
    \end{pmatrix} = 
    \begin{pmatrix}
        BA & 0 \\
        wA + v & 1
    \end{pmatrix} \notag
\end{gather}
in generale non c'è uguaglianza tra questi due elementi. \\ \\ \\
\textbf{5.} Dimostrare se $\mathbf{GL}_n(\mathbf{R})$ è compatto. \\

\textbf{Soluzione:} ricordiamo che
\begin{gather}
    \mathbf{GL}_n(\mathbf{R}) = \{ A \in \mathbf{M}_n(\mathbf{R}) : \det A \neq 0 \} \notag
\end{gather}
Sia $k \in \mathbf{N}$, consideriamo la seguente matrice
\begin{gather}
    M(k) = \begin{pmatrix}
        k & 0 & 0 & \dots & 0 \\
        0 & k & 0 & \dots & 0 \\
        . & & & & . \\
        . & & & & . \\
        0 & 0 & 0 & \dots & k
    \end{pmatrix} \notag
\end{gather}
per ogni $\rho > 0$ esiste $\Tilde{k} \in \mathbf{N}$ tale che
\begin{gather}
    \Tilde{k} > \rho \implies M(\Tilde{k}) \not\subset D_\rho(0) \notag
\end{gather}
ovvero $\mathbf{GL}_n(\mathbf{R})$ non è limitato, pertanto per il \textit{teorema di Heine-Borel} non può essere compatto. \\ \\ \\
\textbf{6.} Dimostrare se $\textbf{SL}_n(\mathbf{R})$ è compatto in $\mathbf{M}_n(\mathbf{R}$). \\

\textbf{Dimostrazione:} sappiamo che $\mathbf{SL}_n(\mathbf{R})$ è un gruppo di matrici e quindi chiuso in $\mathbf{GL}_n(\mathbf{R})$ e pertanto chiuso anche in $\mathbf{M}_n(\mathbf{R})$. Tuttavia non è compatto, infatti sia $k \in \mathbf{N}$, consideriamo la matrice
\begin{gather} A(k) =
    \begin{pmatrix}
        1 & k & 0 & \dots & 0 \\
        0 & 1 & 0 & \dots & 0 \\
        . & & & & . \\
        . & & & & . \\
        0 & 0 & 0 & \dots & 1
    \end{pmatrix} \in \mathbf{SL}_n(\mathbf{R}) \notag
\end{gather}
per ogni $k \in \mathbf{N}$. Inoltre, analogamente all'esercizio precedente, si ha che per ogni $\rho > 0$ esiste $\Tilde{k} \in \mathbf{N}$ tale che
\begin{gather}
    \sqrt{\sum_{i, j = 1}^n a^2_{ij}(\Tilde{k})} = \sqrt{n + \Tilde{k}^2} > \rho \notag
\end{gather}
infatti basta scegliere
\begin{gather}
    \Tilde{k} > \sqrt{\rho^2 - n} \notag
\end{gather}
e si ha
\begin{gather}
    A(\Tilde{k}) \not\subset D_\rho(0) \notag
\end{gather}
pertanto $\mathbf{SL}_n(\mathbf{R})$ non è limitato e dunque è compatto per il \textit{teorema di Heine-Borel}. \\ \\ \\ \\
\textbf{7.} Sia $G := \left\{ A \in \mathbf{M}_n(\mathbf{R}) : \det A \in \{ 2^k : k \in \mathbf{Z} \} \right\}$, dimostrare che $G$ è un gruppo di matrici. \\

\textbf{Dimostrazione:} poniamo $2^{\mathbf{Z}} := \{ 2^k : k \in \mathbf{Z} \}$, ovviamente $\mathbf{I}_n \in G$ poiché $\det \mathbf{I}_n = 1 = 2^0 \in 2^{\mathbf{Z}}$. Inoltre se $A \in G$ si ha che $A^{-1} \in G$, infatti
\begin{gather}
    \det A^{-1} = \frac{1}{\det A} = \frac{1}{2^{\overline{k}}} \text{ per qualche } \overline{k} \in \mathbf{Z} \notag
\end{gather}
pertanto
\begin{gather}
    \det A^{-1} = 2^{-\overline{k}} \text{ e } -\overline{k} \in \mathbf{Z} \notag
\end{gather}
infatti $Z$ è un gruppo rispetto all'addizione, dunque ogni elemento ha il suo inverso. Siano ora $A, B \in G$, si ha che $A \cdot B \in G$, infatti
\begin{gather}
    \det (A \cdot B) = \det A \cdot \det B = 2^p \cdot 2^q = 2^{p + q} \notag
\end{gather}
con $p, q \in \mathbf{Z}$, pertanto anche $p + q \in \mathbf{Z}$ perché $\mathbf{Z}$ è un gruppo rispetto all'addizione. \\ \\
Dunque $G$ è un sottogruppo di $\mathbf{GL}_n(\mathbf{R})$, verifichiamo ora che sia chiuso in $\mathbf{GL}_n(\mathbf{R})$. \\
Si ha che $2^{\mathbf{Z}}$ è non è chiuso in $\mathbf{R}$, infatti $0$ è un punto di accumulazione per $2^{\mathbf{Z}}$ ma $0 \not\in 2^{\mathbf{Z}}$, tuttavia è chiuso in $\mathbf{R} \setminus \{ 0 \}$ che è sufficiente poiché siamo interessati alle matrici invertibili. Si ha che la funzione
\begin{gather}
    \det|_{\mathbf{GL}_n(\mathbf{R})} : \mathbf{GL}_n(\mathbf{R}) \longrightarrow \mathbf{R} \setminus \{ 0 \} \notag
\end{gather}
è continua, pertanto
\begin{gather}
    \mathrm{det}^{-1}(2^{\mathbf{Z}}) = G \notag
\end{gather}
è chiuso in $\mathbf{GL}_n(\mathbf{R})$, infatti $\mathrm{det}^{-1}(\mathbf{R} \setminus \{ 0 \}) = \mathbf{GL}_n(\mathbf{R})$. \\ \\ \\
\textbf{8.} Dimostrare se $Isom(\mathbf{R}^n)$ è un gruppo di matrici. \\

\textbf{Dimostrazione:} sappiamo che
\begin{gather}
    Isom(\mathbf{R}^n) = \left\{ \begin{pmatrix}
        A & 0 \\
        v & 1
    \end{pmatrix} : A \in \mathbf{O}_n(\mathbf{R}), v \in \mathbf{R}^n \right\} \notag
\end{gather}
$\mathbf{I}_{n + 1} \in Isom(\mathbf{R}^n)$ perché $\mathbf{I}_n \in \mathbf{O}_n(\mathbf{R})$. Inoltre
\begin{gather}
    \begin{pmatrix}
        A & 0 \\
        v & 1
    \end{pmatrix}
    \begin{pmatrix}
        B & 0 \\
        u & 1
    \end{pmatrix} =
    \begin{pmatrix}
        AB & 0 \\
        vB + u & 1
    \end{pmatrix} =
    \mathbf{I}_{n + 1} \notag
\end{gather}
se e solo se $B = A^{-1} \text{ e } u = -vB$ e poiché $A^{-1} \in \mathbf{O}_n(\mathbf{R})$ e $-vB \in \mathbf{R}^n$ si ha che $Isom(\mathbf{R}^n)$ contiene l'inverso di ogni suo elemento. Notiamo anche che $Isom(\mathbf{R}^n)$ è chiuso rispetto al prodotto di matrici. Dunque è un gruppo. \\ \\
Dimostriamo ora che è chiuso in $\mathbf{GL}_{n + 1}(\mathbf{R})$. Consideriamo la funzione \\
$f : Isom(\mathbf{R}^n) \longrightarrow \mathbf{M}_{n + 1}(\mathbf{R})$ definita ponendo
\begin{gather}
    f \begin{pmatrix}
        A & 0 \\
        v & 1
    \end{pmatrix} =
    \begin{pmatrix}
        A & 0 \\
        v & 1
    \end{pmatrix}
    \begin{pmatrix}
        A^t & 0 \\
        -vA^t & 1
    \end{pmatrix} =
    \begin{pmatrix}
        AA^t & 0 \\
        0 & 1
    \end{pmatrix} \notag
\end{gather}
ovviamente $f$ è continua, è un polinomio, inoltre
\begin{gather}
    f^{-1}(\{ \mathbf{I}_{n + 1} \} = Isom(\mathbf{R}^n) \notag
\end{gather}
ed essendo $\{ \mathbf{I}_{n + 1} \}$ chiuso in $\mathbf{M}_{n + 1}(\mathbf{R})$, poiché è uno spazio $T_1$, si ha che anche $Isom(\mathbf{R}^n)$ è chiuso in $\mathbf{M}_{n + 1}(\mathbf{R})$, inoltre, poiché $\mathbf{GL}_{n + 1} \cap Isom(\mathbf{R}^n) = Isom(\mathbf{R}^n)$ è chiuso nella topologia relativa e quindi è un gruppo di matrici.



\newpage
\subsection{Algebre di Lie}
\textbf{Definizione 2.4.1:} sia $X \subseteq \mathbf{R}^n$ e sia $p \in X$. Lo \textit{spazio tangente} ad $X$ in $p$ è
\begin{gather}
    T_p(X) := \{ \gamma'(0) : \gamma : (-\varepsilon, \varepsilon) \longrightarrow X, \gamma \text{ differenziabile}, \gamma(0) = p \} \notag
\end{gather}
Non sempre è possibile determinare lo spazio tangente ad un punto di uno spazio $X$ arbitrario. Tale possibilità è determinata dalla geometria locale dello spazio attorno al punto in esame. \\ \\ \\
\textbf{Definizione 2.4.2:} sia $G$ un gruppo di matrici, l'\textit{algebra di Lie} di $G$ è
\begin{gather}
    \mathfrak{g}(G) := T_{\mathbf{I}}(G) \notag
\end{gather}
ovvero è lo spazio tangente a $G$ nell'identità $\mathbf{I}$. \\ \\
Esibire una funzione $\gamma : (-\varepsilon, \varepsilon) \longrightarrow G$ differenziabile significa esibire $n^2$ funzioni scalari differenziabili
\begin{gather}
    a_{ij} : (-\varepsilon, \varepsilon) \longrightarrow \mathbb{K} \notag
\end{gather}
tali che
\begin{gather}
    \begin{pmatrix}
        a_{11}(t) & \dots & a_{1n}(t) \\
        . & & . \\
        . & & . \\
        a_{n1}(t) & \dots & a_{nn}(t) 
    \end{pmatrix} \in G \text{ per ogni } t \in (-\varepsilon, \varepsilon) \notag
\end{gather} \\ \\ \\
\textbf{Proposizione 2.4.1:} siano $\beta, \gamma : (-\varepsilon, \varepsilon) \longrightarrow \mathbf{M}_n(\mathbf{R})$ differenziabili, definiamo
\begin{gather}
    \beta\gamma : (-\varepsilon, \varepsilon) \longrightarrow \mathbf{M}_n(\mathbf{R}) \notag
\end{gather}
ponendo
\begin{gather}
    (\beta\gamma)(t) = \beta(t) \cdot \gamma(t) \notag
\end{gather}
per ogni $t \in (-\varepsilon, \varepsilon)$. Allora $\beta\gamma$ è differenziabile e
\begin{gather}
    (\beta\gamma)'(t) = \beta'(t) \cdot \gamma'(t) \notag
\end{gather} \\

\textbf{Dimostrazione:} si ha che $\beta\gamma$ è differenziabile poiché ogni sua componente è differenziabile dato che è prodotto e somma di funzioni scalari differenziabili, inoltre
\begin{gather}
    \left( \beta(t) \cdot \gamma(t) \right)_{ij} = \sum_h \beta_{ih}(t)\gamma_{hj}(t) \notag 
\end{gather}
quindi
\begin{gather}
    \left( \beta(t) \cdot \gamma(t) \right)'_{ij} = \sum_h \left( \beta'_{ih}(t)\gamma_{hj}(t) + \beta_{ih}(t)\gamma'_{hj}(t) \right) \notag
\end{gather}
per la teoria sulle funzioni scalari ad una (o più) variabile. Inoltre, poiché ciò vale per ogni $i, j = 1, \dots, n$ si ha
\begin{gather}
    (\beta\gamma)'(t) = \beta'(t)\gamma(t) + \beta(t)\gamma'(t) \notag
\end{gather} \qed \\ \\
\textbf{Proposizione 2.4.2:} sia $G$ un gruppo di matrici, allora $\mathfrak{g}(G)$ è uno spazio vettoriale reale. \\

\textbf{Dimostrazione:} sia $A \in \mathfrak{g}(G)$ e sia $\lambda \in \mathbf{R}$, per ipotesi esiste $\gamma$ differenziabile tale che $\gamma'(0) = A$, consideriamo il cammino
\begin{gather}
    \Tilde{\gamma} : \left( -\frac{\varepsilon}{|\lambda|}, \frac{\varepsilon}{|\lambda|} \right) \longrightarrow G, \quad \Tilde{\gamma}(t) = \gamma(\lambda t) \notag
\end{gather}
per ogni $t \in (\varepsilon, \varepsilon)$. $\Tilde{\gamma}$ è differenziabile, $\Tilde{\gamma}(0) = \mathbf{I}_n$ e inoltre
\begin{gather}
    \Tilde{\gamma}'(0) = \lambda \cdot \gamma'(0) = \lambda \cdot A \notag
\end{gather}
pertanto $\lambda \cdot A \in \mathfrak{g}(G)$. \\
Siano $A, B \in \mathfrak{g}(G)$, per ipotesi esistono $\alpha, \beta$ differenziabili tali che $\alpha'(0) = A$ e $\beta'(0) = B$, allora $\alpha\beta$ è differenziabile e
\begin{gather}
    (\alpha\beta)'(0) = \alpha'(0)\beta(0) + \alpha(0)\beta'(0) = A + B \in \mathfrak{g}(G) \notag
\end{gather}
Infine si ha che la matrice nulla appartiene a $\mathfrak{g}(G)$ poiché la funzione costante $\delta(t) = \mathbf{I}_n$ per ogni $t \in (-\varepsilon, \varepsilon)$ è differenziabile e $\delta'(0) = 0 \in \mathfrak{g}(G)$. \qed \\ \\ \\
\textbf{Definizione 2.4.3:} sia $G$ un gruppo di matrici. La dimensione di $G$ è $\dim G = \dim \mathfrak{g}(G)$ come spazio vettoriale reale. \\ \\
Poniamo $\mathfrak{gl}_n(\mathbf{R}) = \mathfrak{g}(\mathbf{GL}_n(\mathbf{R}))$, $\mathfrak{sl}_n(\mathbf{R}) = \mathfrak{g}(\mathbf{SL}_n(\mathbf{R}))$, $\mathfrak{so}_n(\mathbf{R}) = \mathfrak{g}(\mathbf{SO}_n(\mathbf{R}))$. \\ \\ \\
\textbf{Proposizione 2.4.3:} sia $n \in \mathbf{N}$, si ha che $\mathfrak{gl}_n(\mathbf{R}) = \mathbf{M}_n(\mathbf{R})$. \\

\textbf{Dimostrazione:} ovviamente $\mathfrak{gl}_n(\mathbf{R}) \subset \mathbf{M}_n(\mathbf{R})$. Sia $A \in \mathbf{M}_n(\mathbf{R})$, consideriamo il cammino
\begin{gather}
    \gamma : \mathbf{R} \longrightarrow \mathbf{M}_n(\mathbf{R}), \quad \gamma(t) = \mathbf{I}_n + tA \notag
\end{gather}
per ogni $A \in \mathbf{M}_n(\mathbf{R})$. Chiaramente $\gamma$ è differenziabile poiché è un polinomio,  $\gamma(0) = \mathbf{I}_n$ e $\gamma'(0) = A$. \\
Inoltre la funzione $\det$ è continua dunque, per il teorema della permanenza del segno, esiste $\varepsilon > 0$ tale che $\det \gamma(t) > 0$ per ogni $t \in (-\varepsilon, \varepsilon)$, pertanto la tesi segue e $\dim \mathbf{GL}_n(\\mathbf{R}) = n^2$. \qed \\ \\ \\
\textit{Lemma 2.3.1:} sia $\gamma : (-\varepsilon, \varepsilon) \longrightarrow \mathbf{M}_n(\mathbf{R})$ differenziabile tale che $\gamma(0) = \mathbf{I}_n$, si ha che
\begin{gather}
    \frac{\mathrm{d}}{\mathrm{d}t} \det \gamma(t) |_{t = 0} = \mathrm{tr}(\gamma'(0)) \notag
\end{gather} \\

\textit{Dimostrazione:} esplicitando i calcoli so ottiene
\begin{gather}
    \frac{\mathrm{d}}{\mathrm{d}t} \det \gamma(t)|_{t = 0} = \frac{\mathrm{d}}{\mathrm{d}t} \left( \sum_{i = 1}^n (-1)^{i - 1} \gamma_{1i}(t) \det(\gamma[1, i](t)) \right)|_{t = 0} = \notag \\
    = \sum_{i = 1}^n (-1)^{i - 1}\gamma'_{1i}(0) \det(\gamma[1,i](0)) + \sum_{i = 1}^n (-1)^{i - 1} \gamma_{1i}(0) \frac{\mathrm{d}}{\mathrm{d}t} [\det (\gamma[1, i](t))]|_{t = 0} = \notag \\
    = \gamma'_{11}(0) + \frac{\mathrm{d}}{\mathrm{d}t} \det(\gamma[1, 1](t))|_{t = 0} = \notag \\
    = \gamma'_{11} + \mathrm{tr}(\gamma'[1,1](0)) = \sum_{i = j} \gamma'_{ij}(0) \notag
\end{gather}
dove il penultimo passaggio si ottiene ragionando per induzione. \qed \\ \\ \\
\textbf{Teorema 2.4.2:} sia $n \in \mathbf{N}$, si ha che
\begin{gather}
    \mathfrak{sl}_n(\mathbf{R}) = \{ A \in \mathbf{M}_n(\mathbf{R}) : \mathrm{tr} A = 0 \} \notag
\end{gather}
e in particolare $\dim \mathbf{SL}_n(\mathbf{R}) = n^2 - 1$. \\

\textbf{Dimostrazione:} sia $A \in \mathfrak{sl}_n(\mathbf{R})$, per ipotesi esiste $\gamma$ differenziabile tale che $\gamma(0) = \mathbf{I}_n$ e $\gamma'(0) = A$. Inoltre deve soddisfare $\gamma(t) \in \mathbf{SL}_n(\mathbf{R})$ per ogni $t \in (-\varepsilon, \varepsilon)$, pertanto
\begin{gather}
    \frac{\mathrm{d}}{\mathrm{d}t} \det \gamma(t) = \mathrm{tr} (\gamma'(0)) = 0 \notag
\end{gather}
imponendo $\det \gamma(t) = 1$ per ogni $t \in (-\varepsilon, \varepsilon)$. \\ \\
Viceversa, sia $\alpha : (-\varepsilon, \varepsilon) \longrightarrow \mathbf{SL}_n(\mathbf{R})$ definita ponendo
\begin{gather}
    \alpha(t) = \begin{pmatrix}
        \frac{\gamma_{11}(t)}{\det \gamma(t)} & \dots & \frac{\gamma_{1n}(t)}{\det \gamma(t)} \\
        . & & . \\
        . & & . \\
        \gamma_{n1}(t) & \dots & \gamma_{nn}(t)
    \end{pmatrix} \notag
\end{gather}
si ha che $\alpha$ è differenziabile (per teoremi generali sulle funzioni scalari), \\$\alpha(t) \in \mathbf{SL}_n(\mathbf{R})$ per ogni $t \in (-\varepsilon, \varepsilon)$, inoltre $\alpha(0) = \mathbf{I}_n$ e
\begin{gather}
    \frac{\mathrm{d}}{\mathrm{d}t} \left( \frac{\gamma_{1i}}{\det \gamma(t)} \right)|_{t = 0} = \left( \frac{\gamma'_{1i}(0) \det \gamma(0) - \gamma_{1i}(0) \frac{\mathrm{d}}{\mathrm{d}t} (\det \gamma(t))|_{t = 0}}{(\det \gamma(0))^2} \right) = \notag \\
    =  \gamma'_{1i}(0) - \gamma_{11}(0)\mathrm{tr}(A) = \gamma'_{1i}(0) \notag
\end{gather}
per il \textit{lemma 2.4.1} e poiché, per ipotesi, $\mathrm{tr}(\gamma'(0)) = \mathrm{tr}(A) = 0$. Pertanto $\alpha'(0) = \gamma'(0)$ e quindi si ottiene la tesi. \\ \\
In particolare si ha che $\dim \mathbf{SO}_n(\mathbf{R}) = n^2 - 1$ poiché la condizione per appartenere all'algebra di Lie è $\mathrm{tr}(A) = 0$, pertanto è una condizione che influisce solamente sugli $n$ elementi della diagonale, in particolare $n - 1$ sono indipendenti mentre l'$n$-esimo è dipendente poiché deve soddisfare la condizione di traccia nulla, ovvero deve soddisfare la seguente equazione:
\begin{gather}
    \sum_{1 \leq i = j \leq n - 1} a_{ij} = - a_{nn} \notag
\end{gather}
dove sono stati scegli come indipendenti i primi $n - 1$ elementi della diagonale senza perdita di generalità. \qed \\ \\ \\
\textbf{Proposizione 2.4.4:} sia $A \in \mathbf{M}_n(\mathbf{R})$, si ha che le seguenti sono equivalenti: \\
\textbf{1.} $A + A^t = 0$. \\
\textbf{2.} $\langle x, yA \rangle = -\langle xA, y \rangle$ per ogni $x, y \in \mathbf{R}^n$. \\
\textbf{3.} $\langle x, xA \rangle = 0$ per ogni $x \in \mathbf{R}^n$. \\

\textbf{Dimostrazione:} $(1 \implies 2)$ Notiamo che $(1)$ implica
\begin{gather}
    a_{ij} = 0 \text{ per ogni } i = j, 1 \leq i, j \leq n \notag \\
    a_{ij} = - a_{ji} \text{ per ogni } 1 \leq i, j \leq n \notag
\end{gather}
si ha che $x = \sum_i x_i \mathbf{e}_i$ e $y = \sum_j y_j \mathbf{e}_j$, allora
\begin{gather}
    \langle x, yA \rangle = \sum_i \sum_j x_iy_j \langle \mathbf{e}_i, \mathbf{e}_jA \rangle = \sum_i \sum_j x_iy_j \langle \mathbf{e}_i, A_j \rangle \notag = \sum_i \sum_j x_iy_jA_{ji} = \notag \\
    = - \sum_i \sum_j x_iy_j A_{ij} = -\sum_i \sum_j x_iy_j \langle \mathbf{e}_j, \mathbf{e}_iA \rangle = \notag \\
    = -\langle \sum_j y_j\mathbf{e}_j, \sum_i x_i\mathbf{e}_i A \rangle = -\langle y, xA \rangle \notag
\end{gather}
per la bilinearità del prodotto scalare euclideo. \\ \\
$(2 \implies 3)$ Immediato, basta porre $y = x$ e si ha
\begin{gather}
    \langle x, xA \rangle = - \langle xA, x \rangle \implies 2\langle x, xA \rangle = 0 \notag
\end{gather}
per la simmetria del prodotto scalare euclideo. \\ \\
$(3 \implies 2)$ Notiamo che
\begin{align}
    0 = \langle x + y, xA + yA \rangle &= \langle x, xA \rangle + \langle y, yA \rangle + \langle x, yA \rangle + \langle y, xA \rangle \notag \\
    &= \langle x, yA \rangle + \langle y, xA \rangle \notag
\end{align}
pertanto la tesi segue. \\ \\
$(2 \implies 1)$ Come visto prima si ha che
\begin{gather}
    \sum_i \sum_j x_iy_j A_{ji} + \sum_i \sum_j x_iy_j A_{ij} = 0 \iff A_{ij} = -A_{ji} \notag
\end{gather}
per ogni $1 \leq i, j \leq n$. \\ \\
Le matrici che soddisfano le proprietà sopra elencate sono dette \textit{matrici antisimmetriche}. \\
Notiamo che la $(3)$ è equivalente a dire $x \perp xA$ per ogni $A$ antisimmetrica. \qed \\ \\ \\
\textbf{Teorema 2.4.3:} si ha che
\begin{gather}
    \mathfrak{so}_n(\mathbf{R}) = \{ A \in \mathbf{M}_n(\mathbf{R}) : A + A^t = 0 \} \notag
\end{gather}
ovvero l'algebra di Lie del gruppo speciale ortogonale è l'insieme della matrici antisimmetriche a coefficienti reali. \\

\textbf{Dimostrazione:} $(\rightarrow)$ sia $X \in \mathfrak{so}_n(\mathbf{R})$, per ipotesi esiste $\gamma$ differenziabile tale che $\gamma(0) = \mathbf{I}_n$, $\gamma'(0) = X$ e $\gamma(t) \in \mathbf{SO}_n(\mathbf{R})$ per ogni $t \in (-\varepsilon, \varepsilon)$. \\
Prima di tutto, poiché $\mathbf{SO}_n(\mathbf{R}) \leq \mathbf{O}_n(\mathbf{R})$, deve valere che $\gamma(t) \cdot \gamma(t)^t = \mathbf{I}_n$ per ogni $t \in (-\varepsilon, \varepsilon)$, dunque
\begin{gather}
    \frac{\mathrm{d}}{\mathrm{d}t} (\gamma(t) \cdot \gamma(t)^t) = 0 \notag
\end{gather}
ovvero
\begin{gather}
    \gamma'(t) \cdot \gamma(t)^t + \gamma(t) \cdot \gamma'(t)^t = 0 \notag
\end{gather}
in particolare per $t = 0$, pertanto si ottiene
\begin{gather}
    \gamma'(0) \cdot \gamma (0)^t + \gamma(0) \cdot \gamma'(0)^t = X + X^t = 0 \notag
\end{gather}
quindi $\mathfrak{so}_n(\mathbf{R}) \subset \{ A \in \mathbf{M}_n(\mathbf{R}) : A + A^t = 0 \}$. \\ \\
$(\leftarrow)$ Viceversa sia $X \in \mathbf{M}_n(\mathbf{R})$ tale che $X + X^t = 0$, dimostriamo che $X$ appartiene all'algebra di Lie del gruppo speciale ortogonale, ovvero dimostriamo che esiste un cammino differenziabile $\gamma$ tale che $\gamma(0) = \mathbf{I}_n$, $\gamma'(0) = X$ e $\gamma(t) \in \mathbf{SO}_n(\mathbf{R})$ per ogni $t \in (-\varepsilon, \varepsilon)$. \\
Consideriamo il cammino $\gamma : (-\varepsilon, \varepsilon) \longrightarrow \mathbf{SO}_n(\mathbf{R})$, definito ponendo
\begin{gather}
    \gamma(t) = e^{tX} = \sum_{n = 0}^\infty \frac{t^n}{n!}X^n \text{ per ogni } t \in (-\varepsilon, \varepsilon) \notag
\end{gather}
ovviamente $\gamma$ è differenziabile, $\gamma(0) = \mathbf{I}_n$ e $\gamma'(0) = X$. Inoltre per ipotesi si ha
\begin{gather}
    e^{tX} \cdot (e^{tX})^t = e^{tX} \cdot e^{tX^t} = e^{t(X + X^t)} = \mathbf{I}_n \notag
\end{gather}
per ogni $t \in (-\varepsilon, \varepsilon)$, infatti
\begin{gather}
    X + X^t = 0 \notag
\end{gather}
ciò comporta anche che $X$ è invertibile e che $\det X = 1$, pertanto la tesi segue. \qed \\ \\ \\
\textbf{Corollario 2.4.1:} si ha che
\begin{gather}
    \dim \mathbf{SO}_n(\mathbf{R}) = \dim \mathbf{O}_n(\mathbf{R}) = \begin{pmatrix}
        n \\
        2
    \end{pmatrix} \notag
\end{gather}

\textbf{Dimostrazione:} si ha che $1 + 2 + \dots (n - 1) = \begin{pmatrix}
    n \\
    2
\end{pmatrix}$, ovvero il numero delle componenti indipendenti. \qed \\ \\ \\
Questa trattazione può essere eseguita analogamente per matrici a coefficienti complessi, le uniche accortezze sono: \\
\textbf{1.} $\langle x, y \rangle = \sum_i x_i \overline{y}_i$. \\
\textbf{2.} al posto di $A^t$ si ha $A^* = (\overline{a}_{ji})_{1 \leq i, j \leq n}$, trasposta coniugata. \\
\textbf{3.} si ha che $\mathfrak{g}(G)$ è uno spazio vettoriale reale anche se $G$ è un gruppo di matrici a coefficienti complessi. \\
\textbf{4.} si scrive $\mathbf{U}_n$ e $\mathbf{SU}_n$ al posto di $\mathbf{O}_n(\mathbf{C})$ e $\mathbf{SO}_n(\mathbf{C})$, e sono detti \textit{gruppo unitario} e \textit{gruppo speciale unitario}.

\subsubsection{Esercizi}
\textbf{1.} Determinare: 
\begin{equation}
    \mathfrak{g}(Isom\, \mathbf{R}^n) := \{ \gamma ' (0) : \gamma : (- \varepsilon, \varepsilon) \longrightarrow Isom\, \mathbf{R}^n, \, \gamma \text{ diff.}, \, \gamma(0) = \mathbf{I}_n \} \notag
\end{equation} 
dove $\mathbf{I}_n$ è la matrice identità di $\mathbf{M}_n(\mathbf{R})$. \\

\textbf{Soluzione:} sia $\gamma : (- \varepsilon, \varepsilon) \longrightarrow Isom\, \mathbf{R}^n$ definita ponendo:
\begin{equation} \gamma (t) =
    \begin{pmatrix}
        e^{tX} & 0 \\
        v(t) & 1
    \end{pmatrix} \notag
\end{equation}
dove $X \in \mathbb{M}_n$ e $v(t) \in C^1((-\varepsilon, \varepsilon);\mathbf{R}^n)$ tale che $v(0) = \vec{0}$. Per come è stato definito $Isom\, \mathbf{R}^n$ deve valere che $e^{tX} \in \mathbf{O}_n(\mathbf{R})$ per ogni $t \in (-\varepsilon, \varepsilon)$, cioè è sufficiente che:
\begin{gather}
    e^{tX} \cdot (e^{tX})^t = \mathbf{I}_n \iff X + X^t = 0 \notag
\end{gather}
Inoltre per continuità della funzione $\det$ e per il teorema della permanenza del segno si ha che:
\begin{gather}
    \frac{d}{dt} \det (e^{tX})|_{t=0} = 0 = \mathrm{tr}\, \gamma'(0) = \mathrm{tr}\, X \notag
\end{gather}
condizione già implicita in $X + X^t = 0$. Pertanto:
\begin{equation}
    \mathfrak{g}(Isom\, \mathbf{R}^n) = \left\{
    \begin{pmatrix}
        X & 0 \\
        v'(0) & 0
    \end{pmatrix} : X + X^t = 0, \ v \in C^1((-\varepsilon, \varepsilon); \mathbf{R}^n), \ v(0) = 0\right\} \notag
\end{equation}
\\ \\ \\
\textbf{2.} sia $G := \{ A \in \mathbb{M}_n : \det A \in 2^{\mathbf{Z}}\}$, determinare $\mathfrak{g}(G)$. \\

\textbf{Soluzione:} si ha che intorno a $\mathbf{I}_n$ il gruppo di matrici $G$ è localmente identico a $\mathbf{SL}_n(\mathbb{R)}$ e poiché lo spazio tangente è univocamente determinato da caratteristiche locali segue che $\mathfrak{g}(G) = \mathfrak{sl}_n(\mathbf{R})$.

\newpage
\section{Tensori}
\subsection{Quoziente di spazi vettoriali}
\textbf{Definizione 3.1.1:} sia $V$ un $\mathbb{K}$-spazio vettoriale e sia $W \leq V$ (sottospazio vettoriale, lo \textit{spazio vettoriale quoziente} è
\begin{gather}
    V/W := \{ W + v : v \in V \} \notag
\end{gather}
ovvero, fissato $v \in V$, l'insieme dei sottospazi affini dati da
\begin{gather}
    W + v := \{ w + v : w \in W \} \notag
\end{gather}
Poniamo
\begin{gather}
    (W + v_1) + (W + v_2) = W + (v_1 + v_2), \quad \lambda(W + v) = W + \lambda v \notag
\end{gather}
per ogni $v, v_1, v_2 \in V$ e per ogni $\lambda \in \mathbb{K}$. \\ \\
Osserviamo che $W + v = W + u$ se e solo se $v - u \in W$, infatti $W + w = W$ se e solo se $w \in W$ per le proprietà degli spazi vettoriali. \\ \\ \\
\textbf{Proposizione 3.1.1:} siano $V$ uno spazio vettoriale di dimensione finita, $W \leq V$, si ha che $\dim V/W = \dim V - \dim W$. \\

\textbf{Dimostrazione:} sia $\dim V = n$ e $\dim W = k$, $k \leq n$, dimostriamo che $n - k$ vettori linearmente indipendenti generano $V/W$. Consideriamo una base di $W$
\begin{gather}
    \{ a_1, \dots, a_k \} \notag
\end{gather}
completiamola a base di $V$
\begin{gather}
    \{ a_1, \dots, a_k, a_{k + 1}, \dots, a_n \} \notag
\end{gather}
sia $v \in V$, si ha
\begin{gather}
    W + v = W + (\sum_{i = 1}^n \lambda_i a_i) = \sum_{i = 1}^n (W + \lambda_i a_i) = \sum_{i = k + 1}^n (W + \lambda_ia_i) \notag 
\end{gather}
poiché $\sum_{i = 1}^k \lambda_i a_i \in W$, dunque
\begin{gather}
    W + (\sum_{i = 1}^k \lambda_i a_i) = W \notag
\end{gather}
pertanto $\{ W + a_{k + 1}, \dots, W + a_n \}$ è un insieme di vettori linearmente indipendenti che genera $V/W$, ovvero $\dim V/W = n - k$. \qed

\newpage
\subsection{Prodotti e somme dirette}
\textbf{Definizione 3.2.1:} sia $\{ V_i : i \in I \}$ una famiglia di spazi vettoriali reali di dimensione finita. La \textit{somma diretta} di $\{ V_i : i \in I \}$ è lo spazio vettoriale
\begin{gather}
    \bigoplus_{i \in I} V_i \notag
\end{gather}
che ha come elementi famigli di vettori $(x_i)_{i \in I}$ dove $x_i \in V_i$ per ogni $i \in I$ e $x_i = 0$ per quasi tutti gli $i \in I$, ovvero per tutti gli indici $i$ tranne un numero finito. \\ \\
Definiamo le operazioni
\begin{gather}
    (x_i)_{i \in I} + (y_i)_{i \in I} = (x_i + y_i)_{i \in I} \notag
\end{gather}
e
\begin{gather}
    \lambda \cdot (x_i)_{i \in I} = (\lambda x_i)_{i \in I} \notag
\end{gather}
per ogni $(x_i)_{i \in I}, (y_i)_{i \in I} \in \bigoplus_{i \in I} V_i$ e ogni $\lambda \in \mathbf{R}$. \\ \\
La condizione che le componenti non nulle siano finite serve a mantenere la possibilità di trattare combinazioni lineari come somme finite. \\ \\ \\
\textbf{Definizione 3.2.2:} sia $\{ V_i : i \in I \}$ una famiglia di spazi vettoriali reali di dimensione finita. Il \textit{prodotto diretto} della famiglia $\{ V_i : i \in I \}$ è lo spazio vettoriale
\begin{gather}
    \prod_{i \in I} V_i \notag
\end{gather}
che ha come elementi le famiglie $(x_i)_{i \in I}$, dove $x_i \in V_i$ per ogni $i \in I$, con le stesse operazioni definite per la somma diretta di spazi vettoriali reali. \\ \\
Osserviamo che
\begin{gather}
    \bigoplus_{i \in I} V_i \leq \prod_{i \in I} V_i \notag
\end{gather}
in particolare
\begin{gather}
    \bigoplus_{i \in I} V_i = \prod_{i \in I} V_i \iff |I| < +\infty \notag
\end{gather}
uguaglianza per isomorfismo. \\ \\
\textbf{Proposizione 3.2.1:} siano $V$ e $W$ spazi vettoriali di dimensione finita, si ha che $\dim V \oplus W = \dim V + \dim W$. \\
Inoltre se $\{ e_i : i = 1, \dots, n \}$ è una base di $V$ e $\{ f_i : i = 1, \dots, m \}$ è una base di $W$ si ha che l'insieme
\begin{gather}
    \{ (e_i, 0)_{i = 1, \dots, n} \} \cup \{ (0, f_i)_{i = 1, \dots, m} \} \notag
\end{gather}
è una base per $V \oplus W$. \\

\textbf{Dimostrazione:} immediata. \qed \\ \\ \\
\textbf{Definizione 3.2.3:} sia $S$ un insieme, lo \textit{spazio vettoriale di base} $S$ è
\begin{gather}
    \mathbf{R}^S := \bigoplus_{s \in S} \mathbf{R} \notag
\end{gather}
per identificare un elemento di $\mathbf{R}^S$ si usa la notazione di somma formale, ovvero
\begin{gather}
    (x_s)_{s \in S} = \sum_{s \in S} a_s \cdot s \notag
\end{gather}

\newpage
\subsection{Prodotto tensoriale}
\textbf{Definizione 3.3.1:} siano $V$ e $W$ spazi vettoriali reali, il prodotto tensoriale di $V$ e $W$ è lo spazio vettoriale quoziente
\begin{gather}
    V \otimes W = \frac{\mathbf{R}^{V \times W}}{D} \notag
\end{gather}
dove $V \times W$ è inteso come insieme delle coppie $(v, w)$ tali che $v \in V$ e $w \in W$, mentre $D$ è l'insieme generato da
\begin{gather}
    (x_1 + x_2, y) - (x_1, y) - (x_2, y) \notag \\
    (x, y_1 + y_2) - (x, y_1) - (x, y_2) \notag \\
    a \cdot (x, y) - (ax, y) \notag \\
    a \cdot (x, y) - (x, ay) \notag
\end{gather}
per ogni $x, x_1, x_2 \in V$, ogni $y, y_1, y_2 \in W$ e ogni $a \in \mathbf{R}$. \\ \\
Se $V$ è uno spazio vettoriale e $W \leq V$ si ha che $\pi : V \longrightarrow \frac{V}{W}$ definita ponendo
\begin{gather}
    \pi(v) = W + v \notag
\end{gather}
è detta \textit{proiezione canonica}. \\ \\
Siano $x \in V$ e $y \in W$ si scrive $x \otimes y = \pi ((x, y))$, dunque $\pi : \mathbf{R}^{V \times W} \longrightarrow V \otimes W$. \\
Per quanto visto in precedenza $(x, y)$ è un elemento di
\begin{gather}
    \mathbf{R}^{V \times W} = \bigoplus_{(v, w) \in V \times W} \mathbf{R} \notag
\end{gather}
ovvero $(x, y) = (x, y, 0)$, cioè ha un numero finito di componenti non nulle. Pertanto $\pi ((x, y))$ trasla il sottospazio vettoriale $D$ di $(x, y)$. \\ \\
Quindi si ottiene
\begin{gather}
    x \otimes y = \pi((x, y)) = D + (x, y) \notag
\end{gather}
allora
\begin{gather}
    D + ((x_1 + x_2, y) - (x_1, y) - (x_2, y)) = D \notag
\end{gather}
cioè
\begin{gather}
    D + (x_1 + x_2, y) = D + ((x_1, y) + (x_2, y)) \notag
\end{gather}
proprio perché per definizione $(x_1 + x_2, y) - (x_1, y) - (x_2, y) \in D$, per ogni $x_1, x_2 \in V$ e ogni $y \in W$, ovvero
\begin{gather}
    (x_1 + x_2) \otimes y = x_1 \otimes y + x_2 \otimes y \notag
\end{gather}
Con un ragionamento analogo si ottiene
\begin{gather}
    x \otimes (y_1 + y_2) = x \otimes y_1 + x \otimes y_2 \notag \\
    a \cdot (x \otimes y) = ax \otimes y \notag \\
    a \cdot (x \otimes y) = x \otimes ay \notag
\end{gather}
cioè $\otimes : \mathbf{R}^{V \times W} \longrightarrow V \otimes W$ è una forma bilineare. \\ \\ \\
\textbf{Teorema (proprietà universale del prodotto tensoriale):} siano $V, W$ spazi vettoriali di dimensione finita, si ha che esiste \underline{unico} $T$ spazio vettoriale e \underline{unica} $g : V \times W \longrightarrow T$ bilineare tali che per ogni spazio vettoriale $U$ e ogni funzione $f : V \times W \longrightarrow U$ bilineare esiste \underline{unica} $h : T \longrightarrow U$ tale che $h \cdot g = f$. \\
\[
\begin{tikzcd}
&V \times W \arrow[r,"f"] \arrow[d,"g", swap] & U \\
&T \arrow[ur, "h", swap] 
\end{tikzcd}
\]


\textbf{Dimostrazione:} poniamo
\begin{gather}
    T := V \otimes W, \quad g(x, y) := x \otimes y \notag
\end{gather}
$g$ è bilineare per quanto visto in precedenza. Sia $U$ uno spazio vettoriale e $f : V \times W \longrightarrow U$ una funzione bilineare. Definiamo la funzione
\begin{gather}
    \Tilde{f} : \mathbf{R}^{V \times W} \longrightarrow U \notag
\end{gather}
definita ponendo
\begin{gather}
    \Tilde{f}(\sum_{(x, y) \in V \times W} a_{x, y} \cdot (x, y)) = \sum_{(x, y) \in V \times W} a_{x, y} \cdot f(x, y) \notag
\end{gather}
si ha che, per come è stata costruita, $\Tilde{f}$ è lineare. Inoltre $\Tilde{f}(D) = 0$ proprio per la linearità di $\Tilde{f}$. \\ \\
Poniamo $h : V \otimes W \longrightarrow U$ definita ponendo
\begin{gather}
    h(D + \sum_{(x, y) \in V \times W} a_{x, y} \cdot (x, y)) = \Tilde{f}(\sum_{(x, y) \in V \times W} a_{x, y} \cdot (x, y)) \notag
\end{gather}
per ogni $\sum_{(x, y) \in V \times W} a_{x, y} \cdot (x, y) \in \mathbf{R}^{V \times W}$. Notiamo che $D + \alpha = D + \beta$ se e solo se $\alpha - \beta \in D$, pertanto si ottiene
\begin{gather}
    h(D + (\alpha - \beta)) = \Tilde{f}(\alpha - \beta) = \Tilde{f}(\alpha) - \Tilde{f}(\beta) = 0 \notag
\end{gather}
quindi $\Tilde{f}(\alpha) = \Tilde{f}(\beta)$ e $h$ è lineare. Inoltre, se $(x, y) \in V \times W$ si ha
\begin{gather}
    (h \cdot g)(x, y) = h(g(x, y)) = h(x \otimes y) = \Tilde{f}((x, y)) = f(x, y) \notag
\end{gather}
pertanto $h \cdot g = f$. \\ \\
Mostriamo ora che $h$ è unica. Supponiamo che esista $h' : V \otimes W \longrightarrow U$ tale che
\begin{gather}
    (h' \cdot g)(x, y) = h'(g(x, y)) = h'(x \otimes y) = \Tilde{f}((x, y)) = f(x, y) \notag
\end{gather}
per ogni $(x, y) \in V \times W$, segue che $h = h'$ poiché
\begin{gather}
    \{ x \otimes y : x \in V, y \in W \} \notag
\end{gather}
genera $V \otimes W$, ovvero sono funzioni lineari che combaciano per ogni elemento di un'insieme generatore, pertanto devono essere uguali. \\ \\
Mostriamo ora che $T$ e $g$ sono uniche. Supponiamo che esistano $T'$ spazio vettoriale e $g' : V \times W \longrightarrow T'$ che soddisfano le condizioni del teorema. Allora esiste $h : T' \longrightarrow T$ lineare tale che $h \cdot g' = g$ \\
\[
\begin{tikzcd}
&V \times W \arrow[r,"g"] \arrow[d,"g'", swap] & T \\
&T' \arrow[ur, "h", swap] 
\end{tikzcd}
\]
Analogamente esiste $h' : T \longrightarrow T'$ lineare tale che $h' \cdot g = g'$ \\
\[
\begin{tikzcd}
&V \times W \arrow[r,"g'"] \arrow[d,"g", swap] & T' \\
&T \arrow[ur, "h'", swap] 
\end{tikzcd}
\]
inoltre esiste $l : T' \longrightarrow T'$ lineare tale che $l \cdot g' = g'$ \\
\[
\begin{tikzcd}
&V \times W \arrow[r,"g'"] \arrow[d,"g'", swap] & T' \\
&T' \arrow[ur, "l", swap] 
\end{tikzcd}
\]
dunque si ha che $g' = h' \cdot (h \cdot g') = (h' \cdot h) \cdot g'$, ovvero $l = \mathbf{1}_{T'} = h' \cdot h$, pertanto $h$ è invertibile, dunque iniettiva, e suriettiva, ovvero $h$ è un isomorfismo, si ottiene
\begin{gather}
    g' \cong g \quad T' \cong T \notag
\end{gather}
ovvero sono unici. \qed \\ \\
Consideriamo ora il seguente diagramma \\
\[
\begin{tikzcd}
&V \times V^* \arrow[r,"f"] \arrow[d,"g", swap] & \mathbb{K} \\
&V \otimes V^* \arrow[ur, "h", swap] 
\end{tikzcd}
\]
dove $f$ è una funzione bilineare e $V$ è uno spazio vettoriale sul campo $\mathbb{K}$. Per la \textit{proprietà universale del prodotto tensoriale} si ha che
\begin{gather}
    \mathrm{Bil}(V \times V^*, \mathbb{K}) \cong \mathrm{Lin}(V \otimes V^*, \mathbb{K}) \notag
\end{gather}
ovvero
\begin{gather}
    \mathrm{Bil}(V \times V^*, \mathbb{K}) \cong (V \otimes V^*)^* \notag
\end{gather}
per definizione di spazio vettoriale duale. Ciò può essere generalizzato al caso di funzioni \textit{multilineari}, ovvero consideriamo $r$ copie di $V$ e $s$ copie di $V^*$, si ha che
\begin{gather}
    \mathrm{Mult}(V \times \dots \times V \times V^* \times \dots \times V^*, \mathbb{K}) \cong (V \otimes \dots \otimes V \otimes V^* \otimes \dots \otimes V^*)^* \notag
\end{gather}
pertanto, poiché $V \cong (V^*)^*$ si ottiene
\begin{gather}
    \mathrm{Mult}(V \times \dots \times V \times V^* \times \dots \times V^*, \mathbb{K}) \cong (V \otimes \dots \otimes V \otimes V^* \otimes \dots \otimes V^*) \notag
\end{gather}
ovvero un tensore è una funzione multilineare $t : V^{\otimes r} \otimes (V^*)^{\otimes s} \longrightarrow \mathbb{K}$. \\ \\ \\
\textbf{Corollario 3.3.1:} dalla \textit{proprietà universale del prodotto tensoriale} si ottiene: \\
\textbf{1.} $V \otimes W \cong W \otimes V$. \\
\textbf{2.} $(V \otimes W) \otimes U \cong V \otimes (W \otimes U)$. \\
\textbf{3.} $V \otimes (W \oplus U) \cong (V \oplus W) \otimes U$. \\
\textbf{4.} $V \otimes \mathbf{R} \cong V$. \\ \\ \\
\textbf{Proposizione 3.3.1:} siano $V, W$ spazi vettoriali di dimensione finita, siano $\{ e_i : i = 1, \dots, n \}$ una base di $V$ e $\{ f_j : j = 1, \dots, m \}$ una base di $W$, si ha che
\begin{gather}
    \{ e_i \otimes f_j : i = 1, \dots, n; j = 1, \dots, m \} \notag
\end{gather}
è una base di $V \otimes W$. In particolare $\dim (V \otimes W) = \dim V \cdot \dim W$. \\

\textbf{Dimostrazione:} siano $v \in W$ e $w \in W$, per ipotesi si ha
\begin{gather}
    v = \sum_{i = 1}^n \lambda_i e_i, \quad w = \sum_{j = 1}^m \mu_j f_j \notag
\end{gather}
pertanto
\begin{gather}
    v \otimes w = \sum_i \sum_j \lambda_i \mu_j  (e_i \otimes f_j) \notag
\end{gather}
ovvero
\begin{gather}
    \{ e_i \otimes f_j : i = 1, \dots, n; j = 1, \dots, m \} \notag
\end{gather}
genera $V \otimes W$. Inoltre
\begin{gather}
    V = \bigoplus_{i = 1}^n \mathbf{R} \notag \\
    W = \bigoplus_{j = 1}^m \mathbf{R} \notag
\end{gather}
pertanto, per i punti $(3)$ e $(4)$ del \textit{corollario 3.3.1}
\begin{gather}
    V \otimes W \cong \bigoplus_i \bigoplus_j (\mathbf{R} \otimes \mathbf{R}) \cong \bigoplus_i \bigoplus_j \mathbf{R} \notag
\end{gather}
ovvero $\dim (V \otimes W) = n \cdot m$. \qed \\ \\
Ragionando per induzione si può generalizzare la \textit{proposizione 3.3.1} al prodotto tensoriale di una famiglia finita di spazi vettoriali di dimensione finita.

\subsubsection{Esercizi}
\textbf{1.} Sia $V$ uno spazio vettoriale, sia $W \leq V$, determinare $\ker (\pi)$. \\

\textbf{Soluzione:} si ha che $\ker (\pi) = W$, infatti
\begin{gather}
    \pi(w) = W \notag
\end{gather}
per ogni $w \in W$, per definizione di spazio vettoriale e proiezione canonica, inoltre $W$ è l'elemento neutro di $V/W$, infatti
\begin{gather}
    (W + v) + W = (W + v) + (W + 0) = W + (v + 0) = W + v \notag
\end{gather}
pertanto si ottiene la tesi. \\ \\ \\
\textbf{2.} Siano $V, W$ spazi vettoriali, siano $\{ e_i : i \in I \}$ una base di $V$ e $\{ f_j : j \in J \}$ una base di $W$, determinare una base di $\mathrm{Hom} (V, W)$. \\

\textbf{Soluzione:} si ha che
\begin{gather}
    \mathrm{Hom}(V, W) = \{ L : V \longrightarrow W : L \text{ omomorfismo} \} \notag
\end{gather}
siano $v \in V$ e $L \in \mathrm{Hom}(V, W)$, si ha che $v = \sum_i \lambda_i e_i$, inoltre
\begin{gather}
    L(v) = L( \sum_{i \in I} \lambda_i e_i) = \sum_{i \in I} \lambda_i L(e_i) = \sum_{i \in I} \lambda_i ( \sum_{j \in J} \mu^i_j f_j) = \sum_{i \in I} \sum_{j \in J} \lambda_i \mu^i_j f_j  \notag
\end{gather}
pertanto per ogni $k \in I$ la funzione $\phi_{ij} : V \longrightarrow W$ definita ponendo
\begin{gather}
    \phi_{ij}(e_k) = \delta_{ik}f_j \notag
\end{gather}
è un omomorfismo e l'insieme
\begin{gather}
    \{ \phi_{ij} : i \in I, j \in J \} \notag
\end{gather}
genera $\mathrm{Hom}(V, W)$, pertanto $\dim \mathrm{Hom}(V, W) = \dim V \cdot \dim W$. \\ \\ \\
\textbf{3.} Sia $V$ uno spazio vettoriale, dimostrare che $V \otimes \mathbf{R} \cong V$. \\

\textbf{Dimostrazione:} sfruttiamo la \textit{proprietà universale del prodotto tensoriale}, consideriamo il seguente diagramma \\
\[
\begin{tikzcd}
&V \times \mathbf{R} \arrow[r,"f"] \arrow[d,"g", swap] & V \\
&V \otimes \mathbf{R} \arrow[ur, "\Phi", swap] 
\end{tikzcd}
\]
dunque esibiamo una funzione $f : V \times \mathbf{R} \longrightarrow V$ bilineare e dimostriamo che $\Phi$ è un isomorfismo. Come scelta naturale per $f$ si ha
\begin{gather}
    f(v, \lambda) := \lambda \cdot v \notag
\end{gather}
Per ogni $(v, \lambda) \in V \times \mathbf{R}$. Per ipotesi sappiamo che esiste unico $\Phi : V \otimes \mathbf{R} \longrightarrow V$ lineare tale che
\begin{gather}
    \Phi(v \otimes \lambda) = \lambda \cdot v = f(v, \lambda) \notag
\end{gather}
per ogni $(v, \lambda) \in V \times \mathbf{R}$. Esibiamo l'inversa di $\Phi$, consideriamo la funzione $\Psi : V \longrightarrow V \otimes \mathbf{R}$ definita ponendo
\begin{gather}
    \Psi (v) = v \otimes 1 \notag
\end{gather}
allora si ottiene
\begin{gather}
    \Phi \cdot \Psi (v) = \Phi(v \otimes 1) = 1 \cdot v = v \notag
\end{gather}
e
\begin{gather}
    \Psi \cdot \Phi (v \otimes \lambda) = \Psi (\lambda v) = v \otimes \lambda \notag
\end{gather}
dunque $\Phi$ è un isomorfismo e $V \times \mathbf{R} \cong V$. \\ \\
Notiamo che, nonostante sembrerebbe che è possibile utilizzare la proprietà universale del prodotto tensoriale per dimostrare l'esistenza di un isomorfismo tra una qualsiasi coppia $V, W$ di spazi vettoriali, esibire una funzione bilineare $f : V \longrightarrow W$ non è scontato e in alcuni casi potrebbe anche non esistere, perciò è l'esistenza di $f$ che impone un limite alle applicazioni del teorema. \\ \\ \\
\textbf{4.} Dimostrare che $\mathbf{M}_{1, n} \otimes \mathbf{M}_{m, 1} \cong \mathbf{M}_{n, m}$. \\

\textbf{Soluzione:} sfruttiamo la \textit{proprietà universale del prodotto tensoriale}. Sia $\{ e_j : j = 1, \dots, n \}$ una base di $\mathbf{M}_{1, n}$ e sia $\{ f_i : i = 1, \dots, m \}$ una base di $\mathbf{M}_{m, 1}$, inoltre sia $\{ m_{ij} : 1 \leq i \leq m, 1 \leq j \leq n \}$ una base di $\mathbf{M}_{m, n}$. Consideriamo la funzione $f : \mathbf{M}_{1, n} \times \mathbf{M}_{m, 1} \longrightarrow \mathbf{M}_{n, m}$ definita ponendo
\begin{gather}
    f(e_j, f_i) = m_{ij} \notag
\end{gather}
per ogni $1 \leq i \leq m$ e ogni $1 \leq j \leq n$, ovviamente $f$ è bilineare. Per ipotesi esiste unica $\Phi : \mathbf{M}_{1, n} \otimes \mathbf{M}_{m, 1} \longrightarrow \mathbf{M}_{n, m}$ lineare tale che
\begin{gather}
    \Phi(e_j \otimes f_i) = m_{ij} = f(e_j, f_i) \notag
\end{gather}
per ogni $1 \leq i \leq m$ e ogni $1 \leq j \leq n$. Esibiamo l'inversa di $\Phi$, consideriamo la funzione $\Psi : \mathbf{M}_{m, n} \longrightarrow \mathbf{M}_{1, n} \otimes \mathbf{M}_{m, 1}$ definita ponendo
\begin{gather}
    \Psi (m_{ij}) = e_j \otimes f_i \notag
\end{gather}
allora si ha
\begin{gather}
    \Phi \cdot \Psi (m_{ij}) = \Phi(e_j \otimes f_i) = m_{ij} \notag
\end{gather}
e analogamente
\begin{gather}
    \Psi \cdot \Phi(e_j \otimes f_i) = \Psi(m_{ij}) = e_j \otimes f_i \notag
\end{gather}
ovvero $\Phi$ è un isomorfismo e $\mathbf{M}_{1, n} \otimes \mathbf{M}_{m, 1} \cong \mathbf{M}_{n, m}$.

\newpage
\textbf{5.} Siano $V, W$ spazi vettoriali di dimensione finita, dimostrare che $\mathrm{Hom}(V, W) \cong V^* \otimes W$. \\

\textbf{Soluzione:} sfruttiamo la \textit{proprietà universale del prodotto tensoriale}. Sia $\{ e_i : i = 1, \dots, n \}$ una base di $V$ e sia $\{ f_j : j = 1, \dots, m \}$ una base di $W$, inoltre sia $\{ h_{ij} : 1 \leq i \leq n, 1 \leq j \leq m \}$ una base di $\mathrm{Hom}(V, W)$. Consideriamo la funzione $f : V^* \times W \longrightarrow \mathrm{Hom}(V, W)$ definita ponendo
\begin{gather}
    f(e^i, f_j) = h_{ij} \notag
\end{gather}
per ogni $1 \leq i \leq n, 1 \leq j \leq m$. Il resto della dimostrazione è del tutto analogo alla dimostrazione dell'\textit{esercizio 5}.

\newpage
\subsection{Tensori}
\textbf{Definizione 3.4.1:} sia $V$ uno spazio vettoriale, il \textit{duale} di $V$ è lo spazio vettoriale
\begin{gather}
    V^* := \{ f : V \longrightarrow \mathbf{R} : f \text{ lineare} \} \notag
\end{gather}
Inoltre sia $\{ e_i : i = 1, \dots, n \}$ una base di $V$ e sia $e^i \in V^*$ definito ponendo
\begin{gather}
    e^i(e_j) = \delta_{ij} \notag
\end{gather}
si ha che $\{ e^i : i = 1, \dots, n \}$ è una base di $V^*$ ed è detta \textit{base duale}. \\ \\ \\
\textbf{Definizione 3.4.2:} siano $V$ uno spazio vettoriale e $p, q$ interi positivi, un \textit{tensore} di tipo $(p, q)$ è un elemento
\begin{gather}
    t \in T^p_q(V) := V \otimes \dots \otimes V \otimes V^* \otimes \dots \otimes V^* \notag
\end{gather}
$p$ copie di $V$ e $q$ copie di $V^*$. \\ \\
Un tensore di tipo $(p, 0)$ si dice \textit{controvariante} di grado $p$, un tensore di tipo $(0, q)$ si dice \textit{covariante} di grado $q$. Un tensore di tipo $(1, 0)$ si dice \textit{vettore}, un tensore di tipo $(0, 1)$ si dice \textit{covettore} (o \textit{vettore covariante}). \\ \\ 
Sia $\{ e_i : i = 1, \dots, n \}$ una base di $V$ e sia $\{ e^i : i = 1, \dots, n \}$ la base duale, per la \textit{proposizione 3.3.1} si ha che esistono unici $a^{i_1, \dots, i_p}_{j_1, \dots, j_q}$ tali che
\begin{gather}
    t = \sum_{1 \leq i_1, \dots, i_p \leq n} \sum_{1 \leq j_1, \dots, j_q \leq n} a^{i_1, \dots, i_p}_{j_1, \dots, j_q} e_{i_1} \otimes \dots \otimes e_{i_p} \otimes e^{j_1} \otimes \dots \otimes e^{j_q} \notag
\end{gather} \\ \\
\textbf{Proposizione 3.4.1:} sia $V$ uno spazio vettoriale di dimensione finita, sia $\{ f_i : i = 1, \dots, n \}$ una base di $V$ e sia $\{ f^i : i = 1, \dots, n \}$ la base dello spazio duale $V^*$. Consideriamo la base $\{ e_i : i = 1, \dots, n \}$ dello spazio vettoriale $V$ e la rispettiva base duale $\{ e^i : i = 1, \dots, n \}$ tali che
\begin{gather}
    t = \sum_{1 \leq i_1, \dots, i_p \leq n} \sum_{1 \leq j_1, \dots, j_q \leq n} a^{i_1, \dots, i_p}_{j_1, \dots, j_q} e_{i_1} \otimes \dots \otimes e_{i_p} \otimes e^{j_1} \otimes \dots \otimes e^{j_q} \notag
\end{gather}
si ha che
\begin{gather}
    t = \sum_{1 \leq i_1, \dots, i_p \leq n} \sum_{1 \leq j_1, \dots, j_q \leq n} b^{i_1, \dots, i_p}_{j_1, \dots, j_q} f_{i_1} \otimes \dots \otimes f_{i_p} \otimes f^{j_1} \otimes \dots \otimes f^{j_q} \notag
\end{gather}
dove
\begin{gather}
    b^{k_1, \dots, k_p}_{l_1, \dots, l_q} = \sum_{1 \leq i_1, \dots, i_p \leq n} \sum_{1 \leq j_1, \dots, j_q \leq n} a^{i_1, \dots, i_p}_{j_1, \dots, j_q} \cdot \alpha^{l_1}_{i_1} \cdot \dots \cdot \alpha^{l_p}_{i_p} \cdot \beta^{j_1}_{k_1} \cdot \dots \cdot \beta^{j_q}_{k_q} \notag
\end{gather}
e
\begin{gather}
    e_i = \sum_{j} \alpha_i^j f_j, \quad e^i = \sum_j \beta^i_j f^j \notag
\end{gather}
per ogni $i = 1, \dots, n$. \\

\textbf{Dimostrazione:} per la bilinearità del prodotto tensoriale si ha che
\begin{gather*}
    t = \sum_{1 \leq i_1, \dots, i_p \leq n} \sum_{1 \leq j_1, \dots, j_q \leq n} a^{i_1, \dots, i_p}_{j_1, \dots, j_q} e_{i_1} \otimes \dots \otimes e_{i_p} \otimes e^{j_1} \otimes \dots \otimes e^{j_q} = \notag \\
    = \sum_{1 \leq i_1, \dots, i_p \leq n} \sum_{1 \leq j_1, \dots, j_q \leq n} a^{i_1, \dots, i_p}_{j_1, \dots, j_q} (\sum_{l_1 = 1}^n \alpha_{i_1}^{l_1}f_{l_1}) \otimes \dots \otimes (\sum_{l_p = 1}^n \alpha_{i_p}^{l_p}f_{l_p}) \otimes \notag \\
    \otimes (\sum_{k_1 = 1}^n \beta^{j_1}_{k_1} f^{k_1}) \otimes \dots \otimes (\sum_{k_q = 1}^n \beta^{j_q}_{k_q} f^{k_q}) \notag = \\
    = \sum_{1 \leq l_1, \dots, l_p \leq n} \sum_{1 \leq k_1, \dots, k_q \leq n} (\sum_{1 \leq i_1, \dots, i_p \leq n} \sum_{1 \leq j_1, \dots, j_q \leq n} a^{i_1, \dots, i_p}_{j_1, \dots, j_q} \cdot \notag \\ 
    \cdot \alpha^{l_1}_{i_1} \dots \alpha^{l_p}_{i_p} \cdot \beta^{j_1}_{k_1} \dots \beta^{j_q}_{k_q}) \cdot f_{l_1} \otimes \dots \otimes f_{l_p} \otimes f^{k_1} \otimes \dots \otimes f^{k_q} \notag
\end{gather*}
dunque la tesi segue. \qed \\ \\ \\
\textbf{Proposizione 3.4.2:} siano $V, W$ spazi vettoriali di dimensione finita, poniamo
\begin{gather}
    \mathrm{End}(V) := \mathrm{Hom}(V, V) \notag
\end{gather}
si ha che $\mathrm{End}(V) \cong V^* \otimes V$. \\

\textbf{Dimostrazione:} sfruttiamo la \textit{proprietà universale del prodotto tensoriale}, consideriamo il seguente diagramma \\
\[
\begin{tikzcd}
&V^* \times V \arrow[r,"f"] \arrow[d,"g", swap] & \mathrm{End}(V) \\
&V^* \otimes V \arrow[ur, "\Phi", swap] 
\end{tikzcd}
\]
dunque esibiamo $f : V^* \times V \longrightarrow \mathrm{End}(V)$ bilineare e mostriamo che
$\Phi$ è un isomorfismo. Sia $\{ e_i : i = 1, \dots, n \}$ una base di $V$ e sia $\{ e^i : i = 1, \dots, n \}$ la base duale, inoltre sia $\{ \alpha_{ij} : 1 \leq i, j \leq n \}$ una base di $\mathrm{End}(V)$, poniamo
\begin{gather}
    f(e_i, e^j) := \alpha_{ij} \notag
\end{gather}
per ogni $1 \leq i, j \leq n$. Per ipotesi esiste unica $\Phi : V^* \otimes V \longrightarrow \mathrm{End}(V)$ lineare tale che
\begin{gather}
    \Phi(e_i \otimes e^j) = \alpha_{ij} \notag
\end{gather}
Esibiamo l'inversa di $\Phi$, consideriamo la funzione $\Psi : \mathrm{End}(V) \longrightarrow V^* \otimes V$ definita ponendo
\begin{gather}
    \Psi (\alpha_{ij}) = e_i \otimes e^j \notag
\end{gather}
ovviamente si ha che
\begin{gather}
    \Phi \cdot \Psi (\alpha_{ij}) = \Phi (e_i \otimes e^j) = \alpha_{ij} \notag
\end{gather}
e analogamente
\begin{gather}
    \Psi \cdot \Phi (e_i \otimes e^j) = \Psi(\alpha_{ij}) = e_i \otimes e^j \notag
\end{gather}
dunque $\Phi$ è un isomorfismo e $V^* \otimes V \cong \mathrm{End}(V)$. \qed

\newpage
\section{Varietà}
\subsection{Varietà topologiche e differenziabili}
\textbf{Definizione 4.1.1:} sia $X$ uno spazio topologico, $X$ è detto \textit{varietà topologica} se: \\
\textbf{1.} $X$ è $T_2$. \\
\textbf{2.} $X$ soddisfa il secondo assioma di numerabilità. \\
\textbf{3.} per ogni $p \in X$ esiste un intorno aperto $U$ di $p$ tale che $U$ è omeomorfo a $\{ \mathbf{x} \in \mathbf{R}^n : \| \mathbf{x} \| < 1 \}$ per qualche $n$ intero positivo, in tal caso si pone $\dim (p) = n$. \\ \\ \\
\textbf{Proposizione 4.1.1:} sia $X$ una varietà topologica connessa, si ha che esiste $n$ intero positivo tale che $\dim (p) = n$ per ogni $p \in X$. \\

\textbf{Dimostrazione:} la funzione $\dim : X \longrightarrow \mathbf{R}$ è continua (per come è definita), pertanto $\dim (X) \subset \mathbf{R}$ è connesso, segue che $\dim(X) = \{ n \}$. \\ \\
Se $X$ è una varietà topologica connessa e $\dim (p) = n$ per qualche $p \in X$, si ha che $\dim (p) = n$ per ogni $p \in X$, dunque poniamo $\dim X = n$. \qed \\ \\
Se $\dim X = 1$ $X$ è detto \textit{curva}, se $\dim X = 2$ $X$ è detto superficie. \\ \\ \\
\textbf{Definizione 4.1.2:} sia $X$ uno spazio topologico, $X$ è detto \textit{varietà differenziabile} se esiste $\{ U_\mu, f_\mu \}_{\mu \in M}$ tale che: \\
\textbf{1.} $\{ U_\mu : \mu \in M \}$ è un ricoprimento aperto di $X$ per ogni $\mu \in M$. \\
\textbf{2.} per ogni $\mu \in M$ esiste $V_\mu$ aperto di $\mathbf{R}^n$ tale che $f_\mu : U_\mu \longrightarrow V_\mu$ è un omeomorfismo. \\
\textbf{3.} per ogni $\alpha, \beta \in M$ tali che $U_\alpha \cap U_\beta \neq \emptyset$ si ha che $f_\alpha \cdot f^{-1}_\beta$ e $f_\beta \cdot f^{-1}_\alpha$ sono $C^{\infty}$ (differenziabili). \\
\textbf{4.} $X$ è $T_2$. \\
\textbf{5.} $X$ soddisfa il secondo assioma di numerabilità. \\ \\
$(U_\mu, f_\mu)$ è detta \textit{carta locale}, $\{ U_\mu, f_\mu \}_{\mu \in M}$ è detto \textit{atlante}. \\ \\
Sia $f : \mathbf{R}^n \longrightarrow \mathbf{R}$ differenziabile, $\Gamma_f$ è una varietà differenziabile con atlante $(X, p)$ dove $X = \Gamma_f$ e $p$ è definita ponendo
\begin{gather}
    p(x_1, \dots, x_{n + 1}) = (x_1, \dots, x_n) \notag
\end{gather} 
per ogni $(x_1, \dots, x_{n + 1}) \in X$. \\ \\
Sia $\mathbf{S}^2 = \{ \mathbf{x} \in \mathbf{R}^3 : \| \mathbf{x} \| < 1 \}$, si ha che $\mathbf{S}^2$ è una varietà differenziabile con atlante $\{ U_{ijk}, f_{ijk} \}$, dove
\begin{gather}
    U_{100} = \{ \mathbf{x} \in \mathbf{S}^2 : x_1 > 0 \} \dots U_{00-1} = \{ \mathbf{x} \in \mathbf{S}^2 : x_3 < 0 \} \notag
\end{gather}
e
\begin{gather}
    f_{100}(x_1, x_2, x_3) = (x_2, x_3) \dots f_{00-1}(x_1, x_2, x_3) = (x_1, x_2) \notag
\end{gather}
Notiamo che $\{ U_{ijk} \}$ è un ricoprimento aperto di $\mathbf{S}^2$, inoltre $f_{ijk}$ è differenziabile per ogni indice $ijk$. Inoltre
\begin{gather}
    f_{100} \cdot f^{-1}(x_1, x_2) = f_{100}(x_1, x_2, \sqrt{1 - x^2_1 - x^2_2}) = (x_2, \sqrt{1 - x^2_1 - x^2_2}) \notag
\end{gather}
è differenziabile in $\{ \mathbf{x} \in \mathbf{R}^3 : x^2_1 + x^2_2 < 1 \}$. Ragionando analogamente si dimostra la differenziabilità per tutte le composizioni di funzioni. \\ \\ \\
\textbf{Definizione 4.1.3:} sia $X$ una varietà differenziabile con atlante $\{ U_\mu, f_\mu \}_{\mu \in M}$, si ha che $f : X \longrightarrow \mathbf{R}$ è differenziabile se $f \cdot f^{-1}_\mu$ è differenziabile per ogni $\mu \in M$. \\ \\
Analogamente, se $Y$ è una varietà differenziabile con atlante $\{ V_\alpha, g_\alpha \}_{\alpha \in H}$, si ha che $f : X \longrightarrow Y$ è differenziabile se $g_{\alpha} \cdot f \cdot f^{-1}_\mu$ è differenziabile per ogni coppia $(\mu, \alpha) \in M \times H$. \\ \\
Per $i = 1, \dots, n$ sia $x_i : \mathbf{R}^n \longrightarrow \mathbf{R}$ definita ponendo
\begin{gather}
    x_i(y_1, \dots, y_n) = y_i \notag
\end{gather}
per ogni $(y_1, \dots, y_n) \in \mathbf{R}^n$. Sia $p \in X$ e sia $\mu \in M$ tale che $p \in U_\mu$, si ha che
\begin{gather}
    (x_1(f_\mu(p)), \dots, x_n(f_\mu(p))) \notag
\end{gather}
sono dette \textit{coordinate locali} di $X$ rispetto all'atlante $\{ U_\mu, f_\mu \}_{\mu \in M}$. \\ \\ \\
Questa è l'utilità principale della struttura differenziabile definita in questo paragrafo, infatti se uno spazio $X$ possiede una struttura di varietà differenziabile è possibile applicare gli strumenti e i risultati del calcolo infinitesimale riconducendosi, tramite gli omeomorfismi dell'atlante, allo spazio euclideo. \\ \\ \\
\textbf{Definizione 4.1.4:} siano $X, Y$ varietà differenziabili, sia $f : X \longrightarrow Y$, $f$ è detto \textit{diffeomorfismo} se $f$ e $f^{-1}$ sono omeomorfismi differenziabili.

\subsubsection{Esercizi}
\textbf{1.} Sia $V$ uno spazio vettoriale reale di dimensione finita $n$, consideriamo il sottospazio vettoriale
\begin{gather}
    \mathbf{P}(V) = \{ W \leq V : \dim W = 1 \} \notag
\end{gather}
detto \textit{spazio proiettivo}, dimostrare che $\mathbf{P}(V)$ è una varietà differenziabile. \\

\textbf{Dimostrazione:} esibiamo un atlante differenziabile di $\mathbf{P}(V)$. Consideriamo una base $\{ e_i : i = 1, \dots, n \}$ dello spazio vettoriale $V$, poniamo
\begin{gather}
    U_i = \{ W \in \mathbf{P}(V) : \langle w, e_i \rangle \neq 0 \text{ per ogni } w \in W \setminus \{ 0 \} \} \notag
\end{gather}
si ha che $\{ U_i : i = 1, \dots, n \}$ è un ricoprimento aperto di $\mathbf{P}(V)$. Inoltre per ogni $i = 1, \dots, n$ consideriamo la funzione $f_i : U_i \longrightarrow \mathbf{R}^{n - 1}$ definita ponendo
\begin{gather}
    f_i(\langle \mathbf{x} \rangle) = \left( \frac{x_1}{x_i}, \dots, \frac{x_{i - 1}}{x_i}, \frac{x_{i + 1}}{x_i}, \dots, \frac{x_n}{x_i} \right) \notag
\end{gather}
dove con $\langle \mathbf{x} \rangle$ si identifica lo \textit{span} del vettore $\mathbf{x}$ e $x_i = \langle \mathbf{x}, e_i \rangle$. Dunque $\{ (U_i, f_i) : i = 1, \dots, n \}$ è un atlante differenziabile. Inoltre si verifica facilmente che $f_i \cdot f_j^{-1}$ è differenziabile per ogni $1 \leq i, j \leq n$ tali che $U_i \cap U_j \neq \emptyset$.

\newpage
\subsection{Spazio tangente}
Diamo ora una definizione di \textit{spazio tangente} ad una varietà differenziabile $X$ in un punto $p \in X$. Per fare ciò generalizzeremo il concetto di \textit{derivata direzionale} definito nello spazio euclideo. \\ \\
Sia $X = \mathbf{R}^n$, per qualche $n$ intero positivo, e sia $p \in X$, ogni vettore $\mathbf{v}$ di $\mathbf{R}^n$ induce un'operazione sulle funzioni differenziabili $f$ definite in un intorno aperto di $p$ detta derivata direzionale rispetto a $\mathbf{v}$, data da
\begin{gather}
    D_{\mathbf{v}}(f)(p) = \lim_{t \rightarrow 0} \frac{f(p + t\mathbf{v}) - f(p)}{t} \notag
\end{gather}
nel caso in cui $\mathbf{v} = \mathbf{e}_i$ per qualche $i = 1, \dots, n$, l'operazione viene detta \textit{derivata parziale} rispetto alla $i$-esima componente e si denota con
\begin{gather}
    \frac{\partial f}{\partial x_i}(p) = \lim_{t \rightarrow 0} \frac{f(p + t\mathbf{e}_i) - f(p)}{t} \notag
\end{gather}
Notiamo immediatamente che l'operazione di derivata è lineare e soddisfa la proprietà di Leibniz. \\ \\
Nel caso in cui sia possibile determinare una parametrizzazione $\Phi : \mathbf{R}^n \longrightarrow M$ per una determinata varietà differenziabile $M$ si ha che lo spazio tangente a $M$ nel punto $x \in M$ è lo spazio generato dall'insieme
\begin{gather}
    \left\{ \frac{\partial \Phi}{\partial x_i} : i = 1, \dots, n \right\} \notag
\end{gather}
viene dunque naturale, in assenza di parametrizzazioni, sistemi di coordinate e immersioni, definire lo spazio tangente come segue. \\ \\
\textbf{Definizione 4.2.1:} sia $X$ una varietà differenziabili, sia $p \in X$, lo \textit{spazio tangente} a $p$ in $X$ è
\begin{gather}
    T_pX = \{ L : C^{\infty}(X) \longrightarrow \mathbf{R} : L \text{ lineare }, L(fg) = L(f) \cdot g + L(g) \cdot f \} \notag
\end{gather}
Gli elementi di $T_pX$ sono detti \textit{vettori tangenti} o \textit{derivazioni}. \\ \\
Notiamo che defininendo le seguenti operazioni
\begin{gather}
    (L_1 + L_2)(f) = L_1(f) + L_2(f), \quad \lambda \cdot L(f) = L(\lambda f) \notag
\end{gather}
per ogni $f \in C^{\infty}(X)$, ogni $L, L_1, L_2 \in T_pX$ e ogni $\lambda \in \mathbf{R}$, si ottiene che $T_pX$ è uno spazio vettoriale.
\newpage
\textbf{Teorema 4.2.1:} siano $X$ una varietà differenziabile e $p \in X$, si ha che $\dim T_pX = \dim X$. \\

\textbf{Dimostrazione:} sia $\{ (U_\mu, \varphi_\mu) \}_{\mu \in M}$ un atlante di $X$, sia $\mu \in M$ tale che $p \in U_\mu$ e sia $n$ intero positivo tale che $U_\mu \cong V_\mu$ dove $V_\mu$ aperto di $\mathbf{R}^n$. Consideriamo una base $\{ e_i : i = 1, \dots, n \}$ di $\mathbf{R}^n$ e per ogni $i = 1, \dots, n$ e ogni $f \in C^{\infty}(X; \mathbf{R})$ poniamo
\begin{gather}
    L_i(f) = \frac{\partial}{ \partial x_i}(f \cdot \varphi^{-1}_\mu)|_{\varphi_\mu(p)} \notag
\end{gather}
ovviamente $L_i$ è lineare e soddisfa la proprietà di Leibniz per ogni $i = 1, \dots, n$, pertanto $L_i \in T_p X$ per ogni $i = 1, \dots, n$. Notiamo che $f \cdot \varphi^{-1}_\mu : V_\mu \longrightarrow \mathbf{R}$, dobbiamo dimostrare che $\{ L_i : i = 1, \dots, n \}$ è un insieme linearmente indipendente di generatori di $T_p X$, ovvero che ogni $L \in T_p X$ può essere scritto come combinazione lineare degli $L_i$ e che
\begin{gather}
    \sum_{i = 1}^n \alpha_i L_i = 0 \iff \alpha_i = 0 \text{ per ogni } i = 1, \dots, n \notag
\end{gather}
Osserviamo immediatamente che sono linearmente indipendenti, ovvero
\begin{gather}
    \sum_i \alpha_i L_i(f) = \sum_i \alpha_i \frac{\partial}{\partial x_i}(f \cdot \varphi^{-1}_\mu)|_{\varphi_\mu(p)} = 0 \notag
\end{gather}
per ogni $f \in C^{\infty}(X; \mathbf{R})$ se e solo se $\alpha_i = 0$ per ogni $i = 1, \dots, n$. Sia $L \in T_p X$, poiché $L$ ha come dominio lo spazio delle funzioni di classe $C^{\infty}$ si ha che, per intorni abbastanza piccoli di $\varphi_\mu(p)$, posso scrivere una funzione $f \cdot \varphi^{-1}_\mu$, dove $f \in C^{\infty}(X; \mathbf{R})$, come la sua espansione al primo ordine, ovvero posso applicare $L$ all'espansione al primo ordine di $f$, pertanto
\begin{gather}
    L(f) = L \left( f \cdot \varphi^{-1}_\mu(\varphi_\mu(p)) + \sum_{i = 1}^n \frac{\partial}{\partial x_i}(f \cdot \varphi^{-1}_\mu)|_{\varphi_\mu(p)} (x_i - \varphi_\mu(p)) \right) \notag
\end{gather}
inoltre poiché $L$ annulla le funzioni costanti (segue dalla proprietà di Leibniz) si ottiene
\begin{gather}
    L(f) = \sum_{i = 1}^n L(x_i - \varphi_\mu(p))L_i(f) = \sum_{i = 1}^n \lambda_i L_i(f) \notag
\end{gather}
con $\lambda_i = L(x_i)$, pertanto la tesi segue. \qed \\ \\ \\
\textbf{Definizione 4.2.2:} sia $f \in C^{\infty}(X; Y)$, il \textit{differenziale} di $f$ in $p$ è la funzione lineare
\begin{gather}
    \mathrm{d}f_p : T_pX \longrightarrow T_{f(p)}Y \notag
\end{gather}
definita ponendo
\begin{gather}
    \mathrm{d}f_p(L)[g] = L(g \cdot f) \notag
\end{gather}
per ogni $L \in T_pX$ e ogni $g \in C^{\infty}(Y, X)$.
\end{document}

